\chapter[The (2,2n+3) series: coset transfer at every boundary level]{The \texorpdfstring{$(2,2n+3)$}{(2,2n+3)} series: coset transfer of graded unitarity at every boundary level}\label{ch:L29}
\begingroup\setlength{\emergencystretch}{3em}\sloppy

\providecommand{\ltwentyninesl}{\mathfrak{sl}_2}
\providecommand{\ltwentyninegr}{\operatorname{gr}}
\providecommand{\ltwentyninevac}{|0\rangle}
\providecommand{\ltwentynineGU}{\mathrm{GU}}
\providecommand{\ltwentynineJ}{J_\infty}
\providecommand{\ltwentynineinv}{\mathcal I}
\providecommand{\ltwentyninegc}{g}
\providecommand{\ltwentynineRn}{R_n}
\providecommand{\ltwentynineS}{S_{n,r}}
\providecommand{\ltwentynineR}{\mathbb{R}}
\providecommand{\ltwentynineln}[1]{\nolinkurl{#1}}

This chapter was written against an earlier state of the LeanVOA library. It covers the files
\leanpath{VOA/Examples/TwoSeries*.lean} (128 files in this release). Every Lean quotation below was copied from the source of
that state, and for this edition the book's automatic check confirms that every one occurs verbatim (up to whitespace and
comments) in the source files of this release. What the machine confirmed is collected in one place,
Section~\ref{l29:sec:verification}.

\paragraph{Conventions for quotations and status boxes.}
\begin{itemize}
\item \emph{Omitted proofs and docstrings.} A Lean quotation shows a statement, not its proof. A block that ends in
  \lean{:=} or \lean{:= by} with nothing after it is a statement whose proof is omitted; it is not a truncated
  quotation. Documentation comments (\lean{/-- ... -/}) are usually omitted; they are not part of the statement.
\item \emph{``Axioms standard''.} This means that \texttt{\#print axioms} reports exactly
  \lean{[propext, Classical.choice, Quot.sound]}, the three standard axioms of Lean's logic (Chapter~\ref{ch:01},
  Section~\ref{ch01:sec:trust}). Since \lean{VOA.pbw} became a theorem the library has no project
  axiom at all. A clean axiom list does \emph{not} mean ``unconditional'': \texttt{\#print axioms} does not list the
  explicit hypotheses of a theorem (Section~\ref{ch01:sec:uncond}). Every status box below therefore names the
  hypotheses separately, and the word \textbf{unconditional} is reserved for a statement with no hypothesis beyond the
  natural-number parameters.
\item \emph{Re-checked.} Every axiom claim in this chapter was re-run on a scratch file against the library's build
  products for this edition; the list is in Section~\ref{l29:sec:verification}.
\end{itemize}

\paragraph{Prerequisites.} This chapter is the level-generic sequel of Chapter~\ref{ch:L28}, which treats the
first member of the series, the Lee--Yang model, in full. We do not re-derive the Lee--Yang chain; we explain what
changes when the level moves, and we present in detail the parts that are new. We use without re-development:
(1) the graded-unitarity predicate of Arabi Ardehali, Beem, Lemos and Rastelli, its layers and the Virasoro length
flag (Chapter~\ref{ch:L19}, Definitions~\ref{l19:def:gu} and \ref{l19:def:virgu}), the boundary central charges
$c_{2,2m+1}$ and their real structures (Definition~\ref{l19:def:boundarycharge}, Theorem~\ref{l19:thm:boundaryreal});
(2) the boundary admissible levels $k_n=-2+2/(2n+1)$ of $\ltwentyninesl$, the universal and simple affine VOAs
$V^k(\ltwentyninesl)$, $L_k(\ltwentyninesl)$ and their vacuum characters (Chapter~\ref{ch:L20},
Definitions~\ref{l20:def:level}, \ref{l20:def:simple} and \ref{l20:def:character}); (3) the Lee--Yang coset chain of
Chapter~\ref{ch:L28}: the BRST complex at total level $-4$, the Beem--Garner Hodge package, the Koszul step, the
dimension squeeze, and the Rogers--Ramanujan character. Mathlib notions (\lean{Submodule}, \lean{Module.finrank},
\lean{Module.rank}, \lean{MvPolynomial}, \lean{PowerSeries}, \lean{Ideal.span}) are explained in plain words the
first time they appear. \textbf{Reading the Lean.} A reader who has never read Lean should first read the primer of
Chapter~\ref{ch:01}, Section~\ref{ch01:sec:lean}; the dictionaries below assume it.

\paragraph{What this chapter is about, in one paragraph.} For every $n\ge1$ let $k_n=-2+\frac{2}{2n+1}$. The
\emph{coset} $(L_{k_n}(\ltwentyninesl)\otimes L_1(\ltwentyninesl))/L_{k_n+1}(\ltwentyninesl)$ contains a Virasoro
algebra of central charge
\begin{equation}\label{l29:eq:charge}
c_{2,2n+3}\;=\;\frac{3k_n}{k_n+2}+1-\frac{3(k_n+1)}{k_n+3}\;=\;1-\frac{3(2n+1)^2}{2n+3},
\end{equation}
the central charge of the non-unitary minimal model $M(2,2n+3)$. Chapter~\ref{ch:L28} proved, at $n=1$
($k=-4/3$, $c=-22/5$), that \emph{graded unitarity} of the affine VOA implies graded unitarity of the Lee--Yang model.
This chapter proves the same implication for every $n$:
\begin{equation}\label{l29:eq:headline}
\ltwentynineGU\bigl(L_{k_n}(\ltwentyninesl)\bigr)\ \Longrightarrow\ \ltwentynineGU\bigl(M(2,2n+3)\bigr)\qquad(n\ge1),
\end{equation}
as the Lean theorem \ltwentynineln{VOA.TwoSeries.JetPoincareTwo.twoSeries_gradedUnitary_of_affine_final}
(Theorem~\ref{l29:thm:final}). Its only hypothesis is graded unitarity of the affine VOA, which is \textbf{open}; the
theorem is a \emph{transfer}, not a proof that either side is graded unitary.

\paragraph{Why formalize this.} Graded unitarity is the shadow, on a chiral algebra, of unitarity of a
four-dimensional $\mathcal N=2$ superconformal field theory. The level-$k_n$ affine VOA is the chiral algebra of the
Argyres--Douglas theory $(A_1,D_{2n+1})$, and $M(2,2n+3)$ is the chiral algebra of $(A_1,A_{2n})$; so
\eqref{l29:eq:headline} says that a positivity property passes from one family of four-dimensional theories to the
other, along the coset. The proof at $n=1$ (Chapter~\ref{ch:L28}) used one explicit model of the associated graded
of $L_{-4/3}$ and one cyclic dual vector. Neither generalizes (warning on p.~\pageref{l29:warn:cyclic}). The general proof
needed four new ingredients, each of independent interest, and each formalized here with no hypothesis: the vacuum
characters of all $M(2,2n+3)$ via the Andrews--Gordon identity; a jet model of the associated graded of $L_{k_n}$ with
the Chern--Tran--Wakhare fermionic character; a Drinfeld--Sokolov ``slice map'' with the descending half of Arakawa's
theorem; and an ideal-membership theorem for the $n$-equal arrangement (block generators span the relevant ideal).

\begin{remark}[The logical arc and the ledger]\label{l29:rem:arc}
Write $\ltwentynineGU(A)$ for graded unitarity of the affine factor $A=L_{k_n}(\ltwentyninesl)$ (with its canonical
form, compact conjugation and corrected flag, Definition~\ref{l29:def:endpoint}) and $\ltwentynineGU(T)$ for that of
the target $T=L(c_{2,2n+3},0)$ (partition-basis real structure, length flag). The Lean proof of
\eqref{l29:eq:headline} is the chain
\begin{equation}\label{l29:eq:chain}
\begin{aligned}
&\text{(1)}\quad \ltwentynineGU(A)\wedge\mathrm{Hodge}\wedge\mathrm{Squeeze}(\mathrm{coeff})\ \Longrightarrow\ \ltwentynineGU(T);\\
&\text{(2)}\quad \bigl(\dim\ltwentynineinv_{r,N}\le\mathrm{coeff}(N,r)\bigr)\wedge
  \bigl(\textstyle\sum_{r\le p}\mathrm{coeff}(N,r)\le\dim F_pT_N\bigr)\ \Longrightarrow\ \mathrm{Squeeze}(\mathrm{coeff});\\
&\text{(3)}\quad \textstyle\sum_{r\le p}\ltwentyninegc_n(N,r)\le\dim F_pT_N\quad(\text{take }\mathrm{coeff}=\ltwentyninegc_n);\\
&\text{(4)}\quad \text{strict model maps}\wedge\text{model socle data}\ \Longrightarrow\ \text{socle data of }\ltwentyninegr_rA
  \ \Longrightarrow\ \dim\ltwentynineinv_{r,N}\le\ltwentyninegc_n(N,r);\\
&\text{(5)}\quad \mathrm{Detect}_n\wedge\mathrm{ModelChar}_n\ \Longrightarrow\ \text{strict model maps};\\
&\text{(6)}\quad \text{slice map}\wedge\text{SliceArakawa}\wedge\text{VirSymbol}\ \Longrightarrow\ \mathrm{Detect}_n;\\
&\text{(7)}\quad \text{$n$-equal arrangement inclusion}\ \Longrightarrow\ \text{jet-ideal Poincar\'e lemma}
  \ \Longrightarrow\ \text{model socle data}.
\end{aligned}
\end{equation}
Step (1) is Theorem~\ref{l29:thm:transfer}, with the Hodge package supplied by Theorem~\ref{l29:thm:hodge}; (2) is
Theorem~\ref{l29:thm:invariants}; (3) is Theorem~\ref{l29:thm:hpbw}; (4) is Section~\ref{l29:sec:model}
(Theorems~\ref{l29:thm:soclespectrum} and \ref{l29:thm:swc}); (5) is Theorem~\ref{l29:thm:strict}, with
$\mathrm{ModelChar}_n$ proved in Theorem~\ref{l29:thm:modelchar}; (6) is Section~\ref{l29:sec:ds}, all three inputs
proved; (7) is Sections~\ref{l29:sec:poincare} and \ref{l29:sec:arr}, the arrangement inclusion proved in
Theorem~\ref{l29:thm:geninclusion}.
Here $\ltwentynineinv_{r,N}$ is the space of $\ltwentyninesl[t]$-invariant tensors of weight $N$ in
$\ltwentyninegr_rA\otimes L_1(\ltwentyninesl)$ (Definition~\ref{l29:def:invariants}), $F_pT_N$ is the length-$\le p$
part of $T_N$, and $\ltwentyninegc_n(N,r)$ is the number of partitions of $N$ into $r$ parts $\ge2$ with
$f_m+f_{m+1}\le n$ (Definition~\ref{l29:def:gordon}).

The ledger:
\begin{center}\small
\begin{tabular}{p{0.38\textwidth}p{0.14\textwidth}p{0.38\textwidth}}
\toprule
step & form & status\\
\midrule
level data, coset charge \eqref{l29:eq:charge} & equality & proved, every $n$\\
GU of the auxiliary factor $V^{k_U}$ & predicate & proved, every $n$ ($k_U<-2$)\\
Hodge package (Beem--Garner) & construction & proved from $\ltwentynineGU(A)$, every $n$\\
Koszul embedding $H\hookrightarrow\ltwentynineinv$ & injection & proved from $\ltwentynineGU(A)$, every $n$\\
vacuum character of $M(2,2n+3)$ & equality & proved, every $n$ (Andrews--Gordon)\\
refined PBW lower bound & $\le$ & proved, every $n$\\
socle-word count $=$ Gordon count & equality & proved, every $n$\\
model character bound & $\le$ & proved, every $n\ge1$ (Bailey chain)\\
detection (slice map + Arakawa half) & vanishing & proved, every $n\ge1$\\
model socle data (transversality; box count) & $\le$, disjoint & proved, every $n\ge1$\\
jet-ideal Poincar\'e lemma, arrangement inclusion & containment & proved, every $n\ge1$\\
$\ltwentynineGU(A)$ & predicate & \textbf{the one hypothesis; open}\\
\bottomrule
\end{tabular}
\end{center}
So the headline is \texttt{\#print axioms}-clean (standard axioms), \emph{conditional on $\ltwentynineGU(A)$ alone}.
\end{remark}

\paragraph{The findings of this chapter, in brief.}
\begin{enumerate}
\item \textbf{The transfer holds at every boundary level.} \ltwentynineln{twoSeries_gradedUnitary_of_affine_final}
  (Theorem~\ref{l29:thm:final}) concludes graded unitarity of $M(2,2m+5)$ from graded unitarity of
  $L_{-2+2/(2m+3)}(\ltwentyninesl)$, for every $m\in\N$, with no other hypothesis. At $m=0$ it is the Lee--Yang
  theorem of Chapter~\ref{ch:L28} again, by a different route.
\item \textbf{Several standalone theorems, all unconditional:} the vacuum character of every $M(2,2n+3)$
  (Theorem~\ref{l29:thm:character}); the model character bound (Theorem~\ref{l29:thm:modelchar}); the descending half
  of Arakawa's theorem for the Drinfeld--Sokolov slice map, at every level datum whose affine level is not a natural number,
  hence at every $k_n$ (Theorem~\ref{l29:thm:slicearakawa}); the generation of $L_N\cap(\text{vanishing ideal of the $n$-equal arrangement})$ by explicit
  block generators (Theorem~\ref{l29:thm:geninclusion}).
\item \textbf{Two declarations share the short name} \ltwentynineln{twoSeries_gradedUnitary_of_affine}, in different
  namespaces and with different hypotheses (warning on p.~\pageref{l29:warn:clash}). Neither is the headline. Always read the
  fully qualified name.
\item \textbf{Several docstrings are stale}: they still call ``open'' inputs that later files prove
  (warning on p.~\pageref{l29:warn:stale}). The Lean statements are correct; only the prose lags.
\item \textbf{What this chapter does not prove}: graded unitarity of any affine VOA, or of any $M(2,2n+3)$ with
  $n\ge1$. Both remain open (Kaufmann's conjecture for the minimal models, Chapter~\ref{ch:L19}).
\end{enumerate}

\paragraph{Notation used in this chapter.} The sources use three different indices for the same family; we fix
them here.
\begin{center}\small
\begin{longtable}{p{0.24\textwidth}p{0.68\textwidth}}
\toprule
symbol & meaning\\
\midrule\endhead
$n\ge1$ & the member index of the series: level $k_n=-2+\frac2{2n+1}$, target $M(2,2n+3)$. Lean
  \lean{twoSeriesLevels n}.\\
$m\ge0$, $n=m+1$ & the index of the headline theorem: level $-2+\frac2{2m+3}$, target $M(2,2m+5)$. Lean
  \lean{twoSeriesLevels (m + 1)}. Used to exclude $n=0$.\\
$c_{2,2j+1}$ & $1-\frac{3(2j-1)^2}{2j+1}$, Lean \lean{boundaryCharge j}; so the target is $c_{2,2n+3}$
  $=$ \lean{boundaryCharge (n + 1)}, and the field \lean{CosetLevels.m} equals $n+1$.\\
$q=2n+1$ & the denominator of the level; also the formal variable of $q$-series when no confusion is possible.
  In Sections~\ref{l29:sec:char} and \ref{l29:sec:modelchar} $q$ is the formal variable of power series.\\
$k=k_n$, $k_U$ & the boundary level and the auxiliary level $k_U=-5-k_n=-3-\frac{2}{2n+1}<-2$.\\
$A$, $B$, $U$ & $A=L_{k}(\ltwentyninesl)$ (simple), $B=L_1(\ltwentyninesl)$, $U=V^{k_U}(\ltwentyninesl)$ (universal).\\
$T$ & the target $L(c_{2,2n+3},0)$, the simple Virasoro VOA.\\
$F_pT_N$ & the span of Virasoro words $L_{-j_1-2}\cdots L_{-j_\ell-2}\ltwentyninevac$ with $\ell\le p$ and weight $N$
  (the length flag, Definition~\ref{l19:def:lengthflag}).\\
$\ltwentyninegr_rA$ & the $r$-th layer of the corrected PBW filtration of $A$ (Chapter~\ref{ch:L28}).\\
$\ltwentynineinv_{r,N}$ & the $\ltwentyninesl[t]$-invariant tensors of weight $N$ in $\ltwentyninegr_rA\otimes B$
  (Definition~\ref{l29:def:invariants}).\\
$\ltwentyninegc_n(N,r)$ & the Andrews--Gordon count (Definition~\ref{l29:def:gordon}).\\
$\ltwentynineRn$ & $\C[T_0,T_1,\dots]/(\text{jets of }T^n)$, the arc algebra of $M(2,2n+1)$ (Definition~\ref{l29:def:Rn}).
  Note the shift: $\ltwentynineRn$ belongs to $M(2,2n+1)$, the Drinfeld--Sokolov image of $A$, not to the target
  $M(2,2n+3)$.\\
$\ltwentynineS$ & the model: products of $r$ model letters in $\ltwentynineRn[u_j,v_j]$
  (Definition~\ref{l29:def:model}).\\
$e,f,h$; $J^a_m$ & the basis of $\ltwentyninesl$ and the current modes; in Lean the letters are indexed by
  \lean{Fin 3} with $0=e$, $1=f$, $2=h$, and \lean{J k a m} is the mode $J^a_m$ on $V^k$.\\
$\ltwentynineJ$ & the arc (infinite jet) functor.\\
\bottomrule
\end{longtable}
\end{center}

%======================================================================
\section{The mathematics}\label{l29:sec:math}
%======================================================================

\subsection{The boundary levels and the \texorpdfstring{$(2,2n+3)$}{(2,2n+3)} series}

An affine level $k$ of $\ltwentyninesl$ is \emph{admissible} when $k+2=p/p'$ with coprime integers $p\ge2$,
$p'\ge1$. The admissible levels with the smallest possible numerator, $p=2$, are the \emph{boundary} levels
$k=-2+2/p'$; for $p'$ odd we write $p'=2n+1$ and $k_n=-2+2/(2n+1)$. These are the levels of Chapter~\ref{ch:L20},
where the vacuum character of $L_{k_n}(\ltwentyninesl)$ was computed and the VOA was shown to be classically free.
Two other levels attach to $k_n$ in this chapter:
\begin{itemize}
\item the level $k_n+1$ of the diagonal $\ltwentyninesl$ inside $L_{k_n}\otimes L_1$, and
\item the auxiliary level $k_U=-5-k_n$, chosen so that $(k_n+1)+k_U=-4=-2h^\vee$: the total level at which the
  semi-infinite BRST charge of $\ltwentyninesl$ squares to zero.
\end{itemize}

\begin{example}[Level data for $n=1,2,3$]\label{l29:ex:levels}
With $c(k)=3k/(k+2)$ the Sugawara central charge of $\ltwentyninesl$ at level $k$:
\begin{center}\small
\begin{tabular}{cccccccc}
\toprule
$n$ & $k_n$ & $k_n+1$ & $k_U$ & $c(k_n)$ & $c(k_n+1)$ & $c_{2,2n+3}$ & DS image of $A$\\
\midrule
1 & $-4/3$ & $-1/3$ & $-11/3$ & $-6$ & $-3/5$ & $-22/5$ & $M(2,3)$, $c=0$\\
2 & $-8/5$ & $-3/5$ & $-17/5$ & $-12$ & $-9/7$ & $-68/7$ & $M(2,5)$, $c=-22/5$\\
3 & $-12/7$ & $-5/7$ & $-23/7$ & $-18$ & $-5/3$ & $-46/3$ & $M(2,7)$, $c=-68/7$\\
\bottomrule
\end{tabular}
\end{center}
For $n=2$: $c(-8/5)=3(-8/5)/(2/5)=-12$, $c(-3/5)=3(-3/5)/(7/5)=-9/7$, and $-12+1+9/7=-68/7=1-3\cdot25/7$. Also
$(k_2+1)+k_U=-3/5-17/5=-4$. In general $k_n=-4n/(2n+1)$ and $k_n+2=2/(2n+1)$, so
$c(k_n)=3k_n(2n+1)/2=-6n$, and
$c(k_n)+1-c(k_n+1)=1-\frac{3(2n+1)^2}{2n+3}$ (Theorem~\ref{l29:thm:charge}). The last column is the
Drinfeld--Sokolov reduction of $A$: $H_{DS}(L_{k_n}(\ltwentyninesl))=L(c_{2,2n+1},0)$ (Arakawa), with
$c_{\mathrm{DS}}(k)=1-6(k+1)^2/(k+2)$; it reappears in Section~\ref{l29:sec:ds}.
\end{example}

\begin{physics}
Four-dimensional $\mathcal N=2$ superconformal theories have chiral algebras (Beem, Lemos, Liendo, Peelaers, Rastelli
and van Rees). The Argyres--Douglas theory $(A_1,D_{2n+1})$ has flavour symmetry $SU(2)$ and chiral algebra
$L_{k_n}(\ltwentyninesl)$; the Argyres--Douglas theory $(A_1,A_{2n})$ has no flavour symmetry and chiral algebra
$M(2,2n+3)$. Their central charges are $c_{2d}=-12c_{4d}$. The coset of this chapter relates the two families, and
the level $-4=-2h^\vee$ of the BRST step is the level of a gauged $SU(2)$: the coset is, physically, the gauging of
the diagonal flavour symmetry with a compensating auxiliary sector.
\end{physics}

\subsection{The coset and its BRST form}\label{l29:sec:coset}

The GKO construction says that the commutant of the diagonal $\ltwentyninesl$ in $L_k\otimes L_1$ contains the
difference of Sugawara stress tensors $T_{\mathrm{coset}}=T_{(k)}+T_{(1)}-T_{(k+1)}$, a Virasoro field of central
charge $c(k)+1-c(k+1)$. At admissible $k=-2+p/p'$ this coset is the minimal model $M(p,p+p')$; at $p=2$,
$p'=2n+1$ it is $M(2,2n+3)$. The Lean development never identifies the coset VOA with the minimal model. It works
with the \emph{relative semi-infinite BRST complex}
\begin{equation}\label{l29:eq:brst}
C\;=\;\bigl(A\otimes B\otimes U\bigr)\otimes\mathcal G,\qquad \text{total level }(k+1)+k_U=-4,
\end{equation}
where $\mathcal G$ is the $bc$-ghost system of $\ltwentyninesl$ and $U=V^{k_U}(\ltwentyninesl)$ is universal. The
Virasoro algebra maps into $H^0(C)$ through $T_{\mathrm{coset}}$, and its kernel lies in the maximal ideal because
the vacuum class survives; so the \emph{simple} quotient $T=L(c,0)$ is what is reached, with no simplicity statement
about the coset needed. This is exactly the architecture of Chapter~\ref{ch:L28}; the level-generic files
reproduce it with the numbers $-4/3$, $-11/3$, $-22/5$ replaced by the fields of a structure
(Definition~\ref{l29:def:cosetlevels}).

\subsection{What is transferred, and how}\label{l29:sec:howtransfer}

Graded unitarity of a VOA $W$ with invariant form $B$, real structure $\rho$ and filtration $F$
(Chapter~\ref{ch:L19}, Definition~\ref{l19:def:gu}) asks, weight by weight, that the layers
$F_p\cap F_{p-1}^\perp$ be nondegenerate and that the sign-corrected Hermitian form be positive definite on each
layer. The transfer argument of Chapter~\ref{ch:L28}, which this chapter re-runs at every level, has three parts.
\begin{enumerate}
\item \textbf{Hodge theory.} Positivity of the three factors $A$, $B$, $U$ gives a positive Hermitian form on each
  finite piece of the relative complex $C$; with the Beem--Garner Kähler package (a BRST adjoint and an
  $\mathfrak{sp}_2$-type Lefschetz structure) every cohomology class has a unique harmonic representative and the
  form restricted to harmonic classes is positive in the graded sense. The positivity of $A$ is the hypothesis; the
  positivity of $B=L_1$ (genuine unitarity) and of $U$ (below the critical level) are theorems.
\item \textbf{Squeeze.} The images of Virasoro words of length $\le p$ in $H^0$ have harmonic representatives of
  ``Kähler degree'' $\le 2p$. If one knows that the leading cohomology in Kähler degree $2r$ has dimension at most
  $\mathrm{coeff}(N,r)$, and that $T$ has at least $\sum_{r\le p}\mathrm{coeff}(N,r)$ independent words of length
  $\le p$ at weight $N$, then the two filtrations (length on $T$, Kähler degree on $H^0$) coincide dimensionally,
  and the positivity of harmonic classes is exactly clause (vi) of graded unitarity for the length flag.
\item \textbf{The dimension inputs.} (I) the leading cohomology is bounded by an \emph{affine} invariant count
  (Koszul step); (II) the PBW lower bound is a statement about $T$ alone.
\end{enumerate}
At $n=1$ both dimension inputs were proved with $\mathrm{coeff}(N,r)=$ the number of partitions of $N$ into $r$
parts differing by at least $2$ (Rogers--Ramanujan). For general $n$ the coefficient is the Andrews--Gordon count
$\ltwentyninegc_n(N,r)$, (II) follows from the vacuum character of $M(2,2n+3)$, and (I) is the hard new input.

\subsection{What changes with \texorpdfstring{$n$}{n}}\label{l29:sec:whatchanges}

\begin{center}\small
\begin{tabular}{p{0.24\textwidth}p{0.32\textwidth}p{0.34\textwidth}}
\toprule
ingredient & Lee--Yang ($n=1$, Chapter~\ref{ch:L28}) & general $n$ (this chapter)\\
\midrule
levels & literal $-4/3$, $-11/3$, $-22/5$ & fields of \lean{CosetLevels}\\
BRST complex, coset, Hodge, Koszul & concrete files & level-generic ports, same proofs\\
coefficient $\mathrm{coeff}(N,r)$ & parts differing by $\ge2$ & Andrews--Gordon count $\ltwentyninegc_n(N,r)$\\
character of the target & Rogers--Ramanujan & Andrews--Gordon identity (Sec.~\ref{l29:sec:char})\\
model of $\ltwentyninegr_rA$ & jet ring of $u^2,uv,v^2,u'v-uv'$ & same letters over $\ltwentynineRn=\C[T_j]/(T^n)_\infty$\\
dual generators & one cyclic vector $w_r$ & a family indexed by socle words\\
detection & level-one vacuum functional & Drinfeld--Sokolov slice map\\
commutative algebra & none needed & jet-ideal Poincar\'e lemma, $n$-equal arrangement\\
\bottomrule
\end{tabular}
\end{center}
The new coefficient ring $\ltwentynineRn$ deserves a word. At $n=1$, $R_1=\C$ (the jets of $T^1$ are all the
$T_j$), and the model is the jet ring of the Lee--Yang chapter. For $n\ge2$ the stress-tensor letter of the model
carries a nilpotent ``tail'' $T_j\in\ltwentynineRn$: the model knows that $A$ reduces, under Drinfeld--Sokolov, to
$M(2,2n+1)$, whose arc algebra is $\ltwentynineRn$ (classical freeness, Chapter~\ref{ch:L20}).

%======================================================================
\section{The level data}\label{l29:sec:levels}
%======================================================================

The whole level-generic development is parametrized by one structure that bundles the four levels and every
arithmetic fact about them that any later proof uses. This is a design decision worth understanding: the
Lee--Yang files of Chapter~\ref{ch:L28} were written with literal numbers ($-4/3$, $-11/3$, $-22/5$), and every
place where a proof used arithmetic (``$-4/3+1+(-11/3)=-4$'') is now a field of the structure.

\begin{definition}[Coset level data]\label{l29:def:cosetlevels}
A \emph{coset level datum} consists of complex numbers $k_A$ (the level of the boundary factor), $k_{\mathrm{act}}$
(the index of its current action), $k_U$ (the auxiliary level), a real number $k_U^{\ltwentynineR}$, a target central charge
$c$, and a natural number $m$, such that: $k_{\mathrm{act}}=k_A$; $k_A+2\ne0$; $\overline{k_A}=k_A$;
$k_U=k_U^{\ltwentynineR}$ and $k_U^{\ltwentynineR}<-2$; $k_U+2\ne0$; $k_{\mathrm{act}}+3\ne0$; $k_{\mathrm{act}}+1+k_U=-4$;
$\frac{3k_A}{k_A+2}+1-\frac{3(k_{\mathrm{act}}+1)}{k_{\mathrm{act}}+3}=c$; and $c_{2,2m+1}=c$.
\end{definition}

\noindent File \leanpath{VOA/\allowbreak Examples/\allowbreak TwoSeriesLevels.lean}.
\begin{leancode}
structure CosetLevels where
  /-- Level of the boundary affine carrier. -/
  kA : ℂ
  /-- Index of the boundary current action (equal to `kA`). -/
  kAct : ℂ
  /-- Level of the auxiliary universal factor. -/
  kU : ℂ
  /-- The auxiliary level as a real number. -/
  kUr : ℝ
  /-- Target central charge. -/
  c : ℂ
  /-- Index of `c` in the boundary family `c_{2,2m+1}`. -/
  m : ℕ
  kAct_eq : kAct = kA
  aNoncritical : kA + 2 ≠ 0
  aReal : (starRingEnd ℂ) kA = kA
  kU_eq : kU = (kUr : ℂ)
  kUr_lt : kUr < -2
  uNoncritical : kU + 2 ≠ 0
  diagNoncritical : kAct + 1 + 2 ≠ 0
  /-- The total matter level cancels the ghost level `-4`. -/
  sum : kAct + 1 + kU = -4
  /-- Coset central charge `c(k) + c(1) - c(k+1)`. -/
  charge : 3 * kA / (kA + 2) + 1 - 3 * (kAct + 1) / (kAct + 1 + 2) = c
  /-- The target is the `m`th boundary minimal model. -/
  boundary : boundaryCharge m = c
\end{leancode}

\begin{dictionary}
\lean{structure ... where} declares a record type; each line \lean{name : type} is a field. Fields whose type is a
proposition (\lean{kAct_eq}, \lean{aNoncritical}, \dots) are \emph{proofs} that the record must carry, so an
inhabitant of \lean{CosetLevels} is a tuple of numbers \emph{together with} proofs of the listed identities.
\lean{starRingEnd ℂ} is complex conjugation, so \lean{aReal} says $k_A\in\ltwentynineR$. \lean{(kUr : ℂ)} is the coercion
$\ltwentynineR\hookrightarrow\C$. \lean{boundaryCharge m} is $c_{2,2m+1}=1-3(2m-1)^2/(2m+1)$ (Chapter~\ref{ch:L19},
Definition~\ref{l19:def:boundarycharge}).

\textbf{Why two copies of the level.} \lean{kA} is the level at which the carrier $L_{k_A}$ is built, \lean{kAct}
the number that indexes its current action in the diagonal sum. They are provably equal (\lean{kAct_eq}). The
duplication exists only to make the Lee--Yang instance reproduce the existing concrete objects \emph{definitionally}:
there the carrier was built at \lean{level 1} (an expression equal to $-4/3$) and the action at the literal
\lean{-4 / 3}. Mathematically it is one level.

\textbf{Each hypothesis.} $k_A+2\ne0$, $k_U+2\ne0$, $k_{\mathrm{act}}+3\ne0$: the three Sugawara constructions
(levels $k_A$, $k_U$, $k_A+1$) are non-critical. $k_A\in\ltwentynineR$: needed for the compact real structure on $A$.
$k_U<-2$ real: needed for the auxiliary positivity (Theorem~\ref{l29:thm:aux}). The sum $=-4$: nilpotency of the BRST
charge. \lean{charge}: the coset central charge is $c$. \lean{boundary}: $c$ is a boundary minimal-model charge, which
is what selects the real structure on the target. Nothing in the record is a hidden assumption about any VOA; it is
pure arithmetic, and it is inhabited for every $n$ (Definition~\ref{l29:def:twolevels}), so it is not vacuous.

\textbf{Index.} The field \lean{m} is the index of $c$ in the family $c_{2,2m+1}$; for the $n$-th member of the series
it is $n+1$ (Notation table).
\end{dictionary}

\begin{physics}
A coset level datum is the list of numbers that make a gauging consistent: the gauged diagonal $SU(2)$ current
algebra sits at level $k+1$, the auxiliary sector at $k_U$, and their sum must cancel the ghost contribution $-4$
(the level-matching, or anomaly-cancellation, condition of a $G/G$-type BRST reduction). The charge field is the GKO
central-charge bookkeeping $c_{\mathrm{coset}}=c_{\mathrm{UV}}-c_{\mathrm{gauged}}$.
\end{physics}

\begin{status}
Definition. Inhabited for every $n$ (below), so no theorem quantifying over it is vacuous.
\end{status}

\begin{theorem}[The coset central charge]\label{l29:thm:charge}
For every $n\in\N$, with $k=k_n=-2+\frac{2}{2n+1}$,
\begin{equation*}
\frac{3k}{k+2}+1-\frac{3(k+1)}{k+3}\;=\;c_{2,2n+3}\;=\;1-\frac{3(2n+1)^2}{2n+3}.
\end{equation*}
\end{theorem}

\noindent File \leanpath{VOA/\allowbreak Examples/\allowbreak TwoSeriesLevels.lean}.
\begin{leancode}
theorem twoSeries_charge (n : ℕ) :
    3 * level n / (level n + 2) + 1 - 3 * (level n + 1) / (level n + 1 + 2) =
      boundaryCharge (n + 1) := by
\end{leancode}

\begin{dictionary}
\lean{level n} is $k_n=-2+2/(2n+1)$ as a complex number (Chapter~\ref{ch:L20}, Definition~\ref{l20:def:level};
\lean{level_eq} restates it with the denominator written out). \lean{boundaryCharge (n + 1)} is
$c_{2,2(n+1)+1}=c_{2,2n+3}=1-3(2n+1)^2/(2n+3)$. Division in \lean{ℂ} is total (with $x/0=0$), so one must check that no
denominator vanishes: $k_n+2=2/(2n+1)\ne0$ and $k_n+3=(2n+3)/(2n+1)\ne0$ (\lean{diag_noncritical}); hence the
identity is the genuine rational identity. Same as the textbook statement, for every $n\in\N$ including $n=0$
($k_0=0$, $c_{2,3}=0$: the trivial coset $L_0\otimes L_1/L_1$).
\end{dictionary}

\begin{proof}[Proof on paper]
Put $q=2n+1$. Then $k+2=2/q$, $k=-2(q-1)/q$, $k+1=(2-q)/q$, $k+3=(q+2)/q$. So $3k/(k+2)=-3(q-1)$ and
$3(k+1)/(k+3)=3(2-q)/(q+2)$. Hence the left side is
$-3q+3+1-\frac{3(2-q)}{q+2}=\frac{(4-3q)(q+2)-6+3q}{q+2}=\frac{-3q^2+q+2}{q+2}$, and
$1-\frac{3q^2}{q+2}=\frac{q+2-3q^2}{q+2}$, the same.
\end{proof}

\begin{physics}
The GKO formula for the central charge of $(\widehat{\mathfrak{su}}(2)_k\times\widehat{\mathfrak{su}}(2)_1)/
\widehat{\mathfrak{su}}(2)_{k+1}$, here continued to the negative fractional levels of the Argyres--Douglas chiral
algebras. At $n=1$ it is the familiar $-6+1+3/5=-22/5$ of the Lee--Yang edge singularity.
\end{physics}

\begin{status}
Proved, unconditional. Axioms standard (re-checked).
\end{status}

\begin{definition}[The level data of the series]\label{l29:def:twolevels}
For $n\in\N$, \lean{boundaryFamilyLevels n} is the datum $k_A=k_{\mathrm{act}}=k_n$, $k_U=-5-k_n$,
$k_U^{\ltwentynineR}=-3-2/(2n+1)$, $c=c_{2,2n+3}$, $m=n+1$. \lean{leeYangLevels} is the same datum at $n=1$ written with the
literal numbers of Chapter~\ref{ch:L28}. \lean{twoSeriesLevels n} is \lean{boundaryFamilyLevels n}, except at $n=1$
where it is \lean{leeYangLevels}.
\end{definition}

\noindent File \leanpath{VOA/\allowbreak Examples/\allowbreak TwoSeriesLevels.lean}.
\begin{leancode}
def boundaryFamilyLevels (n : ℕ) : CosetLevels where
  kA := level n
  kAct := level n
  kU := -5 - level n
  kUr := uLevelReal n
  c := boundaryCharge (n + 1)
  m := n + 1
  kAct_eq := rfl
  aNoncritical := level_noncritical n
  aReal := level_real n
  kU_eq := uLevel_cast n
  kUr_lt := uLevelReal_lt n
  uNoncritical := uLevel_noncritical n
  diagNoncritical := diag_noncritical n
  sum := by ring
  charge := twoSeries_charge n
  boundary := rfl

def leeYangLevels : CosetLevels where
  kA := level 1
  kAct := -4 / 3
  kU := -11 / 3
  kUr := -11 / 3
  c := LeeYang.leeYangC
  m := 2
  kAct_eq := level_one.symm
  aNoncritical := level_noncritical 1
  aReal := level_real 1
  kU_eq := by push_cast; ring
  kUr_lt := by norm_num
  uNoncritical := by norm_num
  diagNoncritical := by norm_num
  sum := by norm_num
  charge := by rw [level_one]; norm_num [LeeYang.leeYangC]
  boundary := boundaryCharge_two_eq_leeYangC

def twoSeriesLevels : ℕ → CosetLevels
  | 1 => leeYangLevels
  | n => boundaryFamilyLevels n

theorem twoSeriesLevels_one : twoSeriesLevels 1 = leeYangLevels := rfl

theorem twoSeriesLevels_c (n : ℕ) : (twoSeriesLevels n).c = boundaryCharge (n + 1) := by
\end{leancode}

\begin{dictionary}
\lean{uLevelReal n} is $-3-2/(2n+1)\in\ltwentynineR$ and \lean{uLevelReal_lt} proves it is $<-2$. The proof fields are filled by
the lemmas of the same file: e.g. \lean{sum := by ring} is the identity $k+1+(-5-k)=-4$. \lean{LeeYang.leeYangC} is
$-22/5$. The definition of \lean{twoSeriesLevels} is by pattern matching: the first clause \lean{| 1 => ...} catches
$n=1$, the second every other $n$. \lean{twoSeriesLevels_one} holds by \lean{rfl}, i.e.\ by unfolding definitions; and
\lean{boundaryFamilyLevels_one_kA}, \dots, \lean{_one_m} (same file) check field by field that the two descriptions
of $n=1$ agree as numbers. \lean{twoSeriesLevels_c} gives the target charge for every $n$. Companion lemmas give
\lean{twoSeriesLevels_kA n = level n}, \lean{twoSeriesLevels_kU n = -3 - 2/(2n+1)} and \lean{twoSeriesLevels_m n = n + 1}.

\textbf{Why the special case at $n=1$.} So that every general theorem, specialised to $n=1$, is \emph{literally}
(definitionally) a statement about the Lee--Yang objects of Chapter~\ref{ch:L28}, with no transport lemma. This is a
choice of convenience; the numbers are the same either way.
\end{dictionary}

\begin{physics}
The family $n\mapsto(A_1,D_{2n+1})\to(A_1,A_{2n})$ of Argyres--Douglas pairs, with $n=1$ the pair $(A_1,D_3)=(A_1,A_3)$
and Lee--Yang $(A_1,A_2)$.
\end{physics}

\begin{status}
Definitions; the field lemmas are proved, unconditional. Axioms standard for \lean{twoSeriesLevels_c} (re-checked).
\end{status}

\begin{warning}[title={Index conventions, and the member $n=0$}]\label{l29:warn:index}
Three indices coexist (Notation table). (1) Theorems stated for \lean{(n : ℕ)} with no hypothesis $1\le n$ also cover
$n=0$: \lean{twoSeriesLevels 0} has $k_A=0$ and $c=c_{2,3}=0$. That instance is harmless (it is a statement about
the trivial coset and the trivial VOA $L(0,0)=\C$), but it is not a member of the series. (2) The headline and the
slice-map files use \lean{twoSeriesLevels (m + 1)} precisely to exclude it; their $m=0$ is $n=1$, Lee--Yang, and their
target $M(2,2m+5)$ is the $n$-th member. (3) The field \lean{CosetLevels.m} is $n+1$, not $n$ and not the $m$ of the
headline. (4) The coefficient ring $R_n$ of Section~\ref{l29:sec:model} uses the member index $n$ but describes
$M(2,2n+1)$, one step \emph{below} the target. A reader checking a statement should translate every index to $n$
first.
\end{warning}

%======================================================================
\section{The objects at the two ends}\label{l29:sec:endpoints}
%======================================================================

The hypothesis and the conclusion of every theorem in this chapter are instances of the graded-unitarity predicate
\lean{GradedUnitary B R F} of Chapter~\ref{ch:L19} (Definition~\ref{l19:def:gu}). Each instance needs an invariant
form $B$, a real structure $R$ and a flag $F$. For the Virasoro target the predicate is pinned as
\lean{VirGradedUnitary c R} (Definition~\ref{l19:def:virgu}). The level-generic file builds the three data for every
level datum.

\begin{definition}[The data on the boundary factor]\label{l29:def:endpoint}
For a level datum $L$, let $A=L_{k_A}(\ltwentyninesl)$ be the simple quotient of $V^{k_A}(\ltwentyninesl)$ by its
maximal graded ideal. Its data are: the canonical invariant form normalised by $(\ltwentyninevac|\ltwentyninevac)=1$;
the \emph{corrected current-and-stress flag} (the PBW-type filtration of Chapter~\ref{ch:L28} in which a current mode
$J^a_{-j}$ and a Sugawara mode $L_{-j}$ each raise the degree by one, descended to the quotient); and the compact
conjugation (the antilinear involution fixing the compact real form $\mathfrak{su}(2)$), descended from
$V^{k_A}$ to $A$.
\end{definition}

\noindent File \leanpath{VOA/\allowbreak Examples/\allowbreak TwoSeriesAffine.lean}.
\begin{leancode}
abbrev affineForm : InvariantForm (AVOA L) := simpleForm L.kA L.aNoncritical

def affineFlag : GradedFlag (AVOA L) :=
  CorrectedR.correctedAffineRQuotientFlag L.kA L.aNoncritical
    (QuotVOA.maxGradedIdeal (universal L.kA L.aNoncritical))
    (QuotVOA.maxGradedIdeal_isIdeal _)
    (QuotVOA.maxGradedIdeal_isGradedIdeal _)

def affineConj : ASpace L →ₛₗ[starRingEnd ℂ] ASpace L :=
  (QuotVOA.maxGradedIdeal (universal L.kA L.aNoncritical)).mapQ
    (QuotVOA.maxGradedIdeal (universal L.kA L.aNoncritical))
    (affineUniversalConj L) (affineUniversalConj_maxIdeal_stable L)
\end{leancode}

\begin{dictionary}
\lean{AVOA L} is \lean{AffineSl2GeneralRigidity.simple L.kA L.aNoncritical}, the maximal quotient
\lean{QuotVOA.maxQuotientVOA} of the universal affine VOA (Chapter~\ref{ch:L20}, Definition~\ref{l20:def:simple});
\lean{ASpace L} is its underlying space, the quotient of the Wakimoto-model space \lean{WakModule L.kA} by the maximal
graded ideal. So the hypothesis is about the \emph{actual simple} VOA, not a surrogate. \lean{→ₛₗ[starRingEnd ℂ]} is a
semilinear map for complex conjugation, i.e.\ an antilinear map; \lean{mapQ} descends a map to quotients, and
\ltwentynineln{affineUniversalConj_maxIdeal_stable} is the proof that the conjugation preserves the maximal ideal (any
conjugation-stable proper ideal is contained in the maximal one). The real structure \lean{affineRealStructure L}
bundles \lean{affineConj} with the proofs that it is antilinear, respects the vertex operation, fixes the vacuum and
the conformal vector, and squares to the identity (\lean{parity := fun _ => 1}: all states are even). The theorems
\lean{affineForm_leeYang}, \lean{affineFlag_leeYang}, \lean{affineRealStructure_leeYang} prove by \lean{rfl} that at
\lean{leeYangLevels} these are the objects of Chapter~\ref{ch:L28}.

The hypothesis of every transfer theorem is therefore
\begin{equation*}
\lean{hA : GradedUnitary (affineForm L) (affineRealStructure L) (affineFlag L)},
\end{equation*}
graded unitarity of the actual simple $L_{k_A}(\ltwentyninesl)$ with these three data.
\end{dictionary}

\begin{physics}
The flag is the chiral-algebra image of the four-dimensional $SU(2)_R$-charge filtration (the ``$R$-filtration'' of
ABLR), with currents of $R$-charge one and the stress tensor of $R$-charge one. The compact conjugation is the CPT
conjugation of the four-dimensional theory restricted to Schur operators. Graded unitarity of $L_{k_n}$ is then the
positivity that four-dimensional unitarity of $(A_1,D_{2n+1})$ is expected to force.
\end{physics}

\begin{status}
Definitions. The flag is the \emph{corrected} one of Chapter~\ref{ch:L28}; whether it is the physically correct
$R$-filtration is a modelling choice inherited from there, not something this chapter re-examines.
\end{status}

\begin{theorem}[Graded unitarity of the auxiliary factor]\label{l29:thm:aux}
For every level datum $L$, the universal affine VOA $V^{k_U}(\ltwentyninesl)$ with its canonical form, compact
conjugation and ordinary PBW flag is graded unitary.
\end{theorem}

\noindent File \leanpath{VOA/\allowbreak Examples/\allowbreak TwoSeriesAffine.lean}.
\begin{leancode}
abbrev auxiliaryForm : InvariantForm (UVOA L) := universalForm L.kU L.uNoncritical

abbrev auxiliaryRealStructure : RealStructure (UVOA L) :=
  compactRealStructure L.uReal L.uNoncritical

abbrev auxiliaryFlag : GradedFlag (UVOA L) := ordinaryFlag L.kU L.uNoncritical

theorem auxiliaryGradedUnitary :
    GradedUnitary (auxiliaryForm L) (auxiliaryRealStructure L) (auxiliaryFlag L) := by
\end{leancode}

\begin{dictionary}
\lean{UVOA L} is the universal $V^{k_U}(\ltwentyninesl)$ (not a quotient). The proof specialises the theorem
\ltwentynineln{universal_gradedUnitary_belowCritical_unconditional} of the below-critical files (Chapter~\ref{ch:L19}) at the
real level \lean{L.kUr < -2}. No hypothesis besides the level datum; and the datum's field \lean{kUr_lt} is exactly what
is used. Same as the textbook statement.
\end{dictionary}

\begin{proof}[Proof on paper]
On each weight space of $V^k(\ltwentyninesl)$ fix the PBW basis adapted to the ordinary flag. The entries of the Gram
matrix of the canonical form in that basis are polynomials in $k$. The below-critical files show that no leading
principal minor (in the flag-adapted order) vanishes for real $k<-2$ (``no wall''; this is consistent with the
irreducibility of the vacuum module there, whose Kac--Kazhdan reducibility walls lie at $k+2\in\Q_{>0}$, but the
statement needed is the finer one about flag steps). The signs of the minors, hence the signs demanded by
clause (vi), are therefore constant on the half-line, and they are computed at an anchor, in the regime
$k\to-\infty$ where the form is dominated by its leading power of $k$ and each layer behaves like a free
(Heisenberg-type) system whose sign on layer $p$ is $(-1)^p$, which is exactly the sign twist of clause (vi). The
compact conjugation turns the bilinear form into the Hermitian one with the same diagonal on real vectors. The
proof lives in the below-critical files (assigned to Chapter~\ref{ch:L19}); this chapter only specialises it.
\end{proof}

\begin{physics}
At level $k_U<-2$ the vacuum module is ``unitary up to the expected signs'': it behaves like a non-compact current
algebra. In the BRST picture it is the auxiliary sector that absorbs the anomaly of the gauged $SU(2)$; that it has
the right positivity is what lets positivity flow from $A$ alone.
\end{physics}

\begin{status}
Proved, unconditional (for every level datum). Axioms standard (re-checked).
\end{status}

\begin{definition}[The target real structure]\label{l29:def:target}
For a level datum $L$, \lean{targetRealStructure L} is the partition-basis real structure on $L(c,0)$ obtained from the
real Feigin--Fuchs parameter $\lambda$ with $1-12\lambda^2=c$ given by the boundary family (Chapter~\ref{ch:L19},
Theorem~\ref{l19:thm:boundaryreal}).
\end{definition}

\noindent File \leanpath{VOA/\allowbreak Examples/\allowbreak TwoSeriesAffine.lean}.
\begin{leancode}
def targetRealStructure : RealStructure (VirSimple.virasoroSimple L.c) :=
  virasoroSimpleRealStructure L.lam_spec (VirGradedUnitarity.boundaryBackgroundCharge_fixed L.m)
\end{leancode}

\begin{dictionary}
\lean{L.lam_spec} is the identity $1-12\lambda_m^2=c$ derived from the field \lean{boundary} (file
\leanpath{TwoSeriesLevels.lean}), where $\lambda_m=$ \lean{boundaryBackgroundCharge L.m} is real; \lean{_fixed} says
conjugation fixes it. The conclusion of every transfer theorem is
\lean{VirGradedUnitary L.c (targetRealStructure L)}, i.e.\ \lean{GradedUnitary (Bsimp c) R (virGradedFlag c)}:
graded unitarity of the actual simple Virasoro VOA with its canonical form and length flag. By
Theorem~\ref{l19:thm:rhostable} every step of the length flag is stable under every real structure, and the
reduction of Chapter~\ref{ch:L19} is insensitive to the choice of $\rho$.
\end{dictionary}

\begin{physics}
The real structure is the CPT conjugation of the Virasoro chiral algebra in the basis of partition (Feigin--Fuchs)
states; the boundary family is exactly the case where the background charge is real.
\end{physics}

\begin{status}
Definition. At \lean{leeYangLevels} it is \ltwentynineln{LeeYangTransferAssembly.leeYangRealStructure} (\lean{rfl}).
\end{status}

%======================================================================
\section{The level-generic coset chain}\label{l29:sec:chain}
%======================================================================

\subsection{The infrastructure files, described collectively}

About forty files port the Lee--Yang chain of Chapter~\ref{ch:L28} to an arbitrary level datum $L$. Each is a
re-run of a concrete file with the numbers replaced by the fields of $L$; each ends with \lean{rfl} lemmas showing
that at \lean{leeYangLevels} it reproduces the concrete object. We list them by role and do not re-present their
internal lemmas, which are the Lee--Yang lemmas with a different level.
\begin{center}\small
\begin{tabular}{p{0.38\textwidth}p{0.52\textwidth}}
\toprule
files & content\\
\midrule
\leanpath{TwoSeriesFactors} & the factors $A$, $U$; the diagonal current action on $A\otimes B\otimes U$ at level $-4$
  (\lean{matterAction}); the semi-infinite BRST complex \lean{brstComplex L}, with nonzero vacuum class\\
\leanpath{TwoSeriesCoset} & the coset stress tensor \lean{cosetOmega}; the coset VOA \lean{cosetVOA L} on the actual
  diagonal commutant, central charge \lean{L.c} from the field \lean{charge}; its map into the BRST complex\\
\leanpath{TwoSeriesImage} & \lean{virasoroToH0 L}: $V(c,0)\to H^0$, kernel inside the maximal ideal; the finite relative
  complexes \lean{relativeWeight L n}\\
\leanpath{TwoSeriesHodge}, \leanpath{TwoSeriesHodgeGeometry}, \leanpath{-Good}, \leanpath{-Goodness}, \leanpath{-Kahler},
  \leanpath{-BG}, \leanpath{-GhostConj}, \leanpath{-ConjIdentity}, \leanpath{-Harmonic}, \leanpath{-Inputs} & the
  Beem--Garner Kähler package on the actual relative complexes, its positivity from $\ltwentynineGU(A)$, the two BRST
  conjugation identities, the harmonic-metric bridges; assembled in \lean{brstHodgeInputs} (Theorem~\ref{l29:thm:hodge})\\
\leanpath{TwoSeriesAssembly} & the transfer theorem with two packaged inputs (Theorem~\ref{l29:thm:transfer})\\
\leanpath{TwoSeriesKoszulGr}, \leanpath{-U}, \leanpath{-Matter}, \leanpath{-Native}, \leanpath{-Homotopy},
  \leanpath{-Ident}, \leanpath{-Grading}, \leanpath{-Exact}, \leanpath{-Endpoint}, \leanpath{-Eval}, \leanpath{-SymA},
  \leanpath{-Symbol}, \leanpath{-SymbolMaps}, \leanpath{-Parity}, \leanpath{-Invariants} & the symbol module
  $\ltwentyninegr_rA$, the Koszul exactness of the leading BRST differential, the two symbol maps, the vanishing of odd
  and negative degrees; the leading Koszul embedding (Theorem~\ref{l29:thm:koszul})\\
\bottomrule
\end{tabular}
\end{center}
(All paths are \leanpath{VOA/Examples/<name>.lean}; ``\leanpath{-X}'' abbreviates the name with the preceding prefix.)
The definitions and theorems that a reader needs to check the headline statements are the following.

\subsection{The two packaged inputs and the transfer theorem}

\begin{definition}[The Hodge package]\label{l29:def:hodge}
For a level datum $L$, a \emph{BRST/Hodge package} consists of: a Kähler geometry on the finite relative BRST
complexes at each conformal weight (a positive form $P$ and a Kähler-type decomposition whose differential is the
actual BRST residue); a proof that $\ltwentynineGU(A)$ and $\ltwentynineGU(U)$ make every $P$ positive definite; the
graded-unitarity data on harmonic classes; a proof that the classes of Virasoro words of length $\le p$ lie in Kähler
degree $\le 2p$ (generator tests); and three bridges identifying the harmonic metric on the classes of Virasoro words
with the canonical form of $T$ (vacuum normalisation, adjointness of creation operators, reality of words).
\end{definition}

\noindent File \leanpath{VOA/\allowbreak Examples/\allowbreak TwoSeriesAssembly.lean}.
\begin{leancode}
structure BRSTHodgeInputs where
  kahler : BRSTKahlerGeometry L
  weakGU : GradedUnitary (affineForm L) (affineRealStructure L) (affineFlag L) →
    GradedUnitary (auxiliaryForm L) (auxiliaryRealStructure L) (auxiliaryFlag L) →
    ∀ n, IsPosOn (kahler.geometry n).P ⊤
  gu : BRSTGUData L kahler
  generators : ∀ hp : (∀ n, IsPosOn (kahler.geometry n).P ⊤),
    ClassGeneratorFiltration L kahler hp
  vacuumNormalized : HarmonicVacuumNormalized L kahler gu
  creationAdjoint : HarmonicCreationAdjoint L kahler gu
  wordsReal : HarmonicWordsReal L kahler gu
\end{leancode}

\begin{dictionary}
The seven fields are the seven named inputs of the Lee--Yang assembly of Chapter~\ref{ch:L28}, now at level $L$;
their types are defined in \leanpath{TwoSeriesHodge.lean} and are statements about the \emph{actual} relative complexes
\lean{relativeWeight L n} and the actual Virasoro classes, not about abstract surrogates. \lean{IsPosOn P ⊤} says
the Hermitian form $P$ is positive definite on the whole space (\lean{⊤} is the top submodule; Chapter~\ref{ch:L19},
Definition~\ref{l19:def:isposon}). Note that \lean{weakGU} is an \emph{implication} from the two graded-unitarity
statements to positivity: the package does not assume positivity. The package asserts nothing about the target $T$
(no positivity, nondegeneracy, injectivity or filtered isomorphism on $T$); all such statements are derived.
\end{dictionary}

\begin{physics}
The package is the chiral-algebra version of the statement that gauging a unitary theory produces a unitary theory:
the BRST cohomology of a positive complex, computed by harmonic representatives, inherits positivity. The Kähler
($\mathfrak{sp}_2$) structure is what makes harmonic theory compatible with the extra grading that graded unitarity
needs; it is the content of Beem and Garner's analysis of BRST cosets.
\end{physics}

\begin{status}
Definition (a structure type). It is inhabited at every level datum from $\ltwentynineGU(A)$ alone
(Theorem~\ref{l29:thm:hodge}).
\end{status}

\begin{definition}[The squeeze input]\label{l29:def:squeeze}
For a level datum $L$ and a Kähler geometry $S$, a \emph{squeeze input} is a function
$\mathrm{coeff}:\Z\times\N\to\N$ with (I) for every weight $n$ at which the form $P$ of $S$ is positive definite and every
Kähler degree $j$, the leading cohomology in
bidegree $(j,0)$ has dimension at most $\mathrm{coeff}(n,j/2)$ if $j\ge0$ is even, and $0$ otherwise; and (II) for every
$n$ and $p$, $\sum_{0\le j\le 2p}(\text{that bound})\le\dim F_pT_n$.
\end{definition}

\noindent File \leanpath{VOA/\allowbreak Examples/\allowbreak TwoSeriesAssembly.lean}.
\begin{leancode}
def evenCoeff (a : ℤ → ℕ → ℕ) (n j : ℤ) : ℕ :=
  if 0 ≤ j ∧ j % 2 = 0 then a n (j.toNat / 2) else 0

structure SqueezeInput (S : BRSTKahlerGeometry L) where
  /-- `coeff n r`: the expected dimension of the PBW-length-`r` layer at weight `n`. -/
  coeff : ℤ → ℕ → ℕ
  /-- (I) leading-cohomology / invariant-symbol upper bound. -/
  invariantBound : ∀ n (hp : IsPosOn (S.geometry n).P ⊤) j,
    Module.finrank ℂ ((S.model n hp).leadCohom (j, 0)) ≤ evenCoeff coeff n j
  /-- (II) PBW lower bound for the simple target. -/
  pbwLower : ∀ (n : ℤ) (p : ℕ),
    (∑ j ∈ Finset.Icc 0 (2 * (p : ℤ)), evenCoeff coeff n j) ≤
      Module.finrank ℂ ((virGradedFlag L.c).pieceIn p n)
\end{leancode}

\begin{dictionary}
\lean{Module.finrank ℂ X} is $\dim_\C X$ (and $0$ if $X$ is infinite-dimensional; for (I) this matters only as an
upper bound on a finite-dimensional space, since the relative complexes are finite-dimensional at each weight).
\lean{(virGradedFlag L.c).pieceIn p n} is $F_pT_n$, the length-$\le p$ step of the length flag of $T$ at weight $n$
(Definition~\ref{l19:def:lengthflag}). \lean{Finset.Icc 0 (2 * p)} is the integer interval $[0,2p]$. The Kähler degree
$j$ is twice the PBW length, hence \lean{evenCoeff}. Weights are integers (\lean{ℤ}); negative weights are empty.
\end{dictionary}

\begin{physics}
(I) bounds the number of gauge-invariant ``classical'' states in each $R$-charge sector; (II) counts the states of the
minimal model with at most $p$ stress-tensor insertions. When the two counts agree, every state of the minimal model
is accounted for by a positive harmonic class with the right $R$-charge, which is the physical content of graded
unitarity.
\end{physics}

\begin{status}
Definition. Supplied for every $n$ by Theorems~\ref{l29:thm:invariants} and \ref{l29:thm:hpbw} together with the
invariant count of Section~\ref{l29:sec:model}.
\end{status}

\begin{theorem}[Coset transfer for any level datum]\label{l29:thm:transfer}
Let $L$ be a level datum. If $A=L_{k_A}(\ltwentyninesl)$ is graded unitary, and a Hodge package and a squeeze input
are given, then $L(c,0)$ is graded unitary for its partition-basis real structure. In particular, for the $n$-th
member of the series: $\ltwentynineGU(L_{k_n})\wedge\mathrm{Hodge}\wedge\mathrm{Squeeze}\Rightarrow
\ltwentynineGU(M(2,2n+3))$.
\end{theorem}

\noindent File \leanpath{VOA/\allowbreak Examples/\allowbreak TwoSeriesAssembly.lean}.
\begin{leancode}
theorem cosetLevels_gradedUnitary_of_affine
    (hA : GradedUnitary (affineForm L) (affineRealStructure L) (affineFlag L))
    (hodge : BRSTHodgeInputs L) (hsq : SqueezeInput L hodge.kahler) :
    VirGradedUnitary L.c (targetRealStructure L) := by

theorem twoSeries_gradedUnitary_of_affine (n : ℕ)
    (hA : GradedUnitary (affineForm (twoSeriesLevels n)) (affineRealStructure (twoSeriesLevels n))
      (affineFlag (twoSeriesLevels n)))
    (hodge : BRSTHodgeInputs (twoSeriesLevels n))
    (hsq : SqueezeInput (twoSeriesLevels n) hodge.kahler) :
    VirGradedUnitary (twoSeriesLevels n).c (targetRealStructure (twoSeriesLevels n)) :=
  cosetLevels_gradedUnitary_of_affine (twoSeriesLevels n) hA hodge hsq
\end{leancode}

\begin{dictionary}
Full names \ltwentynineln{VOA.TwoSeries.cosetLevels_gradedUnitary_of_affine} and \ltwentynineln{VOA.TwoSeries.twoSeries_gradedUnitary_of_affine}.
Hypotheses: \lean{hA}, graded unitarity of the actual simple $L_{k_A}$ (Definition~\ref{l29:def:endpoint}); \lean{hodge}
and \lean{hsq}, the two packages above (the second depends on the first through its Kähler geometry). Conclusion:
graded unitarity of the actual simple Virasoro VOA at $c=$ \lean{L.c}. The second theorem is the first at
$L=$ \lean{twoSeriesLevels n}, with $(\lean{twoSeriesLevels n}).c=c_{2,2n+3}$ by \lean{twoSeriesLevels_c}. Same as the
textbook statement, with the two packages as explicit hypotheses. It covers $n=0$ (warning on p.~\pageref{l29:warn:index}).
\end{dictionary}

\begin{proof}[Proof on paper]
From \lean{hodge.weakGU} and Theorem~\ref{l29:thm:aux}, every finite relative complex carries a positive form. Fix a
weight $n$. Let $f$ be the map sending a universal Virasoro state to its BRST class, and $g$ the projection to the simple
quotient $T_n$. The kernel of $f$ lies in the kernel of $g$ (the vacuum class survives), so $T_n$ is a quotient of the
image of $f$. The generator tests put the image of $F_p$ (universal words of length $\le p$) into Kähler degree
$\le 2p$; hence $\dim F_pT_n\le\sum_{j\le 2p}\dim(\text{leading cohomology})_{(j,0)}$. Input (I) bounds the right side by
$\sum_{r\le p}\mathrm{coeff}(n,r)$ and (II) bounds that by $\dim F_pT_n$: all inequalities are equalities (the
\emph{squeeze}). Equality means that the length flag of $T_n$ is carried isomorphically onto the Kähler-degree flag of
harmonic classes, layer by layer, and that odd Kähler degrees carry nothing. The metric bridges identify the canonical
form of $T$ with the harmonic metric on these classes, and the latter is positive with the graded sign pattern. Hence
clauses (iv) and (vi) hold at weight $n$.
\end{proof}

\begin{physics}
The physical slogan: gauge-invariant states of a positive theory are positive, \emph{provided} one can match them one to
one with the states of the coset; the squeeze is the counting argument that does the matching.
\end{physics}

\begin{status}
Proved; \texttt{\#print axioms}-clean (re-checked), \textbf{conditional} on \lean{hA}, \lean{hodge}, \lean{hsq}.
\end{status}

\begin{warning}[title={The short name \texttt{twoSeries\_gradedUnitary\_of\_affine} denotes two different theorems}]\label{l29:warn:clash}
\ltwentynineln{VOA.TwoSeries.twoSeries_gradedUnitary_of_affine} (above, in \leanpath{TwoSeriesAssembly.lean}) takes \lean{hA},
\lean{hodge} and \lean{hsq}. \ltwentynineln{VOA.TwoSeries.SliceMap.twoSeries_gradedUnitary_of_affine} (in
\leanpath{TwoSeriesArakawaUnconditional.lean}, Theorem~\ref{l29:thm:slicegu}) takes \lean{hA} and the jet-ideal
Poincar\'e input \lean{JetPoincareCoreT0 (m + 1)}, and is indexed by $m$ with $n=m+1$. Inside the namespace
\lean{VOA.TwoSeries.SliceMap} the short name resolves to the second; elsewhere in \lean{VOA.TwoSeries}, to the first.
\emph{Neither} is the headline: the headline is \ltwentynineln{VOA.TwoSeries.JetPoincareTwo.twoSeries_gradedUnitary_of_affine_final},
which is declared in \leanpath{TwoSeriesGenIncl.lean}, not in \leanpath{TwoSeriesJetPoincareTwo.lean} as its namespace
suggests. These are naming defects (severity S4: documentation, not statements), not mathematical ones; always
cite the fully qualified name.
\end{warning}

\begin{theorem}[The Hodge package exists at every level]\label{l29:thm:hodge}
For every level datum $L$ and every proof $h_A$ of $\ltwentynineGU(A)$ there is a BRST/Hodge package. Consequently
$\ltwentynineGU(A)\wedge\mathrm{Squeeze}\Rightarrow\ltwentynineGU(L(c,0))$.
\end{theorem}

\noindent File \leanpath{VOA/\allowbreak Examples/\allowbreak TwoSeriesHodgeInputs.lean}.
\begin{leancode}
def brstHodgeInputs
    (hA : GradedUnitary (affineForm L) (affineRealStructure L) (affineFlag L)) :
    BRSTHodgeInputs L where

theorem twoSeries_gradedUnitary_of_affine_hodgeFree (n : ℕ)
    (hA : GradedUnitary (affineForm (twoSeriesLevels n)) (affineRealStructure (twoSeriesLevels n))
      (affineFlag (twoSeriesLevels n)))
    (hsq : SqueezeInput (twoSeriesLevels n) (brstHodgeInputs (twoSeriesLevels n) hA).kahler) :
    VirGradedUnitary (twoSeriesLevels n).c (targetRealStructure (twoSeriesLevels n)) :=
  cosetLevels_gradedUnitary_of_affine_hodgeFree (twoSeriesLevels n) hA hsq
\end{leancode}

\begin{dictionary}
\lean{brstHodgeInputs} is a \lean{def}: a construction of a package, with all seven fields filled by theorems
(\lean{endpointKahler}, \lean{concreteKahler_weakGU}, \lean{sp2GUData0}, \lean{interpretedGenerators_std}, and the three
\lean{harmonic..._proved} bridges, all in the Hodge files). The construction takes \lean{hA} as an argument because the
ghost package uses positivity of $A$; the unitarity of $B=L_1(\ltwentyninesl)$ is the proved
\lean{LevelOneUnitary.bUnitary}. The docstring of \lean{BRSTHodgeInputs} in \leanpath{TwoSeriesAssembly.lean} still says
the package is ``proved so far only at Lee--Yang''; that is stale (warning on p.~\pageref{l29:warn:stale}): this \lean{def}
provides it at every level datum.
\end{dictionary}

\begin{proof}[Proof on paper]
The Lee--Yang proof of Chapter~\ref{ch:L28} used the level only through the identities now recorded as fields of $L$
(non-criticality, reality, the level sum $-4$). The ghost-conjugation identities
(\lean{ghostConjRhoD_proved}, \lean{ghostConjBE_proved}) are identities in the BRST algebra in which the level enters
through the cancellation $k+1+k_U=-4$; the harmonic bridges use only the invariance of the forms under the modes.
Re-running each argument with symbolic levels gives the package.
\end{proof}

\begin{physics}
That the BRST--Hodge mechanism is level-independent is expected: the ghost system and the anomaly cancellation do not
care which $(A_1,D_{2n+1})$ theory is being gauged.
\end{physics}

\begin{status}
Proved (construction), \texttt{\#print axioms}-clean (re-checked); the endpoint theorem is \textbf{conditional} on
\lean{hA} and \lean{hsq}.
\end{status}

\subsection{The Koszul step: from cohomology to an affine invariant count}

\begin{definition}[The affine invariant spaces]\label{l29:def:invariants}
For $r\in\N$ and $N\in\Z$, $\ltwentynineinv_{r,N}$ is the space of tensors in $\ltwentyninegr_rA\otimes B$ of total
conformal weight $N$ annihilated by every diagonal non-negative current mode $J^a_m\otimes1+1\otimes J^a_m$
($a\in\{e,f,h\}$, $m\ge0$). Here $\ltwentyninegr_rA$ is the $r$-th layer of the corrected flag of $A$ with its induced
mode action, and $B=L_1(\ltwentyninesl)$.
\end{definition}

\noindent File \leanpath{VOA/\allowbreak Examples/\allowbreak TwoSeriesKoszulInvariants.lean}.
\begin{leancode}
def affineInvariants (r : ℕ) (N : ℤ) : Submodule ℂ (GrA L r ⊗[ℂ] BSpace) :=
  VOA.CyclicEval.tensorWeight (grDeg L r) VOA.CyclicEval.bWeight N ⊓ jointKer (grDiag L r)
\end{leancode}

\begin{dictionary}
\lean{GrA L r} is $\ltwentyninegr_rA$ (file \leanpath{TwoSeriesKoszulGr.lean}); \lean{BSpace} is the space of
$L_1(\ltwentyninesl)$; \lean{⊗[ℂ]} the tensor product over $\C$. \lean{tensorWeight ... N} is the weight-$N$ part for the
sum of the two gradings, \lean{jointKer (grDiag L r)} the common kernel of the diagonal non-negative modes, and
\lean{⊓} the intersection of submodules. The companion lemma \lean{mem_affineInvariants} restates membership as
``weight $N$ and \lean{IsInv}'' (killed by all diagonal modes). Note that the definition involves only the affine
algebra: neither the BRST complex nor $T$.
\end{dictionary}

\begin{physics}
$\ltwentynineinv_{r,N}$ counts the classical (leading-order) gauge-invariant states of the gauged system at
$R$-charge $r$ and weight $N$: invariance under the non-negative modes of the diagonal current algebra is the
classical form of the BRST condition.
\end{physics}

\begin{status}
Definition.
\end{status}

\begin{theorem}[The leading Koszul embedding]\label{l29:thm:koszul}
For every level datum $L$, every $h_A:\ltwentynineGU(A)$ and every weight $n$ at which the modified Hermitian form $P$ of the
Kähler geometry is positive definite (hypothesis $h_p$), the leading cohomology of the relative complex at weight $n$ and Kähler
degree $(2r,0)$ has dimension at most $\dim\ltwentynineinv_{r,n}$. (Under $h_A$ the form is positive definite at every weight, so in the
application $h_p$ is not an extra assumption.)
\end{theorem}

\noindent File \leanpath{VOA/\allowbreak Examples/\allowbreak TwoSeriesKoszulInvariants.lean}.
\begin{leancode}
theorem rank_leadCohom_le_affineInvariants (n : ℤ)
    (hp : IsPosOn ((endpointKahler L hA).geometry n).P ⊤) (r : ℕ) :
    Module.rank ℂ (((endpointKahler L hA).model n hp).leadCohom (2 * (r : ℤ), 0)) ≤
      Module.rank ℂ (affineInvariants L r n) := by
\end{leancode}

\begin{dictionary}
\lean{Module.rank} is the dimension as a cardinal (it can be infinite), so the inequality is meaningful even if
$\ltwentynineinv_{r,n}$ were infinite-dimensional. \lean{endpointKahler L hA} is the Kähler geometry of
Theorem~\ref{l29:thm:hodge}; \lean{hp} its positivity at weight $n$. \lean{hp} is a genuine explicit hypothesis of the Lean
statement (the Kähler complex at weight $n$ is defined only once it is given); in every use it is discharged from \lean{hA}
and \lean{auxiliaryGradedUnitary} through the \lean{weakGU} field of the Hodge package (Theorem~\ref{l29:thm:hodge}). Here
\lean{L} and \lean{hA} are section variables, explicit arguments in the order \lean{(L) (hA)}. The proof is an injective linear map
(\lean{leadKoszulEmbedding}) from the leading cohomology into the invariant space. Same as the textbook statement.
\end{dictionary}

\begin{proof}[Proof on paper]
This is the Koszul argument of Chapter~\ref{ch:L28}, re-run at level $L$. Filter the relative complex by the corrected
PBW degree. The leading part of the BRST differential is a Koszul
differential for the action of the non-negative diagonal modes on $\ltwentyninegr A\otimes B\otimes\ltwentyninegr U$,
twisted by ghosts. The universal factor $\ltwentyninegr U$ is free, so the Koszul complex is exact away from ghost
number zero and the ghost-number-zero leading cohomology injects into the invariants of $\ltwentyninegr A\otimes B$; the
symbol maps $\pi_0,\pi_1$ make this explicit. Odd and negative Kähler degrees have no leading cohomology
(\lean{endpoint_oddZero}, \lean{endpoint_negZero}).
\end{proof}

\begin{physics}
The Koszul step is the statement that, at leading order, BRST cohomology is just the space of gauge invariants: the
ghosts and the auxiliary sector cancel the gauge directions.
\end{physics}

\begin{status}
Proved, conditional on \lean{hA} (which enters only through the Kähler geometry) and on the weightwise positivity \lean{hp}, which
follows from \lean{hA}; \texttt{\#print axioms}-clean (re-checked).
\end{status}

\begin{theorem}[Transfer with an invariant count]\label{l29:thm:invariants}
For every $n$: if $\ltwentynineGU(L_{k_n})$ holds and a function $\mathrm{coeff}$ satisfies
$\dim\ltwentynineinv_{r,N}\le\mathrm{coeff}(N,r)$ for all $N,r$ and
$\sum_{r\le p}\mathrm{coeff}(N,r)\le\dim F_pT_N$ for all $N,p$, then $M(2,2n+3)$ is graded unitary.
\end{theorem}

\noindent File \leanpath{VOA/\allowbreak Examples/\allowbreak TwoSeriesKoszulInvariants.lean}.
\begin{leancode}
theorem twoSeries_gradedUnitary_of_affine_invariants (n : ℕ)
    (hA : GradedUnitary (affineForm (twoSeriesLevels n)) (affineRealStructure (twoSeriesLevels n))
      (affineFlag (twoSeriesLevels n)))
    (coeff : ℤ → ℕ → ℕ)
    (hinv : ∀ (N : ℤ) (r : ℕ),
      Module.rank ℂ (affineInvariants (twoSeriesLevels n) r N) ≤ coeff N r)
    (hpbw : ∀ (N : ℤ) (p : ℕ), ∑ r ∈ Finset.range (p + 1), coeff N r ≤
      Module.finrank ℂ ((virGradedFlag (twoSeriesLevels n).c).pieceIn p N)) :
    VirGradedUnitary (twoSeriesLevels n).c (targetRealStructure (twoSeriesLevels n)) :=
\end{leancode}

\begin{dictionary}
\lean{Finset.range (p + 1)} is $\{0,\dots,p\}$. \lean{hinv} is stated with \lean{Module.rank}, so an
infinite-dimensional invariant space makes \lean{hinv} \emph{false} (a cardinal $\ge\aleph_0$ is not $\le$ a natural
number), never vacuously true; a \lean{_finrank} variant takes \lean{finrank} plus finite-dimensionality. The
proof assembles \lean{squeezeInput_of_invariants} (Theorem~\ref{l29:thm:koszul} plus the parity vanishing) with
Theorem~\ref{l29:thm:hodge}. Both remaining inputs are about separate objects: \lean{hinv} about the affine algebra,
\lean{hpbw} about $T$.
\end{dictionary}

\begin{proof}[Proof on paper]
Set the squeeze coefficient to $\mathrm{coeff}$. Half (I): by Theorem~\ref{l29:thm:koszul} and \lean{hinv} for even
Kähler degree, and by the parity vanishing for odd or negative degree. Half (II) is \lean{hpbw} after rewriting the sum
over Kähler degrees $j\le2p$ as a sum over lengths $r\le p$ (\lean{sum_evenCoeff}). Apply Theorem~\ref{l29:thm:hodge}.
\end{proof}

\begin{physics}
After this step, the transfer needs one count on each side of the coset: the classical gauge invariants of the
$(A_1,D_{2n+1})$ side, and the states of the $(A_1,A_{2n})$ side.
\end{physics}

\begin{status}
Proved, \texttt{\#print axioms}-clean (re-checked), \textbf{conditional} on \lean{hA}, \lean{hinv}, \lean{hpbw}.
\end{status}

%======================================================================
\section{The vacuum characters of \texorpdfstring{$M(2,2n+3)$}{M(2,2n+3)}}\label{l29:sec:char}
%======================================================================

This section supplies the input \lean{hpbw} of Theorem~\ref{l29:thm:invariants} with the Andrews--Gordon count, for
every $n$, and in passing proves the vacuum character of every $M(2,2n+3)$. All of it is about the Virasoro side
only.

\subsection{The Andrews--Gordon count}

A Virasoro creation word $L_{-j_1-2}\cdots L_{-j_\ell-2}\ltwentyninevac$ is encoded by the list $(j_1,\dots,j_\ell)$
of natural numbers; its \emph{energy} (conformal weight) is $\sum_i(j_i+2)$ (\lean{energyVir}), and it is
\emph{sorted} if $j_1\ge j_2\ge\cdots$ (\lean{SortedWord}). The parts of the corresponding partition are $j_i+2\ge2$.

\begin{definition}[Gordon-good words; the Andrews--Gordon count]\label{l29:def:gordon}
A list $w=(w_0,\dots,w_{\ell-1})$ is \emph{Gordon-good of order $n$} if $w_i\ge w_{i+n}+2$ whenever both entries
exist. The \emph{Andrews--Gordon count} $\ltwentyninegc_n(N,r)$ is the number of sorted Gordon-good words of length
exactly $r$ and energy $N$ (and $0$ for $N<0$). Equivalently, it is the number of partitions of $N$ into exactly $r$
parts $\ge2$ in which, writing $f_j$ for the multiplicity of the part $j$, $f_j+f_{j+1}\le n$ for every $j$.
\end{definition}

\noindent Files \leanpath{VOA/\allowbreak Examples/\allowbreak TwoSeriesCharacterSpanning.lean} (Gordon-good) and
\leanpath{VOA/\allowbreak Examples/\allowbreak TwoSeriesCharacterReduction.lean} (the count).
\begin{leancode}
def GordonGood (n : ℕ) (w : List ℕ) : Prop :=
  ∀ i < w.length - n, w.getD (i + n) 0 + 2 ≤ w.getD i 0

def tupleList {N r : ℕ} (a : Fin r → Fin (N + 1)) : List ℕ := List.ofFn fun j => (a j : ℕ)

def gordonWords (n N r : ℕ) : Finset (List ℕ) :=
  ((Finset.univ : Finset (Fin r → Fin (N + 1))).image tupleList).filter
    (fun u => SortedWord u ∧ GordonGood n u ∧ energyVir u = N)

def gordonCoeff (n : ℕ) (N : ℤ) (r : ℕ) : ℕ :=
  if 0 ≤ N then (gordonWords n N.toNat r).card else 0
\end{leancode}

\begin{dictionary}
\lean{w.getD i 0} is the $i$-th entry of $w$ (or $0$ if $i$ is out of range; it is never out of range here because
$i<\ell-n$). \lean{gordonWords n N r} enumerates all functions $\{0,\dots,r-1\}\to\{0,\dots,N\}$ (\lean{Fin r → Fin (N + 1)}),
turns each into a list, and keeps the sorted, Gordon-good lists of energy $N$; since every entry of a word of energy $N$
is $\le N$, nothing is lost by the bound. \lean{.card} is the number of elements of the finite set. So
\lean{gordonCoeff n N r} is exactly $\ltwentyninegc_n(N,r)$; \lean{mem_gordonWords} proves the membership
characterisation (sorted, Gordon-good, energy $N$, length $r$).

\textbf{Why the two descriptions agree.} For a sorted list, $w_i\ge w_{i+n}+2$ says that no window of $n+1$
consecutive entries spans fewer than three consecutive values; equivalently no two adjacent values $j,j+1$ occur more
than $n$ times in total, i.e.\ $f_j+f_{j+1}\le n$. At $n=1$ it is the Rogers--Ramanujan condition ``parts differ by at
least $2$''.
\end{dictionary}

\begin{example}[The counts for $n=1,2,3$]\label{l29:ex:gordon}
Summing over $r$, the first coefficients of $\sum_N\bigl(\sum_r\ltwentyninegc_n(N,r)\bigr)q^N$ are (computed by
enumeration, and agreeing with the alternating sums of Definition~\ref{l29:def:rcq}):
\begin{center}\small
\begin{tabular}{c|cccccccccccccc}
$N$ & 0&1&2&3&4&5&6&7&8&9&10&11&12&13\\ \hline
$n=1$ (Lee--Yang) & 1&0&1&1&1&1&2&2&3&3&4&4&6&6\\
$n=2$ ($M(2,7)$) & 1&0&1&1&2&2&3&3&5&6&8&9&13&14\\
$n=3$ ($M(2,9)$) & 1&0&1&1&2&2&4&4&6&7&10&12&17&19\\
\end{tabular}
\end{center}
By length, at $n=2$: $N=8$ has $\ltwentyninegc_2(8,1)=1$ (the word $(6)$, i.e.\ $L_{-8}\ltwentyninevac$),
$\ltwentyninegc_2(8,2)=3$ (parts $\{6,2\},\{5,3\},\{4,4\}$) and $\ltwentyninegc_2(8,3)=1$ (parts $\{4,2,2\}$; the
partition $\{3,3,2\}$ is excluded because $f_2+f_3=3>2$). At $n=1$ the partition $\{4,4\}$ is also excluded
($f_4=2>1$), and indeed $\ltwentyninegc_1(8,2)=2$.
\end{example}

\begin{physics}
The partitions with $f_j+f_{j+1}\le n$ are the ``Gordon'' or ``$(n+1)$-exclusion'' configurations: stress-tensor
quanta behave like particles of which at most $n$ may occupy two adjacent momentum levels. This is the quasi-particle
(fermionic) description of the $(A_1,A_{2n})$ chiral algebra, and for $n=1$ the Lee--Yang ``hard-hexagon'' exclusion.
\end{physics}

\begin{status}
Definitions.
\end{status}

\subsection{Spanning: the upper bound}

\begin{definition}[Jet relations]\label{l29:def:jetrel}
For $c\in\C$ and $n\in\N$, $V(c,0)$ \emph{has jet relations of order $n+1$} if there is $\alpha\ne0$ such that for every
$k\in\N$ some vector of the maximal ideal of $V(c,0)$ agrees with $\alpha\sum_{e_1+\cdots+e_{n+1}=k}
L_{-e_1-2}\cdots L_{-e_{n+1}-2}\ltwentyninevac$ modulo words of length $\le n$ and weight $2n+2+k$.
\end{definition}

\noindent File \leanpath{VOA/\allowbreak Examples/\allowbreak TwoSeriesCharacterLeading.lean}.
\begin{leancode}
def jetSum : ℕ → ℕ → VirUniv c
  | 0, k => if k = 0 then vacVec c else 0
  | m + 1, k => ∑ p ∈ Finset.antidiagonal k, Lop c (-(p.1 : ℤ) - 2) (jetSum m p.2)

def HasJetRelations (c : ℂ) (n : ℕ) : Prop :=
  ∃ α : ℂ, α ≠ 0 ∧ ∀ k : ℕ, ∃ σ ∈ maxIdealVir c,
    σ - α • jetSum c (n + 1) k ∈ virFil c n ((2 * n + 2 + k : ℕ) : ℤ)
\end{leancode}

\begin{dictionary}
\lean{VirUniv c} is the universal vacuum module $V(c,0)$; \lean{Lop c j} is the mode $L_j$; \lean{vacVec c} the vacuum.
\lean{Finset.antidiagonal k} is the set of pairs $(a,b)$ with $a+b=k$, so \lean{jetSum c m k} is the sum over
compositions $e_1+\cdots+e_m=k$ of $L_{-e_1-2}\cdots L_{-e_m-2}\ltwentyninevac$: the weight-$(2m+k)$ coefficient of the
normally ordered power $(T^m)$, i.e.\ the $k$-th jet of $T^m$. \lean{maxIdealVir c} is the maximal ideal, and
\lean{virFil c n N} the span of words of length $\le n$ and weight $N$. So the definition is the textbook one.
\end{dictionary}

\begin{physics}
The minimal model $M(2,2n+3)$ has a null vector at weight $2n+2$ whose leading term is $T^{n+1}$ (the normally ordered
$(n+1)$-st power of the stress tensor): ``$T^{n+1}=0$ up to derivatives''. Its jets are the relations $\partial^k(T^{n+1})$.
\end{physics}

\begin{status}
Definition.
\end{status}

\begin{theorem}[Jet relations of the series]\label{l29:thm:jetrel}
For every $n\in\N$, $V(c_{2,2n+3},0)$ has jet relations of order $n+1$.
\end{theorem}

\noindent File \leanpath{VOA/\allowbreak Examples/\allowbreak TwoSeriesCharacterEndpoint.lean}.
\begin{leancode}
theorem twoSeries_hasJetRelations_all (n : ℕ) :
    VOA.TwoSeriesCharacter.HasJetRelations (twoSeriesLevels n).c n := by
\end{leancode}

\begin{dictionary}
For $n\ge1$ the relation comes from the vacuum singular vector of Benoit--Saint-Aubin type at level $2n+2$
(\lean{bsaSource_mem_maxIdeal}), whose length-$(n+1)$ part is a nonzero multiple of $L_{-2}^{n+1}\ltwentyninevac$
(\lean{bsaLead_ne_zero}); its divided $L_{-1}$-powers give all jets (\lean{hasJetRelations_of_singular}). For $n=0$,
$c=0$ and $L_{-2}\ltwentyninevac$ itself lies in the maximal ideal. The charge is matched by
\ltwentynineln{twoSeriesLevels_c_eq_bsaCharge}: $c_{2,2n+3}=$ \lean{bsaCharge (2n+2)}. No hypothesis.
\end{dictionary}

\begin{proof}[Proof on paper]
At $c=c_{2,2n+3}$ the vacuum Verma module has a singular vector at level $(2-1)(2n+3-1)=2n+2$ (Kac determinant). Its
expansion in PBW words has leading length-$(n+1)$ term $\alpha L_{-2}^{n+1}\ltwentyninevac$, $\alpha\ne0$ (an explicit
Benoit--Saint-Aubin formula, checked in the library). Applying $L_{-1}^k/k!$ and using $[L_{-1},L_{-j}]=(j-1)L_{-j-1}$
produces, at length $n+1$, exactly $\alpha$ times the $k$-th jet sum, and the maximal ideal is $L_{-1}$-stable.
\end{proof}

\begin{physics}
The null vector $\chi_{2n+2}$ of $M(2,2n+3)$ and its descendants; the leading term says the classical limit of the
model is the arc space of $\C[T]/(T^{n+1})$.
\end{physics}

\begin{status}
Proved, unconditional. Axioms standard (re-checked).
\end{status}

\begin{theorem}[Gordon spanning: the refined upper bound]\label{l29:thm:upper}
If $V(c,0)$ has jet relations of order $n+1$, then for all $N\in\Z$ and $p\in\N$,
$\dim F_pL(c,0)_N\le\sum_{r\le p}\ltwentyninegc_n(N,r)$.
\end{theorem}

\noindent File \leanpath{VOA/\allowbreak Examples/\allowbreak TwoSeriesCharacterReduction.lean}.
\begin{leancode}
theorem finrank_pieceIn_le_gordon (hJ : HasJetRelations c n) (N : ℤ) (p : ℕ) :
    Module.finrank ℂ ((virGradedFlag c).pieceIn p N) ≤
      ∑ r ∈ Finset.range (p + 1), gordonCoeff n N r := by
\end{leancode}

\begin{dictionary}
\lean{(virGradedFlag c).pieceIn p N} is $F_pL(c,0)_N$. The proof rests on \lean{goodFil_eq_gordonSpan}: under the jet
relations, $F_pL(c,0)_N$ is the image of the span of sorted Gordon-good words of length $\le p$ and energy $N$. Then
the dimension is at most the number of such words, which is \lean{card_gordonWordsLe} $=\sum_{r\le p}\ltwentyninegc_n(N,r)$.
\end{dictionary}

\begin{proof}[Proof on paper]
Straightening: a sorted word that is not Gordon-good contains a window $w_i,\dots,w_{i+n}$ with $w_i\le w_{i+n}+1$.
Among all sorted words with the same length and energy, the quantity $\sum w_i^2$ is minimised by the ``flat''
windows, and the jet relation of the appropriate $k$ expresses the flattest window's product of $n+1$ modes as minus
the sum of the other compositions (all with larger $\sum w_i^2$ after sorting) plus shorter words plus the maximal
ideal. Induction on $\sum w_i^2$ (\lean{sqSum_lt_of_not_perm}, \lean{defect}) rewrites every sorted word into
Gordon-good ones of the same energy and no larger length, modulo the maximal ideal.
\end{proof}

\begin{physics}
The null vector, used as a rewriting rule, removes all states violating the exclusion rule: at most $n$ quanta of
$T$ in two adjacent momentum levels.
\end{physics}

\begin{status}
Proved, \texttt{\#print axioms}-clean (re-checked), conditional on \lean{hJ} (supplied for the series by
Theorem~\ref{l29:thm:jetrel}).
\end{status}

\subsection{The lower bound: Rocha-Caridi and Andrews--Gordon}

\begin{definition}[The total lower bound; its refined form]\label{l29:def:gordonlower}
$\mathrm{GordonLowerBound}(c,n)$: for every $N\in\N$, $\sum_{r\le N}\ltwentyninegc_n(N,r)\le\dim L(c,0)_N$.
$\mathrm{GordonPBWLower}(c,n)$: for every $N\in\Z$, $p\in\N$, $\sum_{r\le p}\ltwentyninegc_n(N,r)\le\dim F_pL(c,0)_N$.
\end{definition}

\noindent File \leanpath{VOA/\allowbreak Examples/\allowbreak TwoSeriesCharacterReduction.lean}.
\begin{leancode}
def GordonLowerBound (c : ℂ) (n : ℕ) : Prop :=
  ∀ N : ℕ, (∑ r ∈ Finset.range (N + 1), gordonCoeff n (N : ℤ) r) ≤
    Module.finrank ℂ ((virasoroSimple c).grading (N : ℤ))

def GordonPBWLower (c : ℂ) (n : ℕ) : Prop :=
  ∀ (N : ℤ) (p : ℕ), ∑ r ∈ Finset.range (p + 1), gordonCoeff n N r ≤
    Module.finrank ℂ ((virGradedFlag c).pieceIn p N)
\end{leancode}

\begin{dictionary}
\lean{(virasoroSimple c).grading N} is the weight-$N$ space $L(c,0)_N$ of the simple Virasoro VOA. The range
$r\le N$ in \lean{GordonLowerBound} loses nothing: a word of energy $N$ has length $\le N/2$. \lean{GordonPBWLower} is
the shape of the input \lean{hpbw} of Theorem~\ref{l29:thm:invariants} with $\mathrm{coeff}=\ltwentyninegc_n$.
\end{dictionary}

\begin{physics}
The total lower bound says that the minimal model has \emph{at least} as many states as the exclusion statistics
allows; together with spanning it says it has exactly that many.
\end{physics}

\begin{status}
Definitions (predicates). Both are proved for the series (Theorems~\ref{l29:thm:character} and
\ref{l29:thm:hpbw}). The docstring of \lean{GordonLowerBound} still reads ``OPEN for general $n$ \dots\ not formalised
here''; that is stale (warning on p.~\pageref{l29:warn:stale}).
\end{status}

\begin{theorem}[The refined lower bound is equivalent to the total one]\label{l29:thm:hpbwiff}
If $V(c,0)$ has jet relations of order $n+1$, then $\mathrm{GordonPBWLower}(c,n)\iff\mathrm{GordonLowerBound}(c,n)$.
\end{theorem}

\noindent File \leanpath{VOA/\allowbreak Examples/\allowbreak TwoSeriesCharacterReduction.lean}.
\begin{leancode}
theorem hpbw_iff_gordonLowerBound (hJ : HasJetRelations c n) :
    GordonPBWLower c n ↔ GordonLowerBound c n :=
  ⟨gordonLowerBound_of_hpbw c n, hpbw_of_gordonLowerBound c n hJ⟩
\end{leancode}

\begin{dictionary}
\lean{↔} is logical equivalence; \lean{⟨a, b⟩} pairs the two directions. Only the direction ``total $\Rightarrow$
refined'' uses \lean{hJ}.
\end{dictionary}

\begin{proof}[Proof on paper]
($\Rightarrow$) Take $p=N$: $F_NL(c,0)_N=L(c,0)_N$. ($\Leftarrow$) By Gordon spanning (Theorem~\ref{l29:thm:upper})
the Gordon-good words of length $\le N$ span $L(c,0)_N$; there are $\sum_r\ltwentyninegc_n(N,r)$ of them, and by the
total lower bound this number is at most $\dim L(c,0)_N$. So they form a basis, and in particular the subfamily of
length $\le p$ is linearly independent inside $F_pL(c,0)_N$, which gives the refined bound. This is a squeeze in the
same spirit as Theorem~\ref{l29:thm:transfer}.
\end{proof}

\begin{physics}
Once the total number of states is right, the exclusion-statistics basis is automatically adapted to the number of
stress-tensor quanta.
\end{physics}

\begin{status}
Proved, \texttt{\#print axioms}-clean (re-checked), conditional on \lean{hJ}.
\end{status}

\begin{definition}[The Rocha-Caridi coefficient]\label{l29:def:rcq}
With $p(N)$ the number of partitions of $N$ ($p(N)=0$ for $N<0$) and $e_n(k)=k\bigl((2n+3)k+2n+1\bigr)/2$,
\begin{equation*}
\mathrm{rcq}_n(N)\;=\;\sum_{k\in\Z}(-1)^k\,p\bigl(N-e_n(k)\bigr),
\end{equation*}
a finite sum (the terms with $|k|>N$ vanish).
\end{definition}

\noindent File \leanpath{VOA/\allowbreak Examples/\allowbreak TwoSeriesCharacterKacCount.lean}.
\begin{leancode}
def eeq (k : ℤ) : ℤ := k * ((2 * n + 3) * k + (2 * n + 1)) / 2

def rcq (N : ℕ) : ℤ := ∑ k ∈ Ico (-(N : ℤ) - 1) ((N : ℤ) + 1), VOA.RRLower.sgn k * VOA.RRLower.pp ((N : ℤ) - eeq n k)
\end{leancode}

\begin{dictionary}
Integer division \lean{/ 2} is exact here (\lean{two_mul_eeq}: the numerator $k((2n+3)k+2n+1)$ is always even).
\lean{Ico a b} is $[a,b)$, so $k$ runs over $-N-1\le k\le N$, which contains every $k$ with $e_n(k)\le N$
(\lean{abs_le_eeq}). \lean{VOA.RRLower.sgn k} is $(-1)^k$ and \lean{VOA.RRLower.pp} the partition function extended by
$0$ to negative arguments (both from the Lee--Yang character files of Chapter~\ref{ch:L28}).
\end{dictionary}

\begin{physics}
This is the Rocha-Caridi (Feigin--Fuchs) formula for the vacuum character of $M(2,2n+3)$:
$\chi_0=\frac{1}{(q)_\infty}\sum_k(-1)^kq^{e_n(k)}$, the alternating sum over the embedding diagram of Verma
submodules, here at $p=2$, $p'=2n+3$.
\end{physics}

\begin{status}
Definition.
\end{status}

\begin{theorem}[Rocha-Caridi lower bound]\label{l29:thm:rcq}
For every $n,N\in\N$, $\mathrm{rcq}_n(N)\le\dim L(c_{2,2n+3},0)_N$.
\end{theorem}

\noindent File \leanpath{VOA/\allowbreak Examples/\allowbreak TwoSeriesCharacterGramRank.lean}.
\begin{leancode}
theorem rcq_le_finrank (N : ℕ) :
    rcq n N ≤ (Module.finrank ℂ ((VOA.VirSimple.virasoroSimple
      (VOA.VirModule.bsaCharge (2 * n + 2))).grading (N : ℤ)) : ℤ) := by
\end{leancode}

\begin{dictionary}
\lean{bsaCharge (2n+2)} is $c_{2,2n+3}$ (Theorem~\ref{l29:thm:jetrel}). The inequality is in $\Z$ (the cast \lean{: ℤ}).
Only the lower bound is claimed here; the matching upper bound is Theorem~\ref{l29:thm:upper}.
\end{dictionary}

\begin{proof}[Proof on paper]
$\dim L(c,0)_N$ is the rank of the Shapovalov Gram matrix of $V(c,0)$ at level $N$. Its determinant (Kac) is a product
of factors $(h-h_{r,s}(c))^{p(N-rs)}$. Deforming $c$ away from $c_{2,2n+3}$ along a line and counting the orders of
vanishing of the Kac factors at the special point (\lean{rcq_eq_orders}, \lean{gram_shift_det_q}) bounds the corank by
$\sum_{k\ne0}$ of the partition numbers with alternating signs; the rank is therefore at least $\mathrm{rcq}_n(N)$.
\end{proof}

\begin{physics}
The rank of the Gram matrix is the number of physical (non-null) states; the Kac determinant counts the nulls.
\end{physics}

\begin{status}
Proved, unconditional. \texttt{\#print axioms}-clean as part of Theorem~\ref{l29:thm:character} (re-checked there).
\end{status}

\begin{theorem}[The Andrews--Gordon identity, coefficientwise]\label{l29:thm:ag}
For all $n,N\in\N$, $\sum_{r\le N}\ltwentyninegc_n(N,r)=\mathrm{rcq}_n(N)$. Equivalently,
\begin{equation*}
\sum_{\substack{\lambda\vdash N,\ \text{parts}\ge2\\ f_j+f_{j+1}\le n}}1\;=\;\sum_{k\in\Z}(-1)^k\,p\bigl(N-e_n(k)\bigr).
\end{equation*}
\end{theorem}

\noindent File \leanpath{VOA/\allowbreak Examples/\allowbreak TwoSeriesCharacterAGIdentity.lean}.
\begin{leancode}
theorem gordon_sum_eq_rcq (N : ℕ) :
    ((∑ r ∈ Finset.range (N + 1), gordonCoeff n (N : ℤ) r : ℕ) : ℤ) = KacGram.rcq n N := by
\end{leancode}

\begin{dictionary}
Full name \ltwentynineln{VOA.TwoSeriesCharacter.AG.gordon_sum_eq_rcq}. The left side is cast from \lean{ℕ} to \lean{ℤ}. This is
the identity for one coefficient $N$ at a time, for every $N$ and $n$; it is equivalent to the identity of power
series $\sum_N(\cdots)q^N=\prod_{j\not\equiv0,\pm1\ (\mathrm{mod}\ 2n+3)}(1-q^j)^{-1}$ by Jacobi's triple product,
but the library states and uses only the coefficientwise form against the alternating sum.
\end{dictionary}

\begin{proof}[Proof on paper]
Andrews' proof. Let $G_i(x;q)=\sum_{m,N}\#\{\text{partitions with }m\text{ parts},\ \text{weight }N,\ f_j+f_{j+1}\le n,
\ f_1\le i-1\}\,x^mq^N$ (the shift by one between ``parts $\ge1$ with $f_1\le i-1$'' and ``parts $\ge2$'' is the
bookkeeping of \lean{gC}). Removing the smallest parts gives Andrews' recursion
$G_{a+1}-G_a=x^aq^aG_{n+1-a}(xq)$ (\lean{GG_rec}). On the other side, Andrews' explicit series
$H_i=F\cdot A_i$, with $A_i$ an alternating sum of the Rocha-Caridi shape and $F=1/(xq)_\infty$, satisfies the same
recursion and boundary values (\lean{HH_rec}). The recursion determines the series uniquely (induction on the
$q$-degree, \lean{GG_eq_HH}). Setting $x=1$, $i=1$ and reading off the $q^N$ coefficient gives the identity; the
evaluation at $x=1$ is done on a ring of ``admissible'' two-variable series where it is a ring homomorphism
(\lean{Adm}, \lean{ev}).
\end{proof}

\begin{physics}
The Andrews--Gordon identity is the statement that the quasi-particle (fermionic) and the Feigin--Fuchs (bosonic)
descriptions of the $(A_1,A_{2n})$ vacuum character agree. For $n=1$ it is the first Rogers--Ramanujan identity.
\end{physics}

\begin{status}
Proved, unconditional. Axioms standard (re-checked).
\end{status}

\begin{theorem}[The vacuum character of $M(2,2n+3)$]\label{l29:thm:character}
For all $n,N\in\N$,
\begin{equation*}
\dim L(c_{2,2n+3},0)_N\;=\;\sum_{r\le N}\ltwentyninegc_n(N,r)\;=\;\mathrm{rcq}_n(N).
\end{equation*}
In particular the total lower bound $\mathrm{GordonLowerBound}(c_{2,2n+3},n)$ holds.
\end{theorem}

\noindent File \leanpath{VOA/\allowbreak Examples/\allowbreak TwoSeriesCharacterAGIdentity.lean}.
\begin{leancode}
theorem gordonLowerBound_twoSeries (n : ℕ) :
    VOA.TwoSeriesCharacter.GordonLowerBound (twoSeriesLevels n).c n := by

theorem twoSeries_character (n N : ℕ) :
    Module.finrank ℂ ((VOA.VirSimple.virasoroSimple (twoSeriesLevels n).c).grading (N : ℤ)) =
      ∑ r ∈ Finset.range (N + 1), VOA.TwoSeriesCharacter.gordonCoeff n (N : ℤ) r ∧
    ((∑ r ∈ Finset.range (N + 1), VOA.TwoSeriesCharacter.gordonCoeff n (N : ℤ) r : ℕ) : ℤ) =
      VOA.TwoSeriesCharacter.KacGram.rcq n N :=
\end{leancode}

\begin{dictionary}
\lean{∧} is ``and''. The first conjunct is an equality of natural numbers, the second of integers. Both are about the
\emph{actual simple} Virasoro VOA at $c=c_{2,2n+3}$, for every $n$ including $n=0$ (where $c=0$, every count vanishes
for $N>0$ and $L(0,0)=\C$). No hypothesis.
\end{dictionary}

\begin{proof}[Proof on paper]
$\le$: Gordon spanning (Theorem~\ref{l29:thm:upper}) with $p=N$ and the jet relations (Theorem~\ref{l29:thm:jetrel}).
$\ge$: the Andrews--Gordon identity (Theorem~\ref{l29:thm:ag}) turns the count into $\mathrm{rcq}_n(N)$, which is at
most $\dim L(c,0)_N$ by Theorem~\ref{l29:thm:rcq}. The second conjunct is Theorem~\ref{l29:thm:ag}.
\end{proof}

\begin{physics}
The Schur index of the $(A_1,A_{2n})$ Argyres--Douglas theory, $\prod_{j\not\equiv0,\pm1\,(2n+3)}(1-q^j)^{-1}$, is the
vacuum character of $M(2,2n+3)$: this theorem proves its coefficients one by one from the definition of the simple
Virasoro VOA.
\end{physics}

\begin{status}
Proved, \textbf{unconditional}. Axioms standard (re-checked, both declarations).
\end{status}

\begin{theorem}[The refined PBW character of $M(2,2n+3)$]\label{l29:thm:hpbw}
For every $n\in\N$, $N\in\Z$, $p\in\N$: $\sum_{r\le p}\ltwentyninegc_n(N,r)\le\dim F_pL(c_{2,2n+3},0)_N$. Combined with
Theorem~\ref{l29:thm:upper}, the inequality is an equality.
\end{theorem}

\noindent File \leanpath{VOA/\allowbreak Examples/\allowbreak TwoSeriesCharacterAGIdentity.lean}.
\begin{leancode}
theorem twoSeries_hpbw (n : ℕ) :
    ∀ (N : ℤ) (p : ℕ), ∑ r ∈ Finset.range (p + 1), VOA.TwoSeriesCharacter.gordonCoeff n N r ≤
      Module.finrank ℂ ((virGradedFlag (twoSeriesLevels n).c).pieceIn p N) :=
\end{leancode}

\begin{dictionary}
Exactly the hypothesis \lean{hpbw} of Theorem~\ref{l29:thm:invariants} at $\mathrm{coeff}=\ltwentyninegc_n$
(\lean{gordonCoeff n}, read as a function of $(N,r)$). Proved from Theorems~\ref{l29:thm:hpbwiff},
\ref{l29:thm:jetrel} and \ref{l29:thm:character}.
\end{dictionary}

\begin{physics}
The number of states of $M(2,2n+3)$ with at most $p$ stress-tensor insertions is exactly the number of exclusion
configurations with at most $p$ quanta.
\end{physics}

\begin{status}
Proved, unconditional. Axioms standard (re-checked).
\end{status}

\begin{theorem}[Transfer with the PBW input discharged]\label{l29:thm:gordonfinal}
For every $n$: if $L_{k_n}(\ltwentyninesl)$ is graded unitary and $\dim\ltwentynineinv_{r,N}\le\ltwentyninegc_n(N,r)$
for all $N,r$, then $M(2,2n+3)$ is graded unitary.
\end{theorem}

\noindent File \leanpath{VOA/\allowbreak Examples/\allowbreak TwoSeriesCharacterAGIdentity.lean}.
\noindent File \leanpath{VOA/Examples/TwoSeriesCharacterGU.lean}.
\begin{leancode}
theorem twoSeries_gradedUnitary_of_affine_invariants_gordon_final (n : ℕ)
    (hA : GradedUnitary (affineForm (twoSeriesLevels n)) (affineRealStructure (twoSeriesLevels n))
      (affineFlag (twoSeriesLevels n)))
    (hinv : ∀ (N : ℤ) (r : ℕ), Module.rank ℂ (affineInvariants (twoSeriesLevels n) r N) ≤
      VOA.TwoSeriesCharacter.gordonCoeff n N r) :
    VirGradedUnitary (twoSeriesLevels n).c (targetRealStructure (twoSeriesLevels n)) :=
\end{leancode}

\begin{dictionary}
Theorem~\ref{l29:thm:invariants} with $\mathrm{coeff}=\ltwentyninegc_n$ and \lean{hpbw} discharged by
Theorem~\ref{l29:thm:hpbw}. The remaining hypothesis \lean{hinv} is the \emph{affine invariant count}: it compares a
space built from $L_{k_n}(\ltwentyninesl)$ and $L_1(\ltwentyninesl)$ with a count of partitions attached to
$M(2,2n+3)$. At $n=1$ it was proved in Chapter~\ref{ch:L28}; for $n\ge2$ it is discharged in
Section~\ref{l29:sec:model}.
\end{dictionary}

\begin{physics}
``The classical gauge invariants of $(A_1,D_{2n+1})\otimes(\text{free }SU(2)_1)$ are no more numerous than the states
of $(A_1,A_{2n})$.''
\end{physics}

\begin{status}
Proved, \texttt{\#print axioms}-clean (re-checked), \textbf{conditional} on \lean{hA} and \lean{hinv}.
\end{status}

%======================================================================
\section{The affine invariant count: socle words and the jet model}\label{l29:sec:model}
%======================================================================

The remaining input is the affine invariant count
\begin{equation}\label{l29:eq:hinv}
\dim\ltwentynineinv_{r,N}\;\le\;\ltwentyninegc_n(N,r)\qquad\text{for all }N\in\Z,\ r\in\N.
\end{equation}
This section explains how it is reduced, with proofs, to statements about one commutative algebra; the last two of
those statements are proved in Sections~\ref{l29:sec:ds} and \ref{l29:sec:arr}.

\subsection{Why the Lee--Yang argument does not generalize}

\begin{warning}[title={Dead end: one cyclic dual vector}]\label{l29:warn:cyclic}
At $n=1$ (Chapter~\ref{ch:L28}) the bound \eqref{l29:eq:hinv} was proved by evaluating an invariant tensor
$\Phi\in(\ltwentyninegr_rA\otimes B)^{\ltwentyninesl[t]}$ against \emph{one} cyclic dual vector $w_r$ of
$\ltwentyninegr_rA$, landing in a sector of $B=L_1$ of charge $2r$ whose dimension is the Rogers--Ramanujan count. For
$n\ge2$ this cannot work. At $n=2$, $r=2$, $N=4$ there is a nonzero invariant tensor (the count is
$\ltwentyninegc_2(4,2)=1$; exact linear algebra over $\mathbb F_p$ confirms the invariant), but its evaluation against
$w_2$ would land in the $L_1$ sector of charge $4$ and energy $N-r=2$, which is zero because that sector starts at
energy $4$. So the restricted dual of $\ltwentyninegr_rA$ is \emph{not} cyclic over $\ltwentyninesl[t]$ for $n\ge2$:
several dual generators are needed, one for each $\ltwentyninesl[t]$-singular vector (the \emph{socle}). (Source:
the docstring of \leanpath{TwoSeriesInvariantsSocle.lean} and the project's numerical records; the
computation is numerical evidence, not a proof, and it is not used by any theorem.)
\end{warning}

\subsection{Socle words}

\begin{definition}[Socle words, their degree, and the socle-word count]\label{l29:def:socle}
A \emph{socle word} of size $r$ is a word $x=(x_1,\dots,x_{N_1})$ in the letters $\{0,\dots,n-1\}$ with
$\sum_j(x_j+1)=r$ (a letter $t$ is a cluster of size $t+1$). With $c_t(x)=\#\{j:x_j\ge t\}$ and $\mathrm{inv}(x)$ the
number of inversions, its degree is $\deg x=N_1+\sum_{t=1}^{n-1}(c_t^2+c_t)+\mathrm{inv}(x)$. The \emph{socle-word
count} is
\begin{equation*}
\mathrm{swc}_n(N,r)=\sum_{x}\#\bigl\{\text{partitions of }N-N_1^2-\deg x\text{ into at most }N_1\text{ parts}\bigr\},
\end{equation*}
the sum over socle words of size $r$ (of any length $N_1\le r$).
\end{definition}

\noindent File \leanpath{VOA/\allowbreak Examples/\allowbreak TwoSeriesInvariantsSocle.lean}.
\begin{leancode}
def SocleWord (n r N₁ : ℕ) : Type := {x : Fin N₁ → Fin n // ∑ j, ((x j : ℕ) + 1) = r}

def clusterAtLeast {n N₁ : ℕ} (x : Fin N₁ → Fin n) (t : ℕ) : ℕ :=
  (Finset.univ.filter fun j => t ≤ (x j : ℕ)).card

def inversions {n N₁ : ℕ} (x : Fin N₁ → Fin n) : ℕ :=
  (Finset.univ.filter fun p : Fin N₁ × Fin N₁ => p.1 < p.2 ∧ x p.2 < x p.1).card

def socleDegree {n N₁ : ℕ} (x : Fin N₁ → Fin n) : ℕ :=
  N₁ + ∑ t ∈ Finset.Ico 1 n, (clusterAtLeast x t ^ 2 + clusterAtLeast x t) + inversions x

abbrev SocleIndex (n r : ℕ) : Type := Σ N₁ : Fin (r + 1), SocleWord n r N₁

def socleWordCount (n : ℕ) (N : ℤ) (r : ℕ) : ℕ :=
  ∑ i : SocleIndex n r, LeeYangTransferAssembly.gapTwoCoefficient
    (N - (socleDegree i.2.1 : ℤ) + ((i.1 : ℕ) : ℤ)) (i.1 : ℕ)
\end{leancode}

\begin{dictionary}
\lean{{x : X // P x}} is the subtype of elements satisfying $P$. \lean{Σ N₁ : Fin (r + 1), ...} is a dependent pair:
a length $N_1\le r$ together with a word of that length. \lean{gapTwoCoefficient M s} is the Lee--Yang coefficient of
Chapter~\ref{ch:L28}, the number of partitions of $M-s(s+1)$ into at most $s$ parts (equivalently the coefficient of
$q^M$ in $q^{s(s+1)}/(q)_s$); with $M=N-\deg x+N_1$ and $s=N_1$ this is the number of partitions of $N-\deg x-N_1^2$
into at most $N_1$ parts. The definitions match the textbook ones.
\end{dictionary}

\begin{example}[$n=2$, $r=2$]\label{l29:ex:socle}
The socle words of size $2$ in the letters $\{0,1\}$ are $(0,0)$ (two clusters of size one: $N_1=2$, $c_1=0$, no
inversion, degree $2$) and $(1)$ (one cluster of size two: $N_1=1$, $c_1=1$, degree $1+2=3$). So
$\mathrm{swc}_2(N,2)=\#\{\lambda\vdash N-6,\ \le2\text{ parts}\}+\#\{\lambda\vdash N-4,\ \le1\text{ part}\}$. At $N=4$:
$0+1=1$; at $N=6$: $1+1=2$; at $N=8$: $2+1=3$. These are $\ltwentyninegc_2(N,2)$ from Example~\ref{l29:ex:gordon}.
\end{example}

\begin{physics}
A socle word labels a lowest-weight $\ltwentyninesl[t]$-singular state of the classical limit of $L_{k_n}$: $N_1$
clusters of currents, a cluster of size $t+1$ being a bound state of $t+1$ current quanta. The degree is the energy of
the bound-state configuration; the partitions count its descendants. This is the bound-state (``fermionic'') picture of
the $(A_1,D_{2n+1})$ chiral algebra.
\end{physics}

\begin{status}
Definitions.
\end{status}

\begin{theorem}[The socle-word count is the Andrews--Gordon count]\label{l29:thm:swc}
For all $n,N,r\in\N$, $\mathrm{swc}_n(N,r)=\ltwentyninegc_n(N,r)$.
\end{theorem}

\noindent File \leanpath{VOA/\allowbreak Examples/\allowbreak TwoSeriesSocleCountGordon.lean}.
\begin{leancode}
theorem socleWordCount_eq (n N r : ℕ) :
    socleWordCount n (N : ℤ) r = gordonCoeff n (N : ℤ) r := by
\end{leancode}

\begin{dictionary}
Full name \ltwentynineln{VOA.SocleCount.socleWordCount_eq}. An equality of natural numbers for all $n,N,r$. The corollary
\lean{socleWordGordon n} is the inequality form $\mathrm{swc}_n\le\ltwentyninegc_n$ (for all integer $N$), which the
docstring of \lean{SocleWordGordon} in \leanpath{TwoSeriesInvariantsSocle.lean} still calls ``OPEN''
(warning on p.~\pageref{l29:warn:stale}).
\end{dictionary}

\begin{proof}[Proof on paper]
$\#\{\text{partitions of }M\text{ into}\le s\text{ parts}\}$ is the $q^M$ coefficient of $1/(q)_s$. For a fixed
content (the multiset of letters, i.e.\ the cluster profile), MacMahon's theorem says that the inversion generating
function of the words with that content is the $q$-multinomial coefficient. Summing, the socle-word count becomes the
fixed-$r$ coefficient of the Andrews (fermionic) multisum
$\sum q^{N_1^2+\cdots+N_{n-1}^2+N_1+\cdots+N_{n-1}}/\prod(q)_{N_i-N_{i+1}}$ with an extra variable marking $r$
(\lean{sum_words_eq}, \lean{coeff_F_one_eq}). That multisum satisfies Andrews' recursion and therefore equals the
Gordon series of Theorem~\ref{l29:thm:ag} (\lean{F_eq_GG}).
\end{proof}

\begin{physics}
The bound-state count of the affine side equals the exclusion-statistics count of the Virasoro side: the two
quasi-particle descriptions of the coset agree.
\end{physics}

\begin{status}
Proved, unconditional. Axioms standard (re-checked).
\end{status}

\begin{definition}[Socle data of a mode module]\label{l29:def:soclespectrum}
Let $M$ be a $\C$-vector space with operators $\rho(a,j)$ ($a\in\{e,f,h\}$, $j\ge0$) and a grading $M=\bigoplus_dM_d$.
The \emph{lowest-weight singular vectors} of degree $d$ and charge $c$ are the $s\in M_d$ with $\rho(a,j+1)s=0$ for all
$a,j$, $f_0s=0$ and $h_0s=cs$. Given labels $i$ with degrees $\deg i$ and cluster numbers $\mathrm{cl}(i)$, and a
``Cartan image'' $\mathrm{Cart}(N)\subseteq M$, \emph{socle data} means: (rank) the singular space of degree $d$ and
charge $c$ has dimension at most the number of labels with $\deg i=d$ and $c=-2\,\mathrm{cl}(i)$; (transversality) a
singular vector of charge $-2N$ and degree $d$ that lies in $\mathrm{Cart}(N)+\sum_{d'\ne d}M_{d'}+\operatorname{im}(h_0+2N)$
is zero.
\end{definition}

\noindent File \leanpath{VOA/\allowbreak Examples/\allowbreak TwoSeriesSocleDualCore.lean}.
\begin{leancode}
def sing (d : ℤ) (c : ℂ) : Submodule ℂ M where
  carrier := {s | s ∈ Mdeg d ∧ (∀ a j, ρ (a, j + 1) s = 0) ∧ ρ (1, 0) s = 0 ∧
    ρ (2, 0) s = c • s}

def avoid (d : ℤ) (N : ℕ) : Submodule ℂ M :=
  Cart N ⊔ (⨆ (d' : ℤ) (_ : d' ≠ d), Mdeg d') ⊔
    LinearMap.range (ρ (2, 0) + (2 * (N : ℂ)) • (1 : Module.End ℂ M))

structure SocleSpectrum : Prop where
  rank_le : ∀ (d : ℤ) (c : ℂ), finrank ℂ (sing ρ Mdeg d c) ≤
    Nat.card {i : ι // (deg i : ℤ) = d ∧ c = -(2 * (cl i : ℂ))}
  transverse : ∀ (d : ℤ) (N : ℕ),
    Disjoint (sing ρ Mdeg d (-(2 * (N : ℂ)))) (avoid ρ Mdeg Cart d N)
\end{leancode}

\begin{dictionary}
\lean{ρ (a, j)} with \lean{a : Fin 3} ($0=e$, $1=f$, $2=h$) is the mode $J^a_j$; so \lean{ρ (1, 0)} is $f_0$ and
\lean{ρ (2, 0)} is $h_0$. \lean{⊔} is the sum of submodules and \lean{⨆} an indexed sum; \lean{LinearMap.range} the
image. \lean{Disjoint U W} for submodules means $U\cap W=0$. \lean{Nat.card} of the subtype of labels is the number of
labels with the stated degree and charge. The proof fields of the omitted part of the definition (closure under
addition and scalars) are routine. Same as the textbook definition. (The source writes \lean{sing} with its proof
fields; the quotation omits them, which is why the block stops at \lean{carrier}.)
\end{dictionary}

\begin{physics}
The socle is the space of ``primary'' classical states for the current algebra: those annihilated by all positive
modes and by $f_0$. The rank clause says the bound-state labels exhaust them; transversality says they are detected by
their charge and energy.
\end{physics}

\begin{status}
Definitions.
\end{status}

\begin{theorem}[Transfer from the socle data of $\ltwentyninegr_rA$]\label{l29:thm:soclespectrum}
For every $n$: if $L_{k_n}(\ltwentyninesl)$ is graded unitary and $\ltwentyninegr_rA$ (for every $r$, with the socle
words of size $r$ as labels and the Cartan image $I_N\cdot\ltwentyninegr_rA$) has socle data, then $M(2,2n+3)$ is
graded unitary.
\end{theorem}

\noindent File \leanpath{VOA/\allowbreak Examples/\allowbreak TwoSeriesSocleDual.lean}.
\begin{leancode}
structure TwoSeriesSocleSpectrum (n : ℕ) : Prop where
  rank_le : ∀ (r : ℕ) (d : ℤ) (c : ℂ),
    finrank ℂ (sing (grRho (twoSeriesLevels n) r) (grDeg (twoSeriesLevels n) r) d c) ≤
      Nat.card {i : SocleIndex n r //
        (socleDegree i.2.1 : ℤ) = d ∧ c = -(2 * (((i.1 : ℕ) : ℕ) : ℂ))}
  transverse : ∀ (r : ℕ) (d : ℤ) (N : ℕ),
    Disjoint (sing (grRho (twoSeriesLevels n) r) (grDeg (twoSeriesLevels n) r) d
        (-(2 * (N : ℂ))))
      (avoid (grRho (twoSeriesLevels n) r) (grDeg (twoSeriesLevels n) r)
        (grCart (twoSeriesLevels n) r) d N)

theorem twoSeries_gradedUnitary_of_socleSpectrum (n : ℕ)
    (hA : GradedUnitary (affineForm (twoSeriesLevels n)) (affineRealStructure (twoSeriesLevels n))
      (affineFlag (twoSeriesLevels n)))
    (S : TwoSeriesSocleSpectrum n) :
    VirGradedUnitary (twoSeriesLevels n).c (targetRealStructure (twoSeriesLevels n)) :=
  twoSeries_gradedUnitary_of_socleInput n hA S.toSocleInput
\end{leancode}

\begin{dictionary}
\lean{grRho L r} and \lean{grDeg L r} are the induced mode action and the conformal grading on $\ltwentyninegr_rA$;
\lean{grCart L r N} is the image of the Cartan ideal $I_N$ (polynomials in the Cartan modes $h_{i+1}$ acting through
$H_{i+1}\mapsto-J^h_{i+1}$). \lean{TwoSeriesSocleSpectrum n} is socle data (Definition~\ref{l29:def:soclespectrum})
of the actual $\ltwentyninegr_rA$ for every $r$. Its docstring still calls it ``OPEN''; it is now derivable for every
$n\ge1$ from the results of this and the next two sections, though no declaration states it as a standalone theorem
(warning on p.~\pageref{l29:warn:stale}).
\end{dictionary}

\begin{proof}[Proof on paper]
From the socle data one constructs, by duality, a family of functionals on $\ltwentyninegr_rA$, one per socle word,
that is \emph{jointly cyclic}: their translates by mode words separate points (\lean{exists_socleDuals}). Each is
supported in one degree, has charge $2N_1$ and is killed by the Cartan ideal $I_{N_1}$. Evaluating an invariant tensor
against them maps $\ltwentynineinv_{r,N}$ injectively into a sum of $L_1$ sectors; each sector's Cartan-annihilated
dual is bounded by the Heisenberg separation of $L_1$ (\lean{levelOneHeisenbergSeparation}), giving
$\dim\ltwentynineinv_{r,N}\le\mathrm{swc}_n(N,r)$ (\ltwentynineln{twoSeries_affineInvariants_rank_le_socleWordCount}). By
Theorem~\ref{l29:thm:swc} this is $\ltwentyninegc_n(N,r)$, and Theorem~\ref{l29:thm:gordonfinal} applies.
\end{proof}

\begin{physics}
Each bound-state primary of the $(A_1,D_{2n+1})$ classical limit pairs with a lattice sector of the free $SU(2)_1$
boson; the invariant states are counted sector by sector.
\end{physics}

\begin{status}
Proved, \texttt{\#print axioms}-clean, \textbf{conditional} on \lean{hA} and \lean{S}. (The intermediate
\ltwentynineln{twoSeries_gradedUnitary_of_socleInput}, re-checked: axioms standard.)
\end{status}

\subsection{The jet model}

The socle data of $\ltwentyninegr_rA$ are hard to compute directly. The route taken embeds $\ltwentyninegr_rA$ into an
explicit commutative model, proves the socle data for the model, and pulls them back along the injective intertwiner
(\lean{SocleSpectrum.comap}).

\begin{definition}[The coefficient ring $\ltwentynineRn$]\label{l29:def:Rn}
Let $\C[T]=\C[T_0,T_1,\dots]$ and $T(s)=\sum_jT_js^j$. The \emph{jet ideal} of $T^n$ is the ideal generated by the
coefficients of $T(s)^n$, and $\ltwentynineRn=\C[T]/(\text{jet ideal of }T^n)$, with $t_j$ the image of $T_j$.
\end{definition}

\noindent File \leanpath{VOA/\allowbreak Examples/\allowbreak TwoSeriesStrictEndpoint.lean}.
\noindent File \leanpath{VOA/Examples/TwoSeriesJetRing.lean}.
\begin{leancode}
abbrev TPoly : Type := MvPolynomial ℕ ℂ

def tSer : PowerSeries TPoly := PowerSeries.mk X

def jetIdealT (n : ℕ) : Ideal TPoly :=
  Ideal.span (Set.range fun j => PowerSeries.coeff j (tSer ^ n))

abbrev Rn (n : ℕ) : Type := TPoly ⧸ jetIdealT n

def tn (n : ℕ) (j : ℕ) : Rn n := Ideal.Quotient.mk (jetIdealT n) (X j)
\end{leancode}

\begin{dictionary}
\lean{MvPolynomial ℕ ℂ} is the polynomial ring in countably many variables $X_0,X_1,\dots$; \lean{PowerSeries.mk X}
is the series $\sum_jX_js^j$ with coefficients in that ring; \lean{PowerSeries.coeff j} extracts the coefficient of
$s^j$. \lean{Ideal.span (Set.range f)} is the ideal generated by the values of $f$; \lean{⧸} is the quotient ring.
At $n=1$ the jet ideal contains every $T_j$, so $R_1=\C$ (\lean{jetIdealT_one}).
\end{dictionary}

\begin{example}[$n=2$]
The jet ideal of $T^2$ is generated by $T_0^2$, $2T_0T_1$, $2T_0T_2+T_1^2$, $2T_0T_3+2T_1T_2$, \dots. So in $R_2$:
$t_0^2=0$, $t_0t_1=0$, $t_1^2=-2t_0t_2$, and so on. $R_2$ is the arc algebra of the double point $\C[T]/(T^2)$, which by
classical freeness (Chapter~\ref{ch:L20}) is $\ltwentyninegr^{\mathrm{Li}}$ of $M(2,5)$, the Lee--Yang model.
\end{example}

\begin{physics}
$\ltwentynineRn$ is the ring of classical (arc-space) configurations of the stress tensor of $M(2,2n+1)$: $T^n=0$ with
all its derivatives. It enters because $L_{k_n}$ reduces, under Drinfeld--Sokolov, to $M(2,2n+1)$, and the model
remembers the stress tensor of the reduction.
\end{physics}

\begin{status}
Definitions.
\end{status}

\begin{definition}[The model $\ltwentynineS$]\label{l29:def:model}
Let $\ltwentynineRn[u_j,v_j]$ be the polynomial ring in two sequences of jet variables over $\ltwentynineRn$. Each
generator letter of the corrected PBW basis of $V^k(\ltwentyninesl)$ gets a \emph{model symbol}: a current letter
$(J^a,j)$ gets the $j$-th jet of $u^2$, $uv$ or $v^2$ (the Adamovi\'c symbols of Chapter~\ref{ch:L28}); the Sugawara
letter $(L,j)$ gets the $j$-th jet of $-(u'v-uv')$ plus $t_j$. A word's symbol is the product of its letters' symbols.
The model $\ltwentynineS$ is the span of the symbols of words of length $r$; its degree-$d$ piece is the span of symbols
of words of length $r$ and weight $d$.
\end{definition}

\noindent Files \leanpath{VOA/\allowbreak Examples/\allowbreak TwoSeriesStrictRing.lean} and
\leanpath{VOA/\allowbreak Examples/\allowbreak TwoSeriesStrictWords.lean}.
\begin{leancode}
abbrev MPoly : Type := MvPolynomial (Fin 2 × ℕ) R

def mletter : Generator × ℕ → MPoly R
  | (.current a, n) => ι R (letterSym (.current a, n))
  | (.sugawara, n) => ι R (letterSym (.sugawara, n)) + MvPolynomial.C (t n)

def mword (w : List (Generator × ℕ)) : MPoly R := (w.map (mletter R t)).prod

def modelImage (r : ℕ) : Submodule ℂ (MPoly R) :=
  Submodule.span ℂ {Q | ∃ w : List (Generator × ℕ), w.length = r ∧ mword R t w = Q}

def modelDeg (r : ℕ) (d : ℤ) : Submodule ℂ (MPoly R) :=
  Submodule.span ℂ
    {Q | ∃ w : List (Generator × ℕ), w.length = r ∧ wordWt w = d ∧ mword R t w = Q}
\end{leancode}

\begin{dictionary}
The variables of \lean{MPoly R} are indexed by \lean{Fin 2 × ℕ}: $(0,j)\mapsto u_j$, $(1,j)\mapsto v_j$. \lean{ι R}
includes $\C[u_j,v_j]$ into $R[u_j,v_j]$; \lean{letterSym} is the Adamovi\'c letter symbol (Chapter~\ref{ch:L28});
\lean{MvPolynomial.C (t n)} is the constant $t_n\in R$. \lean{Generator} has the constructors \lean{.current a}
($a$ in \lean{Fin 3}) and \lean{.sugawara}. \lean{wordWt w} is the conformal weight of the word. The ring $R$ and
$t:\N\to R$ are parameters; the model of the series is $R=\ltwentynineRn$, $t=t_\bullet$. So
$\ltwentynineS=$ \lean{modelImage (Rn n) (tn n) r}.
\end{dictionary}

\begin{physics}
This is a free-field (Wakimoto-type) realisation of the classical limit of $L_{k_n}$: the currents are bilinears in a
pair of symplectic bosons $(u,v)$, and the stress tensor is the bilinear $-(u'v-uv')$ corrected by the stress tensor
$T$ of the Drinfeld--Sokolov reduction. The nilpotency $T^n=0$ is what distinguishes level $k_n$.
\end{physics}

\begin{status}
Definitions.
\end{status}

\begin{definition}[Detection and the model character bound]\label{l29:def:detect}
\emph{Detection} at member $n$: for every $r$, if a pair (state $y$, symbol $Q$) lies in the span of the pairs
(state of a length-$r$ word, its model symbol), modulo (lower PBW length)$\times0$, and $y$ lies in the maximal ideal of
$V^{k_n}(\ltwentyninesl)$, then $Q$ evaluated at the base point $(u,v)=(0,1)$ vanishes in $\ltwentynineRn$.
\emph{Model character bound}: for every $N$, the $q^N$ coefficient of the Arakawa--Dai product
$\bigl((q^{2n+1};q^{2n+1})_\infty/(q;q)_\infty\bigr)^3$ is at most $\sum_{r\le N}\dim(\ltwentynineS)_N$.
\end{definition}

\noindent Files \leanpath{VOA/\allowbreak Examples/\allowbreak TwoSeriesStrictWords.lean},
\leanpath{VOA/\allowbreak Examples/\allowbreak TwoSeriesStrictDescent.lean} and
\leanpath{VOA/\allowbreak Examples/\allowbreak TwoSeriesStrictEndpoint.lean}.
\begin{leancode}
def symModM (r : ℕ) : Submodule ℂ (WakModule k × MPoly R) :=
  symPairsM R t k hk r ⊔ (lowerP k hk r).prod ⊥

def Detects (I : Submodule ℂ (WakModule k)) (r : ℕ) : Prop :=
  ∀ (y : WakModule k) (Q : MPoly R), (y, Q) ∈ symModM R t k hk r → y ∈ I → evR R Q = 0

def ModelDetects : Prop :=
  ∀ r, Detects R t L.kA L.aNoncritical (ImaxL L) r

def TwoSeriesModelDetects (n : ℕ) : Prop :=
  ModelDetects (twoSeriesLevels n) (Rn n) (tn n)

def TwoSeriesModelCharacter (n : ℕ) : Prop :=
  ∀ N : ℕ, PowerSeries.coeff N (AffineSl2BoundaryCharacter.product n) ≤
    ∑ r ∈ Finset.range (N + 1), (Module.finrank ℂ (modelDeg (Rn n) (tn n) r (N : ℤ)) : ℚ)
\end{leancode}

\begin{dictionary}
\lean{symPairsM} is the span of the pairs (\lean{correctedWordState} of a length-$r$ word, \lean{mword} of it), and
\lean{(lowerP k hk r).prod ⊥} is $F_{r-1}\times0$ (lower PBW length, zero symbol). \lean{ImaxL L} is the maximal ideal of
$V^{k_A}$. \lean{evR R} is evaluation at $(u_j,v_j)=$ the base point $(0,1)$ (only $v_0=1$, all other variables $0$),
with values in $R$. \lean{AffineSl2BoundaryCharacter.product n} is the Arakawa--Dai product of Chapter~\ref{ch:L20},
Definition~\ref{l20:def:character}, which equals the vacuum character of $L_{k_n}$ for $n\ge1$
(\lean{character_eq_product}). The inequality is in $\Q$.

\textbf{Why detection is the right condition.} Detection is equivalent to the model symbol map being well defined on
the simple quotient $A$ (\lean{symbol_eq_zero_of_detect}; conversely trivially). The base-point evaluation suffices
because the model is cyclic under the modes (\lean{m_cyclic}).
\end{dictionary}

\begin{physics}
Detection says that every null state of $L_{k_n}$ has zero classical symbol after restriction to the Kostant slice
$f+\C e$; i.e.\ the nulls of the affine theory are invisible on the slice where the Drinfeld--Sokolov reduction lives.
The character bound says that the model has at least as many states as $L_{k_n}$.
\end{physics}

\begin{status}
Definitions (predicates). Detection is proved for every $n\ge1$ in Section~\ref{l29:sec:ds}
(Theorem~\ref{l29:thm:detect}); its $T$-free part was proved earlier at every non-integral level
(\lean{evR_aug_eq_zero}); at $n=1$ it is \lean{twoSeriesModelDetects_one} (re-checked). The character bound is
Theorem~\ref{l29:thm:modelchar}.
\end{status}

\begin{definition}[Strict model maps; model socle data]\label{l29:def:strict}
A \emph{strict model map} at $(n,r)$ is an injective linear map $\ltwentyninegr_rA\to\ltwentynineS$ that intertwines the
mode actions and preserves the conformal degree. \emph{Model socle data} at $(n,r)$ are socle data
(Definition~\ref{l29:def:soclespectrum}) for $\ltwentynineS$ with its own mode action and Cartan image, labelled by the
socle words of size $r$.
\end{definition}

\noindent File \leanpath{VOA/\allowbreak Examples/\allowbreak TwoSeriesStrictEndpoint.lean}.
\begin{leancode}
structure StrictModelMap (n r : ℕ) where
  f : GrA (twoSeriesLevels n) r →ₗ[ℂ] modelImage (Rn n) (tn n) r
  injective : Function.Injective f
  intertwine : ∀ k, f ∘ₗ grRho (twoSeriesLevels n) r k = mImgRho (Rn n) (tn n) r k ∘ₗ f
  degree : ∀ d, ∀ m ∈ grDeg (twoSeriesLevels n) r d, f m ∈ mImgDeg (Rn n) (tn n) r d

def ModelSocleSpectrum (n r : ℕ) : Prop :=
  SocleSpectrum (mImgRho (Rn n) (tn n) r) (mImgDeg (Rn n) (tn n) r)
    (fun i : SocleIndex n r => socleDegree i.2.1) (fun i => (i.1 : ℕ))
    (modelCart (mImgRho (Rn n) (tn n) r) (mImg_cartanCommute (Rn n) (tn n) r))
\end{leancode}

\begin{dictionary}
\lean{→ₗ[ℂ]} is a $\C$-linear map, \lean{∘ₗ} composition of linear maps. \lean{mImgRho} is the mode action on the model
(the jet derivations of $u,v$ together with the derivation of $\ltwentynineRn$), \lean{mImgDeg} the degree pieces,
\lean{modelCart} the image of the Cartan ideal on the model. The labels are the socle words of
Definition~\ref{l29:def:socle}: degree \lean{socleDegree}, cluster number $N_1$.
\end{dictionary}

\begin{physics}
A strict model map is a faithful free-field realisation of the classical limit, compatible with the current algebra
and with energy. Model socle data are the bound-state spectrum computed in the free-field model.
\end{physics}

\begin{status}
Definitions.
\end{status}

\begin{theorem}[Transfer from the strict model]\label{l29:thm:strict}
For every $n\ge1$: graded unitarity of $L_{k_n}$, detection, the model character bound and model socle data for every
$r$ imply graded unitarity of $M(2,2n+3)$.
\end{theorem}

\noindent File \leanpath{VOA/\allowbreak Examples/\allowbreak TwoSeriesStrictEndpoint.lean}.
\begin{leancode}
theorem twoSeries_gradedUnitary_of_strictModel (n : ℕ) (hn : 1 ≤ n)
    (hA : GradedUnitary (affineForm (twoSeriesLevels n)) (affineRealStructure (twoSeriesLevels n))
      (affineFlag (twoSeriesLevels n)))
    (det : TwoSeriesModelDetects n) (chi : TwoSeriesModelCharacter n)
    (spec : ∀ r, ModelSocleSpectrum n r) :
    VirGradedUnitary (twoSeriesLevels n).c (targetRealStructure (twoSeriesLevels n)) :=
  twoSeries_gradedUnitary_of_socleModels n hA
    fun r => twoSeriesSocleModel_of_strict hn det chi r (spec r)
\end{leancode}

\begin{dictionary}
Full name \ltwentynineln{VOA.TwoSeries.StrictModel.twoSeries_gradedUnitary_of_strictModel}. The hypothesis $1\le n$ is needed
because the equality of the Arakawa--Dai product with the character is proved only for $n\ge1$ (\lean{character_eq_product}; at
$n=0$ both sides are $1$, but no declaration states it). The strict
maps are \lean{twoSeriesStrictModel hn det chi r}: the model symbol is well defined on $A$ by detection, and injective
by a dimension squeeze against the character bound (\ltwentynineln{ModelSymbol.injective_of_bound}).
\end{dictionary}

\begin{proof}[Proof on paper]
Detection makes the symbol map $\sigma:\ltwentyninegr_rA\to\ltwentynineS$ well defined; it is surjective onto
$\ltwentynineS$ by construction (both are spanned by words). Summed over $r$, $\dim A_N=\sum_r\dim\ltwentyninegr_rA_N$
(\lean{sum_finrank_grDegL}) is the character of $A$, which by the model character bound is at most
$\sum_r\dim(\ltwentynineS)_N$; a surjection between spaces whose dimensions go the other way is an isomorphism, so each
$\sigma$ is injective. Socle data pull back along injective intertwiners (\lean{SocleSpectrum.comap}); apply
Theorem~\ref{l29:thm:soclespectrum}.
\end{proof}

\begin{physics}
If the free-field model is faithful (no affine null state survives in it) and big enough (at least as many states), it
is exact, and its bound-state spectrum is that of $L_{k_n}$.
\end{physics}

\begin{status}
Proved, \texttt{\#print axioms}-clean (re-checked), \textbf{conditional} on \lean{hA}, \lean{det}, \lean{chi},
\lean{spec}.
\end{status}

\subsection{The model character: the Bailey chain}\label{l29:sec:modelchar}

\begin{definition}[The Chern--Tran--Wakhare fermionic series]\label{l29:def:ferm}
Let $\beta^{(0)}_N=\delta_{N,0}$ and $\beta^{(n+1)}_N=\sum_{k\le N}q^{k^2+k}\beta^{(n)}_k/(q)_{N-k}$ (the Bailey chain
started at the unit sequence). The \emph{fermionic series} is
\begin{equation*}
\mathrm{ferm}_n(q)\;=\;\sum_{N\ge0}q^N\Bigl(\sum_{k\le 2N}\begin{bmatrix}2N\\k\end{bmatrix}_q\Bigr)\beta^{(n)}_N .
\end{equation*}
\end{definition}

\noindent File \leanpath{VOA/\allowbreak Examples/\allowbreak TwoSeriesModelCharFermionic.lean}.
\begin{leancode}
def bet : ℕ → ℕ → PowerSeries ℚ
  | 0, N => if N = 0 then 1 else 0
  | n + 1, N => ∑ k ∈ Finset.range (N + 1),
      PowerSeries.X ^ (k * k + k) * reciprocalQFactorial (N - k) * bet n k

def fermTerm (n N : ℕ) : PowerSeries ℚ :=
  PowerSeries.X ^ N * (∑ k ∈ Finset.range (2 * N + 1), gaussian PowerSeries.X (2 * N) k) *
    bet n N

def fermSeries (n : ℕ) : PowerSeries ℚ := VOA.LocallyFinitePowerSeriesSum.sum (fermTerm n)
\end{leancode}

\begin{dictionary}
\lean{PowerSeries.X} is $q$; \lean{reciprocalQFactorial m} is $1/(q)_m$; \lean{gaussian PowerSeries.X M k} the Gaussian
binomial $\left[\begin{smallmatrix}M\\k\end{smallmatrix}\right]_q$. \ltwentynineln{LocallyFinitePowerSeriesSum.sum} is the sum of
a sequence of power series in which each coefficient receives finitely many contributions (here \lean{fermTerm n N}
starts at $q^N$), Chapter~\ref{ch:L20}, Definition~\ref{l20:def:lfsum}. Unwinding the recursion,
$\beta^{(n)}_N=\sum_{N\ge N_1\ge\cdots\ge N_{n-1}\ge0}q^{\sum_s(N_s^2+N_s)}/\prod_s(q)_{N_s-N_{s+1}}$ (with the
conventions $N_0=N$, $N_n=0$, the sums and products over $1\le s\le n-1$ and $0\le s\le n-1$ respectively), a Gordon
colour-profile sum.
\end{dictionary}

\begin{physics}
A quasi-particle formula for the vacuum character of $L_{k_n}(\ltwentyninesl)$: $N$ counts current quasi-particles,
$\sum_k[2N,k]_q$ their $\ltwentyninesl$ weights, and $\beta^{(n)}$ the clustering into bound states of size $\le n$.
\end{physics}

\begin{status}
Definition. At $n=1$ it is the level-one jet series of Chapter~\ref{ch:L28} (\lean{fermSeries_one}).
\end{status}

\begin{theorem}[The fermionic series is the product; the model character bound]\label{l29:thm:modelchar}
For every $n\in\N$, $\mathrm{ferm}_n=\bigl((q^{2n+1};q^{2n+1})_\infty/(q;q)_\infty\bigr)^3$. For every $n\ge1$ the
model character bound holds.
\end{theorem}

\noindent Files \leanpath{VOA/\allowbreak Examples/\allowbreak TwoSeriesModelCharFermionic.lean} and
\leanpath{VOA/\allowbreak Examples/\allowbreak TwoSeriesModelCharEndpoint.lean}.
\begin{leancode}
theorem fermSeries_eq_product (n : ℕ) :
    fermSeries n = AffineSl2BoundaryCharacter.product n := by

theorem twoSeriesModelCharacter_all (n : ℕ) (hn : 1 ≤ n) : TwoSeriesModelCharacter n := by
  intro N
  rw [← fermSeries_eq_product n]
  exact fermSeries_le_model hn N
\end{leancode}

\begin{dictionary}
Full names \ltwentynineln{VOA.TwoSeries.ModelChar.fermSeries_eq_product} and \ltwentynineln{VOA.TwoSeries.ModelChar.twoSeriesModelCharacter_all}.
The first is an equality of power series over $\Q$ (all coefficients), for every $n$. The second rewrites the product
as the fermionic series and applies the lower bound \lean{fermSeries_le_model}: the $q^N$ coefficient of the fermionic
series is at most $\sum_{r\le N}\dim(\ltwentynineS)_N$. No hypothesis beyond $n\ge1$.
\end{dictionary}

\begin{proof}[Proof on paper]
\emph{Product.} Multiply by $(q;q)_\infty^3$. A terminating $q$-binomial identity
$\sum_m[M,m]_q(-1)^mq^{\binom{M-m}{2}}(x;q)_m=q^{\binom M2}x^M$ (\lean{newton_sum}) shows, by induction on $n$, that
the Bailey transform of $\beta^{(n)}$ is a pure power $q^{nj(j+1)+\binom j2}$ (\lean{cTr_eq}). Combined with a moment
identity for $(q)_\infty^3\sum_Nq^N\sum_k[2N,k]_q(\cdots)$ and Jacobi's identity
$(q)_\infty^3=\sum_j(-1)^j(2j+1)q^{j(j+1)/2}$, this gives $(q;q)_\infty^3\,\mathrm{ferm}_n=(q^{2n+1};q^{2n+1})_\infty^3$
(\lean{euler_cube_mul_fermSeries}). \emph{Lower bound.} For each Gordon colour profile one writes down explicit
elements of $\ltwentynineS$ (block covariants $T^cG$ times symmetrised Feigin--Stoyanovsky functions) and functionals
that are triangular against them (\lean{fFS}, \lean{gFS}, \lean{LamG}); triangularity proves linear independence, and
the number of these elements in weight $N$ is the $q^N$ coefficient of the fermionic series
(\lean{fermSeries_le_model}).
\end{proof}

\begin{physics}
The vacuum character of the $(A_1,D_{2n+1})$ chiral algebra has the quasi-particle form of Chern, Tran and Wakhare and
the product form of Kac--Wakimoto/Arakawa--Dai; the theorem identifies the two and shows the free-field model is
big enough to carry all of it.
\end{physics}

\begin{status}
Proved, \textbf{unconditional} (the second for $n\ge1$). Axioms standard (re-checked, both).
\end{status}

\begin{theorem}[Transfer with the character bound discharged]\label{l29:thm:strictdet}
For every $n\ge1$: $\ltwentynineGU(L_{k_n})$, detection and model socle data imply $\ltwentynineGU(M(2,2n+3))$.
\end{theorem}

\noindent File \leanpath{VOA/\allowbreak Examples/\allowbreak TwoSeriesModelCharEndpoint.lean}.
\begin{leancode}
theorem twoSeries_gradedUnitary_of_strictModel_det (n : ℕ) (hn : 1 ≤ n)
    (hA : GradedUnitary (affineForm (twoSeriesLevels n)) (affineRealStructure (twoSeriesLevels n))
      (affineFlag (twoSeriesLevels n)))
    (det : TwoSeriesModelDetects n) (spec : ∀ r, ModelSocleSpectrum n r) :
    VirGradedUnitary (twoSeriesLevels n).c (targetRealStructure (twoSeriesLevels n)) :=
  twoSeries_gradedUnitary_of_strictModel n hn hA det (twoSeriesModelCharacter_all n hn) spec
\end{leancode}

\begin{dictionary}
Theorem~\ref{l29:thm:strict} with \lean{chi} supplied by Theorem~\ref{l29:thm:modelchar}. Full name
\ltwentynineln{VOA.TwoSeries.ModelChar.twoSeries_gradedUnitary_of_strictModel_det}.
\end{dictionary}

\begin{physics}
What is left: the nulls are invisible on the slice (detection), and the bound-state spectrum of the free-field model
is the expected one (socle data).
\end{physics}

\begin{status}
Proved, \texttt{\#print axioms}-clean (re-checked), \textbf{conditional} on \lean{hA}, \lean{det}, \lean{spec}.
\end{status}

\subsection{The model socle data, and the jet-ideal Poincar\'e lemma}\label{l29:sec:poincare}

The model socle data have two halves. \textbf{Transversality} is proved outright for every coefficient ring with a
scaling automorphism (\lean{model_transverse}, re-checked: axioms standard): the functional ``coefficient of $v_0^{2N}$''
kills $\operatorname{im}(h_0+2N)$ and the Cartan image, separates degrees through the scaling $t_j\mapsto2^{j+1}t_j$, and
is injective on singular vectors of charge $-2N$, which all have the form $\tau\,v_0^{2N}$. The \textbf{rank} half is
reduced, through a chain of files, to a statement about $\C[T]$ alone:
\begin{enumerate}
\item A singular vector of degree $d$ and charge $-q$ in $\ltwentynineS$ is $\tau\,v_0^{q}$ with $\tau$ in the
  \emph{membership space} $\{\tau\in\ltwentynineRn:\tau v_0^q\in(\ltwentynineS)_d\}$, injectively; odd $q$ give nothing
  (\leanpath{TwoSeriesModelSocleEndpoint}).
\item The membership space for $N$ clusters lies in the image of $T$-polynomials of degree $r-N$ and weight $d-r$, and the
  cases $N\ge r$ are immediate (\leanpath{TwoSeriesMembershipLine}).
\item \emph{Box count} (proved for every $n,r$: \lean{boxCount_all}, re-checked): the polynomials in $T_0,\dots,T_{N-1}$
  of degree $r-N$ and weight $d-r$ have image in $\ltwentynineRn$ of dimension at most the socle-word count. The proof
  is a fusion-product recursion for symmetric polynomials (\leanpath{TwoSeriesBoxFusion}).
\item \emph{Box inclusion}: the membership space lies in the span of monomials in $T_0,\dots,T_{N-1}$ (the ``box''). It is
  proved directly in the slots with $N\le1$, with $T$-degree $\le n+1$, and with weight below $2N+k$; the remaining
  slots are exactly the jet-ideal Poincar\'e lemma below (\leanpath{TwoSeriesHighReduce}, \leanpath{TwoSeriesHighEndpoint},
  \leanpath{TwoSeriesHighDegSucc}, \leanpath{TwoSeriesPoincareZero}).
\end{enumerate}

\begin{definition}[$T$-monomial spans; the box]\label{l29:def:box}
For a ring $R$ with elements $t_j$: $\mathrm{tMon}(k,e)$ is the span of the products $t_{m_1}\cdots t_{m_k}$ with
$\sum_i(m_i+1)=e$; $\mathrm{box}(N,k,e)$ is the span of those products with every $m_i<N$.
\end{definition}

\noindent Files \leanpath{VOA/\allowbreak Examples/\allowbreak TwoSeriesMembershipLine.lean} and
\leanpath{VOA/\allowbreak Examples/\allowbreak TwoSeriesMembershipEndpoint.lean}.
\begin{leancode}
def tMonW (t : ℕ → R) (k : ℕ) (e : ℤ) : Submodule ℂ R :=
  Submodule.span ℂ
    {x | ∃ l : List ℕ, l.length = k ∧ (((l.map fun m => m + 1).sum : ℕ) : ℤ) = e ∧
      (l.map t).prod = x}

def boxW (t : ℕ → R) (N k : ℕ) (e : ℤ) : Submodule ℂ R :=
  Submodule.span ℂ {x | ∃ l : List ℕ, l.length = k ∧ (((l.map fun m => m + 1).sum : ℕ) : ℤ) = e ∧
    (∀ m ∈ l, m < N) ∧ (l.map t).prod = x}
\end{leancode}

\begin{dictionary}
A list $l=(m_1,\dots,m_k)$ encodes the monomial $t_{m_1}\cdots t_{m_k}$; its weight is $\sum(m_i+1)$ (the jet $T_m$ has
weight $m+1$ in these files). \lean{R} is \lean{TPoly} (with $t=X$) or \lean{Rn n} (with $t=t_\bullet$).
\end{dictionary}

\begin{physics}
The box is the space of classical stress-tensor configurations that use only the first $N$ Fourier modes; the Poincar\'e
lemma says a configuration whose high-mode derivatives are null is, up to nulls, of this low-mode form.
\end{physics}

\begin{status}
Definitions.
\end{status}

\begin{definition}[The jet-ideal Poincar\'e lemma, and its core slots]\label{l29:def:poincare}
$\mathrm{JP}_n(N,k,E)$: every $F\in\C[T]$ of degree $k$ and weight $E$ (a combination of monomials
$T_{m_1}\cdots T_{m_k}$ with $\sum(m_i+1)=E$) such that $\partial F/\partial T_M$ lies in the jet ideal of $T^n$ for every
$M\ge N$ is, modulo the jet ideal, in $\mathrm{box}(N,k,E)$. The \emph{core} $\mathrm{JPCore}_n$ is $\mathrm{JP}_n(N,k,E)$
for all $N\ge2$, $k\ge n+2$, $2N+k\le E$ and $2N+2k<E+n+2$.
\end{definition}

\noindent Files \leanpath{VOA/\allowbreak Examples/\allowbreak TwoSeriesHighJet.lean} and
\leanpath{VOA/\allowbreak Examples/\allowbreak TwoSeriesPoincareZero.lean}.
\begin{leancode}
def JetPoincare (n N k E : ℕ) : Prop :=
  ∀ F ∈ tMonW TPoly MvPolynomial.X k (E : ℤ), (∀ M, N ≤ M → pderiv M F ∈ jetIdealT n) →
    Ideal.Quotient.mk (jetIdealT n) F ∈ boxW (tn n) N k (E : ℤ)

def JetPoincareCoreT0 (n : ℕ) : Prop :=
  ∀ N k E : ℕ, 2 ≤ N → n + 2 ≤ k → 2 * N + k ≤ E → 2 * N + 2 * k < E + n + 2 →
    JetPoincare n N k E
\end{leancode}

\begin{dictionary}
\lean{pderiv M F} is $\partial F/\partial T_M$; \lean{Ideal.Quotient.mk} the projection $\C[T]\to\ltwentynineRn$.
\lean{JetPoincare} lives in the namespace \ltwentynineln{VOA.TwoSeries.StrictModel.HighJet}, \lean{JetPoincareCoreT0} in
\ltwentynineln{VOA.TwoSeries.StrictModel.PoincareZero}. Our $\mathrm{JPCore}_n$ is \lean{JetPoincareCoreT0 n}; the library also
has a predicate \lean{JetPoincareCore n} (\leanpath{TwoSeriesHighDegSucc.lean}) without the last constraint, a different and
stronger statement, converted by \lean{jetPoincareCore_of_T0}. The slots outside the core are proved: $N=0$ (Euler identity,
\lean{jetPoincare_zero}), $N=1$ (\lean{jetPoincare_one}), $E<2N+k$ (\lean{jetPoincare_of_wt_lt}: no monomial has two
indices $\ge N$), $k=n+1$ (\lean{jetPoincare_degSucc}, by an explicit primitive), and $E+n+2\le2N+2k$
(\lean{jetPoincare_of_T0}: every monomial with two high indices contains $T_0^n=(T^n)_0$). The box inclusions with
$T$-degree $k\le n$ never use the lemma; they are proved through a separate ``test ring''
(\lean{boxInclusion_low}, \lean{boxInclusion_degN}). This is why the core starts at $k\ge n+2$ and $N\ge2$.
\end{dictionary}

\begin{physics}
A ``Poincar\'e lemma'' for the arc space of the fat point $T^n=0$: a classical configuration that is closed in the high
directions is exact up to the constraints. It is the classical statement that the null-vector constraints of
$M(2,2n+1)$ do not generate unexpected bound states.
\end{physics}

\begin{status}
Definitions (predicates). $\mathrm{JPCore}_n$ is proved for every $n\ge1$ (Theorem~\ref{l29:thm:poincarecore}).
\end{status}

\begin{theorem}[Transfer from detection and the Poincar\'e core]\label{l29:thm:strictpoincare}
For every $n\ge1$: $\ltwentynineGU(L_{k_n})$, detection and $\mathrm{JPCore}_n$ imply $\ltwentynineGU(M(2,2n+3))$.
\end{theorem}

\noindent File \leanpath{VOA/\allowbreak Examples/\allowbreak TwoSeriesPoincareZero.lean}.
\begin{leancode}
theorem twoSeries_gradedUnitary_of_strictModel_jetPoincareCoreT0 (n : ℕ) (hn : 1 ≤ n)
    (hA : GradedUnitary (affineForm (twoSeriesLevels n)) (affineRealStructure (twoSeriesLevels n))
      (affineFlag (twoSeriesLevels n)))
    (det : TwoSeriesModelDetects n) (h : JetPoincareCoreT0 n) :
    VirGradedUnitarity.VirGradedUnitary (twoSeriesLevels n).c
      (targetRealStructure (twoSeriesLevels n)) :=
  twoSeries_gradedUnitary_of_strictModel_jetPoincareCore n hn hA det (jetPoincareCore_of_T0 h)
\end{leancode}

\begin{dictionary}
Full name \ltwentynineln{VOA.TwoSeries.StrictModel.twoSeries_gradedUnitary_of_strictModel_jetPoincareCoreT0}. The model socle
data are assembled from transversality, box count, and box inclusion (proved slots plus \lean{h}); the character bound
from Theorem~\ref{l29:thm:modelchar}.
\end{dictionary}

\begin{proof}[Proof on paper]
Box inclusion in every slot follows from the proved slots and $\mathrm{JPCore}_n$: given a membership element $\tau$,
lift it to $F\in\C[T]$; the model relations show that $\partial F/\partial T_M$ is in the jet ideal for $M\ge N$
(\leanpath{TwoSeriesHighReduce}); the Poincar\'e lemma puts $\tau$ in the box. With the box count this bounds the
membership spaces by the socle-word count, which is the rank half; transversality is proved. Apply
Theorem~\ref{l29:thm:strictdet}.
\end{proof}

\begin{physics}
Everything about the free-field model reduces to one classical fact about the arc space of $T^n=0$, plus the
statement that the affine nulls are invisible on the slice.
\end{physics}

\begin{status}
Proved, \texttt{\#print axioms}-clean (re-checked), \textbf{conditional} on \lean{hA}, \lean{det}, \lean{h}.
\end{status}

%======================================================================
\section{Detection: the Drinfeld--Sokolov slice map}\label{l29:sec:ds}
%======================================================================

Detection (Definition~\ref{l29:def:detect}) asks that the model symbol of every null vector of $V^{k_n}$ vanish at the
base point. At $n=1$ ($R_1=\C$) it followed from the level-one vacuum functional. For $n\ge2$ the values lie in the
non-reduced ring $\ltwentynineRn$, and the vacuum functional only sees their constant part (\lean{evR_aug_eq_zero}).
The idea that closes the gap is that evaluating at the base point $(u,v)=(0,1)$ is the classical shadow of
\emph{Drinfeld--Sokolov reduction}: restricting to the Kostant slice $f+\C e$. So one builds a quantum map from
$V^{k}$ to the Virasoro vacuum module at the Drinfeld--Sokolov central charge, shows it sends nulls to nulls, and
composes with the leading-symbol map of the Virasoro side, whose target is exactly $\ltwentynineRn$.

\subsection{The Virasoro side: leading symbols into $\ltwentynineRn$}

\begin{definition}[The Virasoro symbol property]\label{l29:def:virsymbol}
For $c\in\C$, a ring $R$ and $t:\N\to R$: there are linear maps $\sigma_b:V(c,0)\to R$ ($b\in\N$) that kill the
maximal ideal, kill every creation word of length $<b$, and send every creation word
$L_{-j_1-2}\cdots L_{-j_b-2}\ltwentyninevac$ of length $b$ (in any order) to $t_{j_1}\cdots t_{j_b}$.
\end{definition}

\noindent File \leanpath{VOA/\allowbreak Examples/\allowbreak TwoSeriesDSVirSymbol.lean}.
\begin{leancode}
def vLen (b : ℕ) : Submodule ℂ (VirUniv c) :=
  Submodule.span ℂ {x | ∃ w : List ℕ, w.length ≤ b ∧ negWordVir c w = x}

def vBelow (b : ℕ) : Submodule ℂ (VirUniv c) :=
  Submodule.span ℂ {x | ∃ w : List ℕ, w.length < b ∧ negWordVir c w = x}

def VirSymbol : Prop :=
  ∃ σ : ℕ → (VirUniv c →ₗ[ℂ] R),
    (∀ b, ∀ v ∈ maxIdealVir c, σ b v = 0) ∧
    (∀ b, ∀ v ∈ vBelow c b, σ b v = 0) ∧
    (∀ w : List ℕ, σ w.length (negWordVir c w) = (w.map t).prod)
\end{leancode}

\begin{dictionary}
\lean{negWordVir c w} is the creation word $L_{-w_1-2}\cdots L_{-w_\ell-2}\ltwentyninevac$ for the list $w$ (no
ordering assumed). \lean{(w.map t).prod} is $t_{w_1}\cdots t_{w_\ell}$. The words in \lean{vLen}, \lean{vBelow} have
any weight and any order. Same as the textbook definition.
\end{dictionary}

\begin{physics}
A leading-symbol (classical-limit) map for the Virasoro chiral algebra with values in a ring where $T^{n}$'s jets vanish:
the quantum states of the minimal model project onto classical stress-tensor configurations.
\end{physics}

\begin{status}
Definition.
\end{status}

\begin{theorem}[Leading symbols of $M(2,2m+3)$]\label{l29:thm:virsymbol}
For every $m\in\N$, the Virasoro symbol property holds for $c=c_{2,2m+3}$ with values in $R_{m+1}$.
\end{theorem}

\noindent File \leanpath{VOA/\allowbreak Examples/\allowbreak TwoSeriesDSVirSymbol.lean}.
\begin{leancode}
theorem virSymbol_twoSeries (m : ℕ) :
    VirSymbol (twoSeriesLevels m).c (Rn (m + 1)) (tn (m + 1)) :=
  virSymbol_of _ _ _ m (twoSeries_hasJetRelations_all m) (gordonLowerBound_twoSeries m)
    (jet_rel_Rn m)
\end{leancode}

\begin{dictionary}
Full name \ltwentynineln{VOA.TwoSeries.DSSlice.virSymbol_twoSeries}. Note the indices: $(\lean{twoSeriesLevels m}).c=c_{2,2m+3}$
and the ring is $R_{m+1}$, the arc algebra of $\C[T]/(T^{m+1})$. The proof uses the jet relations
(Theorem~\ref{l29:thm:jetrel}), the total Gordon lower bound (Theorem~\ref{l29:thm:character}) and the jet relations of
$R_{m+1}$ (\lean{jet_rel_Rn}).

\textbf{Scope.} This is a PBW-length leading-symbol statement into the \emph{presented} ring $R_{m+1}$; it does not
mention Li's filtration or $V/C_2V$. Intrinsic classical freeness of $M(2,2m+3)$ is a separate theorem
(\ltwentynineln{VOA.VirasoroClassicalFreeness.classicallyFree}, Chapter~\ref{ch:L20}) that uses this one as its injectivity
input. The docstring of \leanpath{TwoSeriesDSSlice.lean} still calls this input ``classical freeness of $M(2,2n+1)$''
(warning on p.~\pageref{l29:warn:stale}).
\end{dictionary}

\begin{proof}[Proof on paper]
The Gordon-good sorted words form a basis of $L(c,0)$ (spanning plus the character). Define $\sigma_b$ on this basis:
a Gordon-good word of length $b$ goes to the product of its jets, every other basis word to $0$. It remains to see
that $\sigma_b$ has the product property on \emph{every} word. Straighten simultaneously: a pair (Virasoro word, its
$T$-monomial) is rewritten, by the same jet relation applied on both sides (in $V(c,0)$ the null vector of
Theorem~\ref{l29:thm:jetrel}, in $R_{m+1}$ the jets of $T^{m+1}$), into a combination of pairs of Gordon-good words of
the same length, modulo (shorter words $+$ maximal ideal)$\times0$ (\lean{pair_mem_pairGood}). Both sides of the product
property are linear and compatible with this rewriting.
\end{proof}

\begin{physics}
The null-vector relation of $M(2,2m+3)$ and the classical constraint $T^{m+1}=0$ have the same leading term, so the
quantum straightening and the classical one run in parallel.
\end{physics}

\begin{status}
Proved, \textbf{unconditional}. Axioms standard (re-checked).
\end{status}

\subsection{The slice map and its Arakawa property}

\begin{definition}[A Drinfeld--Sokolov slice map; the slice Arakawa property]\label{l29:def:dsslice}
For a central charge $c$ and a level datum $L$, a \emph{slice map} is a linear map $W:V^{k_A}(\ltwentyninesl)\to
\bigoplus_{a\in\Z}V(c,0)$ (graded by the Rees parameter $a$) with: $W(\ltwentyninevac)=\ltwentyninevac$ in degree $0$;
the Whittaker relations $W\circ f_{-1}=-(\text{shift by one})\circ W$ and $W\circ f_{-j-2}=0$; the filtration property
that a corrected PBW word of length $\ell$ has degree-$a$ image of Virasoro length $\le\ell-a$, one less if it starts with
an $e$- or $h$-letter; and the Sugawara letter acting as $L_{-m-2}$ up to that lower piece. The \emph{slice Arakawa
property} of $W$: $W$ maps the maximal ideal of $V^{k_A}$ into the maximal ideal of $V(c,0)$, in every degree.
\end{definition}

\noindent File \leanpath{VOA/\allowbreak Examples/\allowbreak TwoSeriesDSSlice.lean}.
\begin{leancode}
structure DSSliceMap where
  /-- The map, with its Rees-parameter grading. -/
  W : WakModule L.kA →ₗ[ℂ] (ℤ →₀ VirUniv c)
  vacuum : W (correctedWordState L.kA L.aNoncritical []) = Finsupp.single 0 (vacVec c)
  f_one : ∀ (y : WakModule L.kA) (a : ℤ), W (J L.kA 1 (-1) y) a = -(W y (a - 1))
  f_high : ∀ (j : ℕ) (y : WakModule L.kA), W (J L.kA 1 (-(j : ℤ) - 2) y) = 0
  filtered : ∀ (w : List (Generator × ℕ)) (a : ℤ),
    W (correctedWordState L.kA L.aNoncritical w) a ∈ vLenZ c ((w.length : ℤ) - a)
  current_lower : ∀ (b : Fin 3), b ≠ 1 → ∀ (m : ℕ) (w : List (Generator × ℕ)) (a : ℤ),
    W (correctedWordState L.kA L.aNoncritical ((Generator.current b, m) :: w)) a ∈
      vLenZ c ((w.length : ℤ) - a)
  sugawara : ∀ (m : ℕ) (w : List (Generator × ℕ)) (a : ℤ),
    W (correctedWordState L.kA L.aNoncritical ((Generator.sugawara, m) :: w)) a -
      Lop c (-(m : ℤ) - 2) (W (correctedWordState L.kA L.aNoncritical w) a) ∈
        vLenZ c ((w.length : ℤ) - a)

def SliceArakawa (S : DSSliceMap c L) : Prop :=
  ∀ y ∈ ImaxL L, ∀ a : ℤ, S.W y a ∈ maxIdealVir c
\end{leancode}

\begin{dictionary}
\lean{ℤ →₀ VirUniv c} is the space of finitely supported functions $\Z\to V(c,0)$, i.e.\ $\bigoplus_{a\in\Z}V(c,0)$;
\lean{W y a} is the degree-$a$ component, \lean{Finsupp.single 0 v} the element $v$ in degree $0$. \lean{J L.kA 1 (-1)}
is $f_{-1}$ (\lean{Fin 3}: $1=f$), \lean{J L.kA 1 (-j-2)} is $f_{-j-2}$. \lean{correctedWordState k hk w} is the
corrected PBW word of $V^k$ for the list of letters $w$ (Chapter~\ref{ch:L28}). \lean{vLenZ c b} is the span of
Virasoro words of length $\le b$ (and $0$ for $b<0$). \lean{ImaxL L} is the maximal ideal of $V^{k_A}$. All the fields
of \lean{DSSliceMap} are statements about the \emph{universal} objects; only \lean{SliceArakawa} mentions the maximal
ideals.
\end{dictionary}

\begin{physics}
$W$ is the quantum Hamiltonian reduction $V^k(\ltwentyninesl)\to\mathrm{Vir}^{c_{\mathrm{DS}}(k)}$ made concrete on the
vacuum module: impose $f(z)=1$ (the Whittaker character, $f_{-1}\mapsto-1$, higher $f$-modes $\mapsto0$), and read the
result in the Virasoro algebra generated by the improved Sugawara field. The slice Arakawa property says the reduction
of the simple affine VOA is a quotient of the simple Virasoro VOA's universal cover by something containing its nulls:
the ``nulls go to nulls'' half of Arakawa's theorem $H_{DS}(L_k(\ltwentyninesl))\cong L(c_{\mathrm{DS}}(k),0)$.
\end{physics}

\begin{status}
Definitions. Inhabited for every $c$ and $L$ (Definition~\ref{l29:def:dsslicemap}); the Arakawa property is proved
at every level datum whose affine level is not a natural number (Theorem~\ref{l29:thm:slicearakawa}).
\end{status}

\begin{theorem}[Detection from a slice map]\label{l29:thm:detectreduce}
For every $m$: a slice map $V^{k_{m+1}}\to\bigoplus_aV(c_{2,2m+3},0)$ with the slice Arakawa property gives detection at
member $n=m+1$.
\end{theorem}

\noindent File \leanpath{VOA/\allowbreak Examples/\allowbreak TwoSeriesDSSlice.lean}.
\begin{leancode}
theorem twoSeriesModelDetects_of_dsSlice (m : ℕ)
    (S : DSSliceMap (twoSeriesLevels m).c (twoSeriesLevels (m + 1))) (hS : SliceArakawa S) :
    TwoSeriesModelDetects (m + 1) :=
  modelDetects_of_dsSlice S hS (virSymbol_twoSeries m)
\end{leancode}

\begin{dictionary}
The Virasoro charge is $c_{2,2m+3}=c_{\mathrm{DS}}(k_{m+1})$, one step below the target $c_{2,2m+5}$ of the member
$n=m+1$ (Example~\ref{l29:ex:levels}, last column). The Virasoro symbol input is Theorem~\ref{l29:thm:virsymbol}.
\end{dictionary}

\begin{proof}[Proof on paper]
Put $\Lambda_r=\sum_a\sigma_{r-a}\circ W_a:V^k\to R_{m+1}$. It kills the maximal ideal (Arakawa property, then $\sigma$
kills the Virasoro maximal ideal), kills corrected words of length $<r$ (filtration property), and on a word of length
$r$ equals the model symbol evaluated at the base point (Whittaker relations for the $f$-letters, which give
$v_0\mapsto1$; the lower-term properties for $e$, $h$; the Sugawara property for the $L$-letters, whose symbol is $t_j$
at the base point). Hence the base-point evaluation of the symbol of a null vector is $\Lambda_r$ of it, which is $0$.
\end{proof}

\begin{physics}
Evaluating classical symbols on the Kostant slice is the classical limit of Drinfeld--Sokolov reduction; detection is
the classical limit of ``nulls reduce to nulls''.
\end{physics}

\begin{status}
Proved, \texttt{\#print axioms}-clean (re-checked), conditional on \lean{S} and \lean{hS} (both supplied below).
\end{status}

\begin{definition}[The constructed slice map]\label{l29:def:dsslicemap}
For every $c$ and non-critical $k$: $W_0:V^k\to\C[X_0,X_1,\dots]$ is the Whittaker coinvariant map (on the
$(f,h,e)$-ordered PBW basis, $f^\lambda h^\nu e^\mu\ltwentyninevac\mapsto\chi(f^\lambda)[\nu=0]X^\mu$); $\psi$ is the
inverse of the isomorphism $V(c,0)\to\C[X]$ sending a Virasoro word to the corresponding word of descended Sugawara
operators applied to $1$; $W_1=\psi\circ W_0$; and $W$ places the image of a word of charge $\#f-\#e=a$ in degree $a$.
\end{definition}

\noindent Files \leanpath{VOA/\allowbreak Examples/\allowbreak TwoSeriesSliceMap.lean} and
\leanpath{VOA/\allowbreak Examples/\allowbreak TwoSeriesArakawaPositive.lean}.
\begin{leancode}
def W1 : WakModule k →ₗ[ℂ] VirUniv c := (psi c k hk).toLinearMap.comp (W0 k hk)

def Wslice : WakModule k →ₗ[ℂ] (ℤ →₀ VirUniv c) := chargeLift k hk (W1 c k hk)

def dsSliceMap (L : CosetLevels) : DSSliceMap c L where
  W := Wslice c L.kA L.aNoncritical
  vacuum := Wslice_vacuum c L.kA L.aNoncritical
  f_one := Wslice_f_one c L.kA L.aNoncritical
  f_high := Wslice_f_high c L.kA L.aNoncritical
  filtered := Wslice_filtered c L.kA L.aNoncritical
  current_lower := Wslice_current_lower c L.kA L.aNoncritical
  sugawara := fun m w a => by
    rw [Wslice_sugawara c L.kA L.aNoncritical m w a]
    exact Submodule.zero_mem _

def cDS (k : ℂ) : ℂ := 1 - 6 * (k + 1) ^ 2 / (k + 2)
\end{leancode}

\begin{dictionary}
\lean{psi c k hk} is a linear \emph{equivalence} (\lean{phi_injective}, \lean{phi_surjective} in
\leanpath{TwoSeriesSliceMapVir.lean}); \lean{.toLinearMap.comp} is composition. \lean{chargeLift} implements the Rees
grading. Every field of the structure is a proved lemma; the Sugawara property holds \emph{exactly} (difference $0$).
\lean{cDS k} is the Drinfeld--Sokolov central charge $1-6(k+1)^2/(k+2)$. The construction works for every $c$; only the
Arakawa property needs $c=c_{\mathrm{DS}}(k)$.
\end{dictionary}

\begin{physics}
The Whittaker coinvariants of the vacuum module are a polynomial ring in the $e$-modes, and the improved Sugawara
modes act on them triangularly, which identifies them with the Virasoro vacuum module: the vacuum-module version of
Kostant's theorem on Whittaker vectors.
\end{physics}

\begin{status}
Construction, proved; \texttt{\#print axioms} of \lean{dsSliceMap} standard (re-checked).
\end{status}

\begin{theorem}[The slice form of Arakawa's theorem, unconditionally]\label{l29:thm:slicearakawa}
For every level datum $L$ (Definition~\ref{l29:def:cosetlevels}) whose affine level $k=k_A$ is not a natural number, the slice map
at $c=c_{\mathrm{DS}}(k)$ sends the maximal ideal of $V^k(\ltwentyninesl)$ into the maximal ideal of $V(c_{\mathrm{DS}}(k),0)$, in every
Rees degree. For an arbitrary non-critical $k\notin\N$ (no level datum) the same holds for the ungraded map $W_1$. In particular, for every $m$, the slice map of
the member $n=m+1$ at $c=c_{2,2m+3}$ has the slice Arakawa property.
\end{theorem}

\noindent File \leanpath{VOA/\allowbreak Examples/\allowbreak TwoSeriesArakawaUnconditional.lean}.
\begin{leancode}
theorem W1_maxIdeal (hkN : ∀ i : ℕ, k ≠ i) {y : WakModule k} (hy : y ∈ Imax k hk) :
    W1 (cDS k) k hk y ∈ maxIdealVir (cDS k) := by

theorem sliceArakawa_unconditional (L : CosetLevels) (hkN : ∀ i : ℕ, L.kA ≠ i) :
    SliceArakawa (dsSliceMap (cDS L.kA) L) := by

theorem twoSeries_kA_not_nat (m : ℕ) (i : ℕ) : (twoSeriesLevels (m + 1)).kA ≠ i := by

theorem twoSeries_sliceArakawa (m : ℕ) :
    SliceArakawa (dsSliceMap (twoSeriesLevels m).c (twoSeriesLevels (m + 1))) := by
\end{leancode}

\begin{dictionary}
\lean{hkN : ∀ i : ℕ, k ≠ i} says $k\notin\{0,1,2,\dots\}$ (as complex numbers). \lean{Imax k hk} is the maximal graded
ideal of $V^k$. For generic $k$ that ideal is $0$ and the statement is trivial; it has content exactly at levels where
$V^k$ is reducible, such as $k_n$. \lean{twoSeries_kA_not_nat}: $k_{m+1}=-2+2/(2m+3)\in(-2,0)$ is not a natural number.
\lean{twoSeries_sliceArakawa} rewrites $c_{2,2m+3}=c_{\mathrm{DS}}(k_{m+1})$ (\lean{twoSeries_c_eq_cDS}) and applies the
general theorem. No other hypothesis.

\textbf{Same, stronger or weaker than Arakawa?} Weaker, deliberately: Arakawa's theorem identifies $H_{DS}(L_k)$ with
$L(c_{\mathrm{DS}}(k),0)$ (including non-vanishing and surjectivity statements about cohomology). Here only the
``descends'' half is proved, for this explicit map on the vacuum module, and only it is used. No result of Arakawa's is
assumed.
\end{dictionary}

\begin{proof}[Proof on paper]
Let $I$ be the maximal ideal of $V^k$ and $S=W_1(I)\subseteq V(c,0)$, $c=c_{\mathrm{DS}}(k)$. (1) $S$ is stable under the
Virasoro creation modes, because $W_1\circ L^{\mathrm{sug}}_{n}=L_n\circ W_1$ for $n\le-1$ and $I$ is stable under the
Sugawara modes. (2) $S$ is stable under $L_1,L_2$, hence under all $L_{n\ge1}$: there are explicit lifts
$Q_1=L_1+h_1+kf_0+\frac12f_0h_0$ and $Q_2$ (a similar quadratic expression) with $W_1\circ Q_i=L_i\circ W_1$, and
$Q_iI\subseteq I$ since $I$ is an ideal. (3) No vector of $S$ has a vacuum component: the vacuum coefficient of $W_1(y)$
only sees the PBW words of Drinfeld--Sokolov weight zero, the powers $f_{-1}^N\ltwentyninevac$; a polynomial projector in
the commuting operators $L_0,h_0$ keeps $y$ in $I$ and isolates such a component, and
$e_1^Nf_{-1}^N\ltwentyninevac=N!\,k(k-1)\cdots(k-N+1)\ltwentyninevac\ne0$ for $k\notin\N$, while
$\ltwentyninevac\notin I$; so that component is $0$ (\lean{eps_maxIdeal}). (4) A subspace stable under all modes with no
vacuum component lies in the radical of the Shapovalov form, which is the maximal ideal of $V(c,0)$ (Li). The graded
components of $W=$ \lean{chargeLift} $W_1$ are separated by a Lagrange projector in $h_0$.
\end{proof}

\begin{physics}
``Under quantum Hamiltonian reduction, null states of $L_k(\ltwentyninesl)$ go to null states of the Virasoro algebra.''
For $k_n$ this is the statement that the reduction of the $(A_1,D_{2n+1})$ chiral algebra by its $SU(2)$ lands in
$M(2,2n+1)$, the chiral algebra of $(A_1,A_{2n-2})$.
\end{physics}

\begin{status}
Proved, \textbf{unconditional} for \lean{twoSeries_sliceArakawa}. The general \lean{sliceArakawa_unconditional} is stated for a
level datum $L$ (a discrete set of levels, since $c(k)$ must be a boundary charge) with $k_A\notin\N$; only the ungraded
\lean{W1_maxIdeal} is stated for every non-critical $k\notin\N$. Axioms standard (re-checked: \lean{sliceArakawa_unconditional},
\lean{twoSeries_sliceArakawa}).
\end{status}

\begin{theorem}[Detection holds for every member]\label{l29:thm:detect}
For every $n\ge1$, detection holds at member $n$.
\end{theorem}

\noindent This theorem is not a separate declaration of the library: it is the term
\lean{twoSeriesModelDetects_of_dsSlice m (dsSliceMap _ _) (twoSeries_sliceArakawa m)} with $n=m+1$. For this edition the composite
was checked as a standalone statement in a scratch file importing \lean{VOA.Examples.TwoSeriesGenIncl} (with the namespaces
\lean{VOA.TwoSeries}, \lean{.DSSlice}, \lean{.SliceMap}, \lean{.StrictModel} opened):
\begin{plaincode}
theorem detect_check (m : ℕ) : TwoSeriesModelDetects (m + 1) :=
  twoSeriesModelDetects_of_dsSlice m (dsSliceMap _ _) (twoSeries_sliceArakawa m)

'detect_check' depends on axioms: [propext, Classical.choice, Quot.sound]
\end{plaincode}
\noindent It is used in this form inside Theorem~\ref{l29:thm:slicegu} (file \leanpath{VOA/\allowbreak Examples/\allowbreak TwoSeriesDSSlice.lean},
through \ltwentynineln{twoSeries_gradedUnitary_of_dsSlice}). The quoted components are Theorems~\ref{l29:thm:detectreduce} and
\ref{l29:thm:slicearakawa} and Definition~\ref{l29:def:dsslicemap}.

\begin{physics}
The free-field model of the classical limit of $L_{k_n}$ is faithful on the null vectors.
\end{physics}

\begin{status}
Proved as a composite term, not as a named declaration; unconditional; the composite and each component re-checked with
standard axioms. The same holds for ``model socle data for every $n\ge1$'' and for \lean{TwoSeriesSocleSpectrum} at every $n\ge1$:
they are obtained inside the proof of the headline, not as standalone declarations.
\end{status}

\begin{theorem}[Transfer conditional on the Poincar\'e core]\label{l29:thm:slicegu}
For every $m$: $\ltwentynineGU(L_{-2+2/(2m+3)})\wedge\mathrm{JPCore}_{m+1}\Rightarrow\ltwentynineGU(M(2,2m+5))$.
\end{theorem}

\noindent File \leanpath{VOA/\allowbreak Examples/\allowbreak TwoSeriesArakawaUnconditional.lean}.
\begin{leancode}
theorem twoSeries_gradedUnitary_of_affine (m : ℕ)
    (hA : GradedUnitary (affineForm (twoSeriesLevels (m + 1)))
      (affineRealStructure (twoSeriesLevels (m + 1))) (affineFlag (twoSeriesLevels (m + 1))))
    (h : PoincareZero.JetPoincareCoreT0 (m + 1)) :
    VirGradedUnitary (twoSeriesLevels (m + 1)).c (targetRealStructure (twoSeriesLevels (m + 1))) :=
  twoSeries_gradedUnitary_of_sliceArakawa m hA (twoSeries_sliceArakawa m) h
\end{leancode}

\begin{dictionary}
Full name \ltwentynineln{VOA.TwoSeries.SliceMap.twoSeries_gradedUnitary_of_affine}: the \emph{second} declaration with this short
name (warning on p.~\pageref{l29:warn:clash}). Its hypotheses are \lean{hA} and \lean{h}; detection has been discharged.
\end{dictionary}

\begin{physics}
At this point the transfer depends on graded unitarity of the $(A_1,D_{2n+1})$ chiral algebra and one classical fact
about the arc space of $T^n=0$.
\end{physics}

\begin{status}
Proved, \texttt{\#print axioms}-clean (re-checked), \textbf{conditional} on \lean{hA} and \lean{h}.
\end{status}

%======================================================================
\section{The jet-ideal Poincar\'e lemma and the \texorpdfstring{$n$}{n}-equal arrangement}\label{l29:sec:arr}
%======================================================================

The last input, $\mathrm{JPCore}_n$, is pure commutative algebra. It is proved by duality: a linear functional on
$\C[T]$ of weight $E$ and degree $k$ is the same as a symmetric polynomial $f(x_0,\dots,x_{k-1})$ (``polarisation'':
$T_{m_1}\cdots T_{m_k}\leftrightarrow$ the symmetrised monomial $x^{m}$); a functional that kills the jet ideal of
$T^n$ corresponds to an $f$ that vanishes whenever $n$ of the variables coincide; and a functional that kills the box
corresponds to an $f$ in the ideal $(x_0^N,\dots,x_{k-1}^N)$. The Poincar\'e lemma then becomes an ideal-membership
statement for the \emph{$n$-equal arrangement}, the union of the diagonals $\{x_{s_1}=\cdots=x_{s_n}\}$.

\subsection{The arrangement inclusion}

\begin{definition}[Diagonals and the arrangement inclusion]\label{l29:def:arrangement}
In $\C[x_0,\dots,x_{k-1}]$: for $S\subseteq\{0,\dots,k-1\}$, $Y$ \emph{vanishes on the diagonal of $S$} if substituting
$x_s\mapsto x_{s_0}$ for all $s\in S$ (any $s_0\in S$) gives $0$. $L_N=(x_0^N,\dots,x_{k-1}^N)$. $W_{N,n}$ is the span of
the products $x_l^NY$ with $Y$ vanishing on every $n$-element diagonal not containing $l$. The \emph{arrangement
inclusion} $\mathrm{AI}_n$: every symmetric $f\in L_N$ vanishing on all $n$-element diagonals lies in $W_{N,n}$ (for all
$k,N$). Its \emph{narrowed} form $\mathrm{SAI}_n$ asks this only for $f$ homogeneous of degree $d$ in the window
$N\ge2$, $k\ge n+2$, $2N\le d$, $2N+k<d+n+2$.
\end{definition}

\noindent Files \leanpath{VOA/\allowbreak Examples/\allowbreak TwoSeriesJetPoincareVdm.lean},
\leanpath{VOA/\allowbreak Examples/\allowbreak TwoSeriesJetPoincareArr.lean} and
\leanpath{VOA/\allowbreak Examples/\allowbreak TwoSeriesArrangementSym.lean}.
\begin{leancode}
abbrev XP (k : ℕ) : Type := MvPolynomial (Fin k) ℂ

def Lid (k N : ℕ) : Ideal (XP k) := Ideal.span (Set.range fun i : Fin k => (X i ^ N : XP k))

def collapseTo (s₀ : Fin k) (S : Finset (Fin k)) (i : Fin k) : Fin k := if i ∈ S then s₀ else i

def VanishOn (S : Finset (Fin k)) (Y : XP k) : Prop := ∀ s₀ ∈ S, rename (collapseTo s₀ S) Y = 0

def Wn (k N n : ℕ) : Submodule ℂ (XP k) :=
  Submodule.span ℂ {Z | ∃ (l : Fin k) (Y : XP k),
    (∀ S : Finset (Fin k), #S = n → l ∉ S → VanishOn S Y) ∧ Z = X l ^ N * Y}

def ArrangementInclusion (n : ℕ) : Prop :=
  ∀ (k N : ℕ) (f : XP k), (∀ σ : Equiv.Perm (Fin k), rename σ f = f) → f ∈ Lid k N →
    (∀ S : Finset (Fin k), #S = n → VanishOn S f) → f ∈ Wn k N n

def SymArrangementInclusion (n : ℕ) : Prop :=
  ∀ (N k d : ℕ), 2 ≤ N → n + 2 ≤ k → 2 * N ≤ d → 2 * N + k < d + n + 2 →
    ∀ f : XP k, f.IsHomogeneous d → (∀ σ : Equiv.Perm (Fin k), rename σ f = f) →
      f ∈ Lid k N → (∀ S : Finset (Fin k), #S = n → VanishOn S f) → f ∈ Wn k N n
\end{leancode}

\begin{dictionary}
\lean{rename g Y} substitutes $x_i\mapsto x_{g(i)}$. So \lean{rename (collapseTo s₀ S) Y = 0} is $Y|_{x_S=x_{s_0}}=0$;
quantifying over $s_0\in S$ is harmless (all choices give the same substitution up to renaming). \lean{#S = n} is
$|S|=n$. \lean{rename σ f = f} for all permutations $\sigma$ says $f$ is symmetric. \lean{f.IsHomogeneous d}: all
monomials of $f$ have total degree $d$. \lean{Wn k N n} is $W_{N,n}$; note $W_{N,n}\subseteq L_N$ and every element of
$W_{N,n}$ vanishes on all $n$-element diagonals, so $\mathrm{AI}_n$ is the nontrivial inclusion
``$L_N\cap(\text{vanishing})\cap(\text{symmetric})\subseteq W_{N,n}$''. \ltwentynineln{symArrangementInclusion_of_arrangementInclusion}
shows $\mathrm{AI}_n\Rightarrow\mathrm{SAI}_n$.
\end{dictionary}

\begin{physics}
The vanishing on $n$-fold diagonals is the functional (Feigin--Stoyanovsky) form of the exclusion rule ``at most $n-1$
quanta at the same point'', i.e.\ of $T^n=0$. The inclusion says that the exclusion rule and the mode cutoff
$x^N$ interact in the simplest possible way.
\end{physics}

\begin{status}
Definitions (predicates). Both are proved for every $n\ge1$ (Theorem~\ref{l29:thm:geninclusion}).
\end{status}

\begin{theorem}[The Poincar\'e core from the arrangement inclusion]\label{l29:thm:poincarecore}
For every $n\ge1$, $\mathrm{SAI}_n\Rightarrow\mathrm{JPCore}_n$.
\end{theorem}

\noindent File \leanpath{VOA/\allowbreak Examples/\allowbreak TwoSeriesArrangementSym.lean}.
\begin{leancode}
theorem jetPoincareCoreT0_of_symArrangement {n : ℕ} (hn : 1 ≤ n)
    (hA : SymArrangementInclusion n) : PoincareZero.JetPoincareCoreT0 n :=
  fun _ _ _ hN hk hE hE' => jetPoincare_of_symArrangement hn hA hN hk hE hE'
\end{leancode}

\begin{dictionary}
Here \lean{hA} is the arrangement hypothesis (not graded unitarity; the letter is reused). The window of
$\mathrm{SAI}_n$ is exactly the image of the core window $2\le N$, $n+2\le k$, $2N+k\le E$, $2N+2k<E+n+2$ under $d=E-k$.
\end{dictionary}

\begin{proof}[Proof on paper]
Suppose $F$ of degree $k$, weight $E$ has $\partial_MF\in J_n$ for $M\ge N$, and suppose its image in $\ltwentynineRn$ is
not in the box. Then some functional $\psi$ on $\C[T]_{k,E}$ kills $J_n$ and the box but not $F$. Its dual polynomial
$f$ (symmetric, homogeneous of degree $d=E-k$) vanishes on the $n$-fold diagonals (because $\psi$ kills the jets of
$T^n$; \lean{dot_pol_jet_n}) and lies in $L_N$ (because $\psi$ kills the box). By $\mathrm{SAI}_n$, $f=\sum_lx_l^NY_l$
with $Y_l$ vanishing on the diagonals avoiding $l$. Pairing $F$ against $x_l^NY_l$ is pairing $\partial$-derivatives of
$F$ in the high variables against $Y_l$, and $Y_l$'s vanishing makes this a pairing of an element of $J_n$ with a
functional that kills $J_n$ (\lean{dot_pol_Wn}); so $\psi(F)=0$, a contradiction.
\end{proof}

\begin{physics}
The Poincar\'e lemma on the arc space and the exclusion statistics of the functional realisation are dual descriptions of
the same fact.
\end{physics}

\begin{status}
Proved, \texttt{\#print axioms}-clean (re-checked), conditional on $\mathrm{SAI}_n$ (supplied below).
\end{status}

\subsection{The Gr\"obner basis of the \texorpdfstring{$n$}{n}-equal ideal}

\begin{definition}[Block generators; the generation input]\label{l29:def:geninclusion}
For $k$ variables, a \emph{block labelling} is a map $c:\{0,\dots,k-1\}\to\{0,\dots,n-2\}$ (there are $n-1$ blocks). For
a strictly increasing list of points $p_0<\cdots<p_r$ inside one block and a power $m\ge N$, the \emph{block generator}
is
\begin{equation*}
g=\prod_{\substack{a<b,\ c(a)=c(b)\\ \{a,b\}\not\subseteq\{p_0,\dots,p_r\}}}(x_b-x_a)\;\cdot\;
\det\bigl(x_{p_i}^{0},x_{p_i}^{1},\dots,x_{p_i}^{r-1},x_{p_i}^{m}\bigr)_{i=0}^{r}.
\end{equation*}
The \emph{generation input} $\mathrm{GI}_n$: every $f\in L_N$ vanishing on all $n$-element diagonals lies in the ideal
generated by the block generators. The \emph{vanishing ideal} is the ideal of polynomials vanishing on all $n$-fold
diagonals.
\end{definition}

\noindent Files \leanpath{VOA/\allowbreak Examples/\allowbreak TwoSeriesJetPoincareVdm.lean},
\leanpath{VOA/\allowbreak Examples/\allowbreak TwoSeriesArrangementGen.lean} and
\leanpath{VOA/\allowbreak Examples/\allowbreak TwoSeriesGenIncl.lean}.
\begin{leancode}
def pprod (P : Finset (Fin k × Fin k)) : XP k := ∏ q ∈ P, (X q.2 - X q.1)

def dmat (pt : Fin (r + 1) → Fin k) (m : ℕ) : Matrix (Fin (r + 1)) (Fin (r + 1)) (XP k) :=
  Matrix.of fun p q => X (pt p) ^ (if q = Fin.last r then m else (q : ℕ))

def dD (pt : Fin (r + 1) → Fin k) (m : ℕ) : XP k := (dmat pt m).det

def blockPairs {M : ℕ} (c : Fin k → Fin M) (pt : Fin (r + 1) → Fin k) : Finset (Fin k × Fin k) :=
  univ.filter fun q => q.1 < q.2 ∧ c q.1 = c q.2 ∧ ¬ (q.1 ∈ univ.image pt ∧ q.2 ∈ univ.image pt)

def blockProd {M : ℕ} (c : Fin k → Fin M) (pt : Fin (r + 1) → Fin k) : XP k :=
  pprod (blockPairs c pt)

def genSet (n k N : ℕ) : Set (XP k) :=
  {g | ∃ (c : Fin k → Fin (n - 1)) (r : ℕ) (pt : Fin (r + 1) → Fin k) (m : ℕ),
    StrictMono pt ∧ (∀ p, c (pt p) = c (pt 0)) ∧ N ≤ m ∧ g = blockProd c pt * dD pt m}

def GenInclusion (n : ℕ) : Prop :=
  ∀ (k N : ℕ) (f : XP k), f ∈ Lid k N → (∀ S : Finset (Fin k), #S = n → VanishOn S f) →
    f ∈ Ideal.span (genSet n k N)
\end{leancode}

\begin{dictionary}
\lean{Matrix.of} builds the $(r+1)\times(r+1)$ matrix whose row $p$ is $(x_{p_p}^0,\dots,x_{p_p}^{r-1},x_{p_p}^m)$
(\lean{Fin.last r} is the last column); \lean{.det} its determinant, an alternant. \lean{StrictMono pt} says the points
are strictly increasing; \lean{c (pt p) = c (pt 0)} that they lie in one block. \lean{univ.image pt} is the set of
points. Note that $\mathrm{GI}_n$ does not ask $f$ to be symmetric: it is a statement about the full ideal.
\lean{vanIdeal n k} (in \leanpath{TwoSeriesGenIncl.lean}) is the vanishing ideal, and
\lean{span_genSet_le_vanIdeal} shows every block generator vanishes on the $n$-fold diagonals (by pigeonhole, $n$
coinciding points among $n-1$ blocks put two of them in one block, killing a Vandermonde factor or the alternant).
\end{dictionary}

\begin{physics}
These are, up to normalisation, the generators of the ideal of the $n$-equal arrangement found by Li and Li; the
alternant factor implements the cutoff $x^N$.
\end{physics}

\begin{status}
Definitions.
\end{status}

\begin{example}[$n=2$]
At $n=2$ there is one block, the labelling is constant, and a block generator is the Vandermonde product over all pairs
not both among the points, times an alternant in the points. With $r=k-1$ (all points) it is the single alternant
$\det(x_i^0,\dots,x_i^{k-2},x_i^m)$, which is divisible by the Vandermonde $\Delta=\prod_{a<b}(x_b-x_a)$; the
case $n=2$ of the inclusion is ``$f\in L_N$ vanishing on all $2$-fold diagonals (hence, being symmetric, divisible by
$\Delta^2$) lies in $\sum_lx_l^NQ_l$'', proved in \leanpath{TwoSeriesJetPoincareVdm.lean} by Vandermonde division
(\lean{arrangementInclusion_two}).
\end{example}

\begin{theorem}[The generation input for every $n\ge2$; the arrangement inclusions for every $n\ge1$]\label{l29:thm:geninclusion}
For every $n\ge2$, $\mathrm{GI}_n$ holds; hence $\mathrm{AI}_n$ holds. For every $n\ge1$, $\mathrm{SAI}_n$ holds.
\end{theorem}

\noindent Files \leanpath{VOA/\allowbreak Examples/\allowbreak TwoSeriesGenIncl.lean} (and
\leanpath{VOA/\allowbreak Examples/\allowbreak TwoSeriesArrangementGen.lean} for the passage $\mathrm{GI}\Rightarrow\mathrm{AI}$).
\begin{leancode}
theorem genInclusion {n : ℕ} (hn : 2 ≤ n) : GenInclusion n := by

theorem arrangementInclusion {n : ℕ} (hn : 2 ≤ n) : ArrangementInclusion n :=
  arrangementInclusion_of_genInclusion hn (genInclusion hn)

theorem symArrangementInclusion {n : ℕ} (hn : 1 ≤ n) :
    ArrangementSym.SymArrangementInclusion n := by
\end{leancode}

\begin{dictionary}
Full names \ltwentynineln{VOA.TwoSeries.GenIncl.genInclusion}, \lean{.arrangementInclusion}, \lean{.symArrangementInclusion}.
\ltwentynineln{arrangementInclusion_of_genInclusion} uses \lean{span_genSet_le_Wn}: every block generator is in $W_{N,n}$
(expand the alternant along its last column: each term is $x_{p}^m$ times a cofactor that, together with the block
Vandermonde, vanishes on the diagonals avoiding $p$). At $n=1$, $\mathrm{SAI}_1$ is immediate: vanishing on all
$1$-element ``diagonals'' (the substitution $x_i\mapsto x_i$) means $f=0$ (\lean{symArrangementInclusion_one}).
\textbf{What is not formalized.} The Lean statements are ideal \emph{membership}: every $f\in L_N\cap(\text{vanishing ideal})$ lies in
the span of the block generators. The stronger statement that the block generators form a lexicographic Gr\"obner basis of
$L_N\cap(\text{vanishing ideal})$ (their leading monomials generate its initial ideal) can be read off the paper proof below, since the
division remainder is $0$ for every such $f$; it is a consequence on paper and is not a Lean declaration. Note also that the ideal is
$L_N\cap(\text{vanishing ideal})$, not the vanishing ideal of the $n$-equal arrangement itself.
\end{dictionary}

\begin{proof}[Proof on paper]
Divide $f$ by the nonzero block generators in the lexicographic monomial order (\lean{MonomialOrder.div_set}): $f=\sum
a_bg_b+\rho$ with no monomial of the remainder $\rho$ divisible by a leading monomial of a generator. The remainder
still lies in $L_N$ and in the vanishing ideal. Suppose $\rho\ne0$ with leading exponent $e$. Since $L_N$ is a monomial
ideal, some $e_i\ge N$. Since $\rho$ vanishes at every point with an $n$-fold repeated value (\lean{eval_eq_zero_of_Q}),
a combinatorial ``point game'' on $e$ (\lean{not_win_of_vanish}, \lean{exists_blocks_of_not_win}) produces a block
labelling $c$ with $\#\{j>i':c_j=c_{i'}\}\le e_{i'}$ for every $i'$. From $c$, the suffix of the block of $i$ starting at
$i$, and the power $\max(N,r)$, one builds a nonzero block generator whose leading exponent is $\le e$ componentwise
(\lean{exists_small_gen}); this contradicts the remainder property. So $\rho=0$.
\end{proof}

\begin{physics}
The functional realisation of the $(A_1,A_{2n-2})$ arc space has a clean, explicit basis of exclusion-rule
constraints; nothing beyond the obvious Vandermonde-type relations is needed.
\end{physics}

\begin{status}
Proved, \textbf{unconditional}. Axioms standard (re-checked: all three).
\end{status}

%======================================================================
\section{The headline theorem}\label{l29:sec:headline}
%======================================================================

\begin{theorem}[Graded unitarity transfers along the coset, at every boundary level]\label{l29:thm:final}
For every $m\in\N$: if $L_{-2+2/(2m+3)}(\ltwentyninesl)$, with its canonical form, compact conjugation and corrected
current-and-stress flag, is graded unitary, then the minimal model $M(2,2m+5)=L(c_{2,2m+5},0)$, with its canonical form,
partition-basis real structure and length flag, is graded unitary.
\end{theorem}

\noindent File \leanpath{VOA/\allowbreak Examples/\allowbreak TwoSeriesGenIncl.lean}.
\begin{leancode}
theorem twoSeries_gradedUnitary_of_affine_final (m : ℕ)
    (hA : GradedUnitary (affineForm (twoSeriesLevels (m + 1)))
      (affineRealStructure (twoSeriesLevels (m + 1))) (affineFlag (twoSeriesLevels (m + 1)))) :
    VirGradedUnitary (twoSeriesLevels (m + 1)).c (targetRealStructure (twoSeriesLevels (m + 1))) :=
  twoSeries_gradedUnitary_of_affine_symArrangement m hA
    (GenIncl.symArrangementInclusion (by omega))
\end{leancode}

\begin{dictionary}
Full name \ltwentynineln{VOA.TwoSeries.JetPoincareTwo.twoSeries_gradedUnitary_of_affine_final}, declared in
\leanpath{TwoSeriesGenIncl.lean} (the file reopens the namespace \lean{VOA.TwoSeries.JetPoincareTwo} at its end).
\begin{itemize}
\item \textbf{Index.} $n=m+1\ge1$: $(\lean{twoSeriesLevels (m + 1)}).kA=k_{m+1}=-2+2/(2m+3)$ and
  $(\lean{twoSeriesLevels (m + 1)}).c=c_{2,2m+5}$ (Definition~\ref{l29:def:twolevels}). At $m=0$ the objects are
  definitionally those of the Lee--Yang chapter.
\item \textbf{The one hypothesis} \lean{hA}: graded unitarity (Chapter~\ref{ch:L19}, Definition~\ref{l19:def:gu}) of the
  actual simple VOA \lean{AVOA} with the data of Definition~\ref{l29:def:endpoint}. It is open for every $m$.
\item \textbf{The conclusion}: \lean{VirGradedUnitary c R} is \lean{GradedUnitary (Bsimp c) R (virGradedFlag c)}
  (Definition~\ref{l19:def:virgu}) for the actual simple Virasoro VOA, with $R$ the target real structure
  (Definition~\ref{l29:def:target}).
\item \textbf{The proof term} composes: Theorem~\ref{l29:thm:geninclusion} $\to$ Theorem~\ref{l29:thm:poincarecore}
  $\to$ Theorem~\ref{l29:thm:slicegu} (slice map, Arakawa property, Virasoro symbol) $\to$
  Theorem~\ref{l29:thm:strictpoincare} $\to$ strict model maps and model socle data $\to$
  Theorem~\ref{l29:thm:soclespectrum} (socle words, Theorem~\ref{l29:thm:swc}) $\to$
  Theorem~\ref{l29:thm:gordonfinal} (characters, Section~\ref{l29:sec:char}) $\to$ Theorem~\ref{l29:thm:invariants}
  (Koszul) $\to$ Theorem~\ref{l29:thm:hodge} (Hodge package) $\to$ Theorem~\ref{l29:thm:transfer}.
\end{itemize}
\textbf{Could it be vacuous?} Only if \lean{hA} were unsatisfiable, i.e.\ if $L_{k_n}$ provably failed graded
unitarity; then the theorem would say nothing. That is not known either way. The predicate itself is not degenerate:
its bilinear-positivity version, which \emph{was} unsatisfiable, was repaired (Chapter~\ref{ch:L19},
Theorem~\ref{l19:thm:bilpos}), and its weight-zero slice is checked to have the right sign
(Theorem~\ref{l19:thm:vacuum}). Numerical evidence (not a proof; computed outside the library, not re-run here):
$L_{-4/3}$, $L_{-8/5}$ and $L_{-12/7}$ pass the corrected-flag graded-unitarity test through weight $7$.
\textbf{Could the conclusion be weaker than intended?} No: the conclusion is about the simple quotient, with the pinned
length flag, at every weight.
\end{dictionary}

\begin{proof}[Proof on paper]
The chain \eqref{l29:eq:chain} of Remark~\ref{l29:rem:arc}, with every input other than $\ltwentynineGU(A)$ discharged:
the Hodge package and the Koszul embedding from $\ltwentynineGU(A)$ (Section~\ref{l29:sec:chain}); the refined PBW
character of $M(2,2n+3)$ (Section~\ref{l29:sec:char}); the invariant count $\dim\ltwentynineinv_{r,N}\le\ltwentyninegc_n(N,r)$
from the socle data of $\ltwentyninegr_rA$, obtained by pulling back the socle data of the jet model along strict model
maps (Section~\ref{l29:sec:model}), with detection from the slice map (Section~\ref{l29:sec:ds}) and the model socle data
from the jet-ideal Poincar\'e lemma (Section~\ref{l29:sec:arr}).
\end{proof}

\begin{physics}
If the chiral algebra of the Argyres--Douglas theory $(A_1,D_{2m+3})$ satisfies the positivity constraints that
four-dimensional unitarity imposes, so does the chiral algebra of $(A_1,A_{2m+2})$. Equivalently, a violation of graded
unitarity in any $M(2,2n+3)$ would rule out graded unitarity of $L_{k_n}(\ltwentyninesl)$. Since Kaufmann's conjecture
(graded unitarity of the boundary minimal models) is open, the theorem reduces the whole $(2,2n+3)$ family to the affine
side.
\end{physics}

\begin{status}
Proved implication; \texttt{\#print axioms} $=$ \lean{[propext, Classical.choice, Quot.sound]} (re-checked for this edition
against the existing build products; the aggregate \texttt{lake build} was not re-run, Section~\ref{l29:sec:verification});
\textbf{conditional on \lean{hA} only}. Not unconditional: \lean{hA} is open for every $m$.
\end{status}

\begin{corollary}[The first two members, and the arrangement form]\label{l29:cor:m27}
$\ltwentynineGU(L_{-8/5})\Rightarrow\ltwentynineGU(M(2,7))$ and $\ltwentynineGU(L_{-4/3})\Rightarrow\ltwentynineGU(M(2,5))$,
each with no other hypothesis; and for every $m$, $\ltwentynineGU(L_{k_{m+1}})\wedge\mathrm{AI}_{m+1}\Rightarrow
\ltwentynineGU(M(2,2m+5))$.
\end{corollary}

\noindent File \leanpath{VOA/\allowbreak Examples/\allowbreak TwoSeriesJetPoincareTwoEndpoint.lean}.
\begin{leancode}
theorem twoSeries_gradedUnitary_M27_of_affine
    (hA : GradedUnitary (affineForm (twoSeriesLevels 2))
      (affineRealStructure (twoSeriesLevels 2)) (affineFlag (twoSeriesLevels 2))) :
    VirGradedUnitary (twoSeriesLevels 2).c (targetRealStructure (twoSeriesLevels 2)) :=
  SliceMap.twoSeries_gradedUnitary_of_affine 1 hA jetPoincareCoreT0_two

theorem twoSeries_gradedUnitary_M25_of_affine
    (hA : GradedUnitary (affineForm (twoSeriesLevels 1))
      (affineRealStructure (twoSeriesLevels 1)) (affineFlag (twoSeriesLevels 1))) :
    VirGradedUnitary (twoSeriesLevels 1).c (targetRealStructure (twoSeriesLevels 1)) :=
  SliceMap.twoSeries_gradedUnitary_of_affine 0 hA jetPoincareCoreT0_one
\end{leancode}

\begin{dictionary}
These predate the general theorem: they use the direct proofs of $\mathrm{JPCore}_2$ (\lean{jetPoincareCoreT0_two},
by the $n=2$ Vandermonde argument) and $\mathrm{JPCore}_1$. \lean{twoSeriesLevels 2}: $k=-8/5$, $c=-68/7$. The arrangement
form \ltwentynineln{twoSeries_gradedUnitary_of_affine_arrangement} (same file) takes $\mathrm{AI}_{m+1}$ as a hypothesis, now
discharged by Theorem~\ref{l29:thm:geninclusion}.
\end{dictionary}

\begin{physics}
$(A_1,D_5)\to(A_1,A_4)$ and $(A_1,D_3)\to(A_1,A_2)$.
\end{physics}

\begin{status}
Proved, \texttt{\#print axioms}-clean (re-checked), conditional on \lean{hA} only (and \lean{hArr} for the arrangement
form).
\end{status}

%======================================================================
\section{Consistency at \texorpdfstring{$n=1$}{n=1}}\label{l29:sec:leeyang}
%======================================================================

Five files check that the general machinery reproduces the Lee--Yang results of Chapter~\ref{ch:L28}:
\leanpath{TwoSeriesLeeYang} (the general Theorem~\ref{l29:thm:transfer} at $n=1$, with the Hodge package and the squeeze
input taken from the concrete Lee--Yang proofs), \leanpath{TwoSeriesKoszulLeeYang} (the invariant count at $n=1$ from one
cyclic vector), \leanpath{TwoSeriesSpectrumLeeYang} (the socle data of $\ltwentyninegr_rL_{-4/3}$, proved directly),
\leanpath{TwoSeriesModelSocleLevelOne} and \leanpath{TwoSeriesStrictLevelOne} (the model at $n=1$, where $R_1=\C$, and its
character bound). Each ends in a Lee--Yang theorem conditional only on graded unitarity of $L_{-4/3}$. We quote the first.

\begin{theorem}[Lee--Yang from the general theorem]\label{l29:thm:leeyang}
$\ltwentynineGU(L_{-4/3}(\ltwentyninesl))\Rightarrow\ltwentynineGU(M(2,5))$, via Theorem~\ref{l29:thm:transfer} at $n=1$.
\end{theorem}

\noindent File \leanpath{VOA/\allowbreak Examples/\allowbreak TwoSeriesLeeYang.lean}.
\begin{leancode}
theorem leeYang_gradedUnitary_of_affine_viaTwoSeries
    (hA : GradedUnitary LeeYangTransferAssembly.affineForm
      LeeYangTransferAssembly.affineRealStructure LeeYangTransferAssembly.affineFlag) :
    VirGradedUnitarity.VirGradedUnitary LeeYang.leeYangC leeYangRealStructure :=
  twoSeries_gradedUnitary_of_affine 1 hA (leeYangHodgeInputs hA) (twoSeriesSqueezeInput_one hA)
\end{leancode}

\begin{dictionary}
The hypothesis and conclusion are stated with the \emph{concrete} Lee--Yang objects of Chapter~\ref{ch:L28}; the proof
applies the general \ltwentynineln{VOA.TwoSeries.twoSeries_gradedUnitary_of_affine} at $n=1$ (the short name resolves to the
\leanpath{TwoSeriesAssembly} declaration inside \lean{VOA.TwoSeries}). That this typechecks is the content of the
\lean{rfl} lemmas: \lean{twoSeriesLevels 1} \emph{is} \lean{leeYangLevels}, and every generic object at it is the concrete
one. The squeeze coefficient at $n=1$ is \lean{gapTwoCoefficient} (Rogers--Ramanujan).
\end{dictionary}

\begin{physics}
The general machine, run at $(A_1,D_3)\to(A_1,A_2)$, reproduces the Lee--Yang result.
\end{physics}

\begin{status}
Proved, \texttt{\#print axioms}-clean (re-checked), conditional on \lean{hA} only.
\end{status}

%======================================================================
\section{The other endpoints}\label{l29:sec:endpointstable}
%======================================================================

The library keeps every intermediate endpoint of the reduction, each with its own extra hypotheses. They are all
\texttt{\#print axioms}-clean and all \textbf{conditional} on \lean{hA} plus the listed inputs; none is needed once the
headline is available, but each records an honest reduction and some are reused by others. Namespaces abbreviate
\lean{VOA.TwoSeries} as \lean{TS}; $n\ge1$ is assumed where marked.
\begin{center}\small
\begin{longtable}{p{0.40\textwidth}p{0.20\textwidth}p{0.30\textwidth}}
\toprule
declaration (file) & index & extra hypotheses\\
\midrule\endhead
\ltwentynineln{TS.cosetLevels_gradedUnitary_of_affine} (Assembly) & any $L$ & Hodge, squeeze\\
\ltwentynineln{TS.twoSeries_gradedUnitary_of_affine} (Assembly) & $n$ & Hodge, squeeze\\
\ltwentynineln{TS.twoSeries_gradedUnitary_of_affine_hodgeFree} (HodgeInputs) & $n$ & squeeze\\
\ltwentynineln{TS.twoSeries_gradedUnitary_of_affine_invariants} (KoszulInvariants) & $n$ & coeff, \lean{hinv}, \lean{hpbw}\\
\ltwentynineln{TS.twoSeries_gradedUnitary_of_affine_invariants_gordon} (CharacterEndpoint) & $n$ & \lean{hinv}, \lean{GordonLowerBound}\\
\ltwentynineln{TS.twoSeries_gradedUnitary_of_affine_invariants_gordon_final} (CharacterAGIdentity) & $n$ & \lean{hinv}\\
\ltwentynineln{TS.twoSeries_gradedUnitary_of_socle} (InvariantsSocle) & $n$ & socle input, \lean{SocleWordGordon}\\
\ltwentynineln{TS.twoSeries_gradedUnitary_of_socleInput} (SocleCountGordon) & $n$ & socle input\\
\ltwentynineln{TS.twoSeries_gradedUnitary_of_socleSpectrum} (SocleDual) & $n$ & socle data of $\ltwentyninegr A$\\
\ltwentynineln{TS.twoSeries_gradedUnitary_of_socleModels} (SpectrumModel) & $n$ & socle models\\
\ltwentynineln{TS.StrictModel.twoSeries_gradedUnitary_of_strictModel} (StrictEndpoint) & $n\ge1$ & det, chi, spec\\
\ltwentynineln{TS.ModelChar.twoSeries_gradedUnitary_of_strictModel_det}, \lean{_of_fermionicLower}, \lean{_two_of_strictModel} (ModelCharEndpoint) & $n\ge1$; $n=2$ & det, spec (and ferm)\\
\ltwentynineln{TS.StrictModel.twoSeries_gradedUnitary_of_evalFunctional}, \lean{_of_coinv} (DetectFunctional) & $n\ge1$ & eval.\ functionals, spec\\
\ltwentynineln{..._of_strictModel_membership}, \lean{_of_coinv_membership} (ModelSocleEndpoint) & $n\ge1$ & det (or functionals), membership\\
\lean{..._of_strictModel_testBound}, \lean{_box} (MembershipEndpoint) & $n\ge1$ & det, test bound / box inclusion and count\\
\lean{..._of_strictModel_incl} (BoxFusion), \lean{_symBox} (BoxSym), \lean{_box_two} (BoxTwo) & $n\ge1$; $n=2$ & det, box inclusion (and \lean{SymBoxBound})\\
\lean{..._of_strictModel_inclHigh}, \lean{_level_one} (InclDegOne); \lean{_inclTail} (InclLow) & $n\ge1$; $n=1$ & det, inclusion variants\\
\lean{..._of_strictModel_singular} (TailWitt); \lean{_singularLift} (SingTailLift); \lean{_singularLift_gt} (SingTailDiag); \lean{_singularLiftHigh}, \ltwentynineln{twoSeries_gradedUnitary_one_of_strictModel} (LiftDegN) & $n\ge1$; $n=1$ & det, singular-tail / lift inputs\\
\ltwentynineln{..._of_strictModel_jetPoincare} (HighEndpoint); \lean{_jetPoincare2}, \lean{_jetPoincareCore} (HighDegSucc); \lean{_jetPoincareCoreT0} (PoincareZero) & $n\ge1$ & det, a Poincar\'e input\\
\ltwentynineln{TS.DSSlice.twoSeries_gradedUnitary_of_dsSlice} (DSSlice) & $n=m+1$ & slice map, Arakawa, \lean{JetPoincareCoreT0}\\
\ltwentynineln{TS.SliceMap.twoSeries_gradedUnitary_of_sliceArakawa} (SliceMap) & $n=m+1$ & Arakawa, \lean{JetPoincareCoreT0}\\
\ltwentynineln{TS.SliceMap.twoSeries_gradedUnitary_of_maximalGeneration}, \lean{_of_adamovicMilas} (ArakawaBottom); \lean{_of_singularBottom}, \lean{_of_singularBottomInput} (ArakawaEndpoint) & $n=m+1$ & generation of the maximal ideal / Adamovi\'c--Milas data / bottom input, \lean{JetPoincareCoreT0}\\
\ltwentynineln{TS.SliceMap.twoSeries_gradedUnitary_of_affine} (ArakawaUnconditional) & $n=m+1$ & \lean{JetPoincareCoreT0}\\
\ltwentynineln{TS.JetPoincareTwo.twoSeries_gradedUnitary_of_affine_arrangement} (JetPoincareTwoEndpoint) & $n=m+1$ & \lean{ArrangementInclusion}\\
\ltwentynineln{TS.JetPoincareTwo.twoSeries_gradedUnitary_of_affine_symArrangement} (ArrangementSym) & $n=m+1$ & \lean{SymArrangementInclusion}\\
\ltwentynineln{TS.JetPoincareTwo.twoSeries_gradedUnitary_M27_of_affine}, \lean{_M25_of_affine} & $n=2$, $1$ & none\\
\ltwentynineln{TS.JetPoincareTwo.twoSeries_gradedUnitary_of_affine_final} (GenIncl) & $n=m+1$ & \textbf{none}\\
\ltwentynineln{TS.leeYang_gradedUnitary_of_affine_viaTwoSeries}, \lean{_viaInvariants}, \ltwentynineln{leeYang_gradedUnitary_viaSocleSpectrum} & $n=1$ & none\\
\bottomrule
\end{longtable}
\end{center}
The \leanpath{TwoSeriesArakawaBottom}/\leanpath{TwoSeriesArakawaEndpoint} endpoints are an earlier route to the slice
Arakawa property through a description of the maximal ideal by a singular vector (Adamovi\'c--Milas); the unconditional
route of Theorem~\ref{l29:thm:slicearakawa} made it unnecessary. The remaining files not mentioned so far
(\leanpath{TwoSeriesSliceMapBasis}, \leanpath{-Whittaker}, \leanpath{-Sugawara}, \leanpath{-Vir};
\leanpath{TwoSeriesArakawaClosure}, \leanpath{-Intertwine}, \leanpath{-Positive}, \leanpath{-Sugawara},
\leanpath{-Transport}, \leanpath{-Triplet}; \leanpath{TwoSeriesModelChar*}; \leanpath{TwoSeriesCharacter*};
\leanpath{TwoSeriesSocleCount*}, \leanpath{TwoSeriesSocleDualCore}, \leanpath{TwoSeriesSpectrumCore};
\leanpath{TwoSeriesBox*}, \leanpath{TwoSeriesIncl*}, \leanpath{TwoSeriesHigh*}, \leanpath{TwoSeriesStrict*},
\leanpath{TwoSeriesMembership*}, \leanpath{TwoSeriesJetPoincare*}, \leanpath{TwoSeriesArrangementGen},
\leanpath{TwoSeriesGenInclLex}) are the internal lemmas of the proofs sketched above: PBW bases and Whittaker
coinvariants for the slice map; the lifts $Q_0,Q_1,Q_2$ and the vacuum-row argument; the Bailey chain, the
Feigin--Stoyanovsky functionals and the profile counting of the model character; Andrews' recursion and the
Kac-determinant rank bound; MacMahon's theorem and the fusion recursion; the Vandermonde and alternant calculus and the
point game.

%======================================================================
\section{Warnings: stale documentation and dead ends}\label{l29:sec:deadends}
%======================================================================

\begin{warning}[title={Stale docstrings that call proved inputs ``open''}]\label{l29:warn:stale}
The Lean statements are correct; the following documentation comments lag behind later files. A reader who trusts the
docstring over the declaration will under-estimate what is proved.
\begin{enumerate}
\item \lean{BRSTHodgeInputs} (\leanpath{TwoSeriesAssembly.lean}): ``in Lean proved so far only at Lee--Yang''. It is
  provided at every level datum by \lean{brstHodgeInputs} (Theorem~\ref{l29:thm:hodge}).
\item \lean{GordonLowerBound} and the header of \leanpath{TwoSeriesCharacterReduction.lean}: ``OPEN for general $n$ \dots\
  not formalised here''. Proved for the series (\lean{gordonLowerBound_twoSeries}, Theorem~\ref{l29:thm:character}).
\item \leanpath{TwoSeriesInvariantsSocle.lean}: \lean{SocleWordGordon} ``OPEN as a Lean statement''; proved
  (\lean{socleWordGordon}, Theorem~\ref{l29:thm:swc}). \lean{TwoSeriesSocleInput} and \lean{TwoSeriesSocleSpectrum}
  (\leanpath{TwoSeriesSocleDual.lean}) are labelled ``OPEN''; for every $n\ge1$ they now follow from the chain of
  Sections~\ref{l29:sec:model}--\ref{l29:sec:arr}, although no declaration states them as standalone theorems.
\item \leanpath{TwoSeriesJetPoincareArr.lean} and \leanpath{TwoSeriesArrangementSym.lean}: ``\lean{ArrangementInclusion n}
  (resp.\ \lean{SymArrangementInclusion n}) for $n\ge3$ is the named open input''. Proved for every $n\ge2$ (resp.\ $\ge1$)
  in \leanpath{TwoSeriesGenIncl.lean} (Theorem~\ref{l29:thm:geninclusion}).
\item \leanpath{TwoSeriesDSSlice.lean} header: calls the Virasoro symbol input ``classical freeness of $M(2,2n+1)$''.
  \lean{virSymbol_twoSeries} is only a PBW-length leading-symbol theorem into the presented ring; intrinsic classical
  freeness is a separate theorem (Theorem~\ref{l29:thm:virsymbol}, dictionary). The docstring of
  \lean{virSymbol_twoSeries} itself was corrected in a later revision of that file.
\item \leanpath{TwoSeriesSliceMap.lean} header: ``the only universal-object input used is the sanctioned library PBW
  machinery''. Since \lean{VOA.pbw} is a theorem there is no sanctioned axiom; the phrase is historical, and
  \texttt{\#print axioms} of \lean{dsSliceMap} is standard.
\end{enumerate}
These are documentation defects (severity S4: a docstring or note is wrong or stale; no statement is affected).
\end{warning}

\begin{warning}[title={Dead ends}]\label{l29:warn:deadends}
The following routes were tried for the invariant count and the detection input and failed; they are recorded so that
they are not re-attempted. The evidence is the project's record (exact linear algebra modulo a large prime, computed outside the
library); it was not re-run for this chapter.
\begin{enumerate}
\item \emph{One cyclic dual vector} (warning on p.~\pageref{l29:warn:cyclic}): fails at $n=2$, $r=2$, $N=4$.
\item \emph{Test substitutions in the membership bound}: with the plain substitution $T\mapsto J$ (no correction term)
  the test space is strictly larger than the socle count at $T$-degree $\ge n$ ($n=2$); with $T\mapsto J+\tilde z$ it is
  also too large; with $\tilde x$ or $\tilde y$ it is sharp (\leanpath{TwoSeriesMembershipEndpoint.lean} docstring).
\item \emph{Detection with the wrong ring}: the model symbols satisfy every relation of $\ltwentyninegr_rL_{k_n}$ modulo the
  jets of $T^n$, but \emph{not} modulo the jets of $T^{n+1}$ ($n=2$, weight $\le8$; $n=3$, weight $\le7$). So $\ltwentynineRn$,
  the arc algebra of $M(2,2n+1)$, is forced.
\item \emph{The slice map at the wrong central charge}: the inclusion of the image of the affine maximal ideal into the
  Virasoro maximal ideal holds at $c_{2,2n+1}$ and fails at $c_{2,2n+3}$ ($n=2,3$, weights $\le10$).
\item \emph{Other detection attempts}: a weight-zero Witt operator, single-$\varepsilon$
  specialisations at $k\ge n$, a test using only $J$, and Riccati-type maps. None gave detection; the slice map did.
\end{enumerate}
\end{warning}

\begin{warning}[title={The converse is not proved}]\label{l29:warn:converse}
Nothing in this chapter goes from $M(2,2n+3)$ back to $L_{k_n}$. At $n=1$ there is no coset route back: the coset partner
$L_{-1/3}(\ltwentyninesl)$ fails graded unitarity at weight $3$ (a norm $-7/15$; the project's numerical record, not
machine-checked), and BRST positivity only flows forward.
No converse implication is stated or claimed in this edition.
\end{warning}

%======================================================================
\section{Verification record}\label{l29:sec:verification}
%======================================================================

\paragraph{Files.} The chapter was written against an earlier state of the library; the files it documents are
\leanpath{VOA/Examples/TwoSeries*.lean} (126 files at that state, 128 in this release). All are in the import closure of the umbrella \leanpath{VOA/Examples.lean}.

\paragraph{Quotations.} Every \texttt{leancode} block of this chapter was compared, piece by piece with comments and whitespace
normalized, with the current Lean sources for this edition: all occur verbatim. Every Lean identifier and Lean file
path cited in the chapter resolves to a declaration or file of the current library.

\paragraph{What was re-run for this edition.} \texttt{\#print axioms} on one scratch file importing
\lean{TwoSeriesGenIncl}, \lean{TwoSeriesArakawaUnconditional}, \lean{TwoSeriesCharacterAGIdentity}, \lean{TwoSeriesModelCharEndpoint},
\lean{TwoSeriesJetPoincareTwoEndpoint}, \lean{TwoSeriesLeeYang}, \lean{TwoSeriesSocleCountGordon}, \lean{TwoSeriesBoxFusion},
\lean{TwoSeriesHighDegSucc}, \lean{TwoSeriesModelSocleCartan} and \lean{TwoSeriesKoszulLeeYang} (all under \lean{VOA.Examples}), run as
\texttt{lake env lean -DmaxSynthPendingDepth=3} from the repository root against the build products already in the checkout (no
\texttt{lake build} was run). Every declaration below, and the composite \lean{detect_check} of Theorem~\ref{l29:thm:detect}, reported
exactly \lean{[propext, Classical.choice, Quot.sound]}, with no error and no \lean{sorryAx}:
\begin{plaincode}
JetPoincareTwo.twoSeries_gradedUnitary_of_affine_final, GenIncl.genInclusion,
GenIncl.arrangementInclusion, GenIncl.symArrangementInclusion,
ArrangementSym.jetPoincareCoreT0_of_symArrangement,
JetPoincareTwo.twoSeries_gradedUnitary_of_affine_symArrangement,
SliceMap.twoSeries_sliceArakawa, SliceMap.sliceArakawa_unconditional, SliceMap.W1_maxIdeal,
SliceMap.twoSeries_gradedUnitary_of_affine, SliceMap.dsSliceMap,
DSSlice.twoSeriesModelDetects_of_dsSlice, DSSlice.virSymbol_twoSeries,
StrictModel.twoSeries_gradedUnitary_of_strictModel_jetPoincareCoreT0,
StrictModel.twoSeriesModelDetects_one, gordonLowerBound_twoSeries, twoSeries_character,
twoSeries_hpbw, twoSeries_gradedUnitary_of_affine_invariants_gordon_final,
twoSeries_hasJetRelations_all, twoSeries_gradedUnitary_of_affine_invariants,
rank_leadCohom_le_affineInvariants, twoSeries_gradedUnitary_of_affine_hodgeFree,
brstHodgeInputs, twoSeries_gradedUnitary_of_affine, cosetLevels_gradedUnitary_of_affine,
auxiliaryGradedUnitary, twoSeries_charge, twoSeriesLevels_c,
ModelChar.twoSeriesModelCharacter_all, ModelChar.fermSeries_eq_product,
ModelChar.twoSeries_gradedUnitary_of_strictModel_det,
StrictModel.twoSeries_gradedUnitary_of_strictModel,
JetPoincareTwo.twoSeries_gradedUnitary_M27_of_affine,
JetPoincareTwo.twoSeries_gradedUnitary_M25_of_affine,
JetPoincareTwo.twoSeries_gradedUnitary_of_affine_arrangement,
JetPoincareArr.arrangementInclusion_two, leeYang_gradedUnitary_of_affine_viaTwoSeries,
leeYang_gradedUnitary_of_affine_viaInvariants, twoSeries_gradedUnitary_of_socleInput,
StrictModel.boxCount_all, StrictModel.symBoxBound,
StrictModel.twoSeries_gradedUnitary_of_strictModel_jetPoincareCore,
StrictModel.model_transverse
  (the names above relative to VOA.TwoSeries)
TwoSeriesCharacter.AG.gordon_sum_eq_rcq, TwoSeriesCharacter.finrank_pieceIn_le_gordon,
TwoSeriesCharacter.hpbw_iff_gordonLowerBound, TwoSeriesCharacter.KacGram.rcq_le_finrank,
SocleCount.socleWordCount_eq, SocleCount.socleWordGordon
  (these relative to VOA)
\end{plaincode}
The same run printed the type of the headline as
\begin{plaincode}
VOA.TwoSeries.JetPoincareTwo.twoSeries_gradedUnitary_of_affine_final : ∀ (m : ℕ),
  VOA.GradedUnitary (VOA.TwoSeries.affineForm (VOA.TwoSeries.twoSeriesLevels (m + 1)))
      (VOA.TwoSeries.affineRealStructure (VOA.TwoSeries.twoSeriesLevels (m + 1)))
      (VOA.TwoSeries.affineFlag (VOA.TwoSeries.twoSeriesLevels (m + 1))) →
    VOA.VirGradedUnitarity.VirGradedUnitary (VOA.TwoSeries.twoSeriesLevels (m + 1)).c
      (VOA.TwoSeries.targetRealStructure (VOA.TwoSeries.twoSeriesLevels (m + 1)))
\end{plaincode}
so, besides the natural number $m$, its only argument is the hypothesis \lean{hA}. The build products used are those already in the
checkout; they were not rebuilt for this edition.

\paragraph{Records of earlier builds (not re-run).} The project's records give, for the state in which the headline was first
obtained, an aggregate \texttt{lake build} green at 13\,672 jobs and the maintained axiom audit passing, with the headline
reporting the three standard axioms; and a later aggregate \texttt{lake build VOA} green at 13\,868 jobs. The aggregate build was not
re-run for this edition. The numerical evidence quoted in the warnings on pp.~\pageref{l29:warn:cyclic} and \pageref{l29:warn:deadends}
and in the dictionary of Theorem~\ref{l29:thm:final} is the project's record (computations outside the library), not machine-checked here.

\paragraph{Forbidden constructs.} A search of the \leanpath{TwoSeries*} files for \lean{sorry}, \lean{admit}, \lean{native_decide},
\lean{implemented_by} (as words) and for lines beginning with \lean{axiom} finds no occurrence; this is consistent with the axiom
checks above.

%======================================================================
\section{Referee's checklist}\label{l29:sec:checklist}
%======================================================================

\begin{center}\small
\begin{longtable}{p{0.29\textwidth}p{0.26\textwidth}p{0.17\textwidth}p{0.16\textwidth}}
\toprule
declaration & intended meaning & hypotheses & where checked\\
\midrule\endhead
\lean{TS.JetPoincareTwo.}\allowbreak\ltwentynineln{twoSeries_gradedUnitary_of_affine_final} & $\ltwentynineGU(L_{k_{m+1}})\Rightarrow\ltwentynineGU(M(2,2m+5))$ & \lean{hA} only; conditional & Thm.~\ref{l29:thm:final}\\
\lean{TS.twoSeries_charge} & coset charge $=c_{2,2n+3}$ & none & Thm.~\ref{l29:thm:charge}\\
\lean{TS.twoSeriesLevels} & level data, $=$ Lee--Yang at $n=1$ & -- & Def.~\ref{l29:def:twolevels}\\
\lean{TS.auxiliaryGradedUnitary} & GU of $V^{k_U}$, $k_U<-2$ & none & Thm.~\ref{l29:thm:aux}\\
\ltwentynineln{TS.cosetLevels_gradedUnitary_of_affine} & coset transfer, packaged & \lean{hA}, Hodge, squeeze & Thm.~\ref{l29:thm:transfer}\\
\ltwentynineln{TS.twoSeries_gradedUnitary_of_affine} (Assembly) & same, at member $n$ & same & Thm.~\ref{l29:thm:transfer}, warning p.~\pageref{l29:warn:clash}\\
\lean{TS.brstHodgeInputs} & Hodge package at every level & \lean{hA} & Thm.~\ref{l29:thm:hodge}\\
\ltwentynineln{TS.rank_leadCohom_le_affineInvariants} & Koszul embedding & \lean{hA} & Thm.~\ref{l29:thm:koszul}\\
\ltwentynineln{TS.twoSeries_gradedUnitary_of_affine_invariants} & transfer from counts & \lean{hA}, \lean{hinv}, \lean{hpbw} & Thm.~\ref{l29:thm:invariants}\\
\ltwentynineln{TS.twoSeries_hasJetRelations_all} & null vector $T^{n+1}$ & none & Thm.~\ref{l29:thm:jetrel}\\
\lean{VOA.TwoSeriesCharacter.}\allowbreak\lean{finrank_pieceIn_le_gordon} & Gordon spanning & jet relations & Thm.~\ref{l29:thm:upper}\\
\lean{...AG.gordon_sum_eq_rcq} & Andrews--Gordon identity & none & Thm.~\ref{l29:thm:ag}\\
\lean{TS.twoSeries_character} & vacuum character of $M(2,2n+3)$ & none; \textbf{unconditional} & Thm.~\ref{l29:thm:character}\\
\lean{TS.twoSeries_hpbw} & refined PBW character & none & Thm.~\ref{l29:thm:hpbw}\\
\ltwentynineln{TS.twoSeries_gradedUnitary_of_affine_invariants_gordon_final} & transfer from \lean{hinv} & \lean{hA}, \lean{hinv} & Thm.~\ref{l29:thm:gordonfinal}\\
\ltwentynineln{VOA.SocleCount.socleWordCount_eq} & socle count $=$ Gordon count & none & Thm.~\ref{l29:thm:swc}\\
\ltwentynineln{TS.twoSeries_gradedUnitary_of_socleSpectrum} & transfer from socle data & \lean{hA}, socle data & Thm.~\ref{l29:thm:soclespectrum}\\
\lean{TS.StrictModel.}\allowbreak\ltwentynineln{twoSeries_gradedUnitary_of_strictModel} & transfer from model & \lean{hA}, det, chi, spec & Thm.~\ref{l29:thm:strict}\\
\ltwentynineln{TS.ModelChar.fermSeries_eq_product}, \lean{twoSeriesModelCharacter_all} & fermionic $=$ product; model big enough & none ($n\ge1$) & Thm.~\ref{l29:thm:modelchar}\\
\lean{TS.StrictModel.}\allowbreak\lean{..._jetPoincareCoreT0} & transfer from det $+$ Poincar\'e & \lean{hA}, det, \lean{h} & Thm.~\ref{l29:thm:strictpoincare}\\
\ltwentynineln{TS.DSSlice.virSymbol_twoSeries} & Virasoro leading symbols into $R_{m+1}$ & none & Thm.~\ref{l29:thm:virsymbol}\\
\lean{TS.DSSlice.}\allowbreak\ltwentynineln{twoSeriesModelDetects_of_dsSlice} & detection from slice map & slice map, Arakawa & Thm.~\ref{l29:thm:detectreduce}\\
\lean{TS.SliceMap.dsSliceMap} & slice map constructed & -- & Def.~\ref{l29:def:dsslicemap}\\
\ltwentynineln{TS.SliceMap.twoSeries_sliceArakawa} & nulls go to nulls under DS & none & Thm.~\ref{l29:thm:slicearakawa}\\
\ltwentynineln{TS.SliceMap.twoSeries_gradedUnitary_of_affine} & transfer from Poincar\'e core & \lean{hA}, \lean{h} & Thm.~\ref{l29:thm:slicegu}, warning p.~\pageref{l29:warn:clash}\\
\lean{TS.ArrangementSym.}\allowbreak\ltwentynineln{jetPoincareCoreT0_of_symArrangement} & Poincar\'e core by duality & $\mathrm{SAI}_n$ & Thm.~\ref{l29:thm:poincarecore}\\
\lean{TS.GenIncl.genInclusion}, \lean{symArrangementInclusion} & block generators span $L_N\cap(\text{vanishing ideal})$ & $2\le n$ (resp.\ $1\le n$) & Thm.~\ref{l29:thm:geninclusion}\\
\lean{TS.JetPoincareTwo.}\allowbreak\ltwentynineln{twoSeries_gradedUnitary_M27_of_affine} & $(A_1,D_5)\to(A_1,A_4)$ & \lean{hA} & Cor.~\ref{l29:cor:m27}\\
\ltwentynineln{TS.leeYang_gradedUnitary_of_affine_viaTwoSeries} & Lee--Yang recovered & \lean{hA} & Thm.~\ref{l29:thm:leeyang}\\
\bottomrule
\end{longtable}
\end{center}
(\lean{TS} abbreviates \lean{VOA.TwoSeries}.)

%======================================================================
\section{Notes and references}\label{l29:sec:notes}
%======================================================================

\paragraph{Graded unitarity and the four-dimensional origin.} A.~Arabi Ardehali, C.~Beem, M.~Lemos and L.~Rastelli,
\emph{Graded unitarity in the SCFT/VOA correspondence}, arXiv:2507.23781 (2025). The SCFT/VOA correspondence: C.~Beem,
M.~Lemos, P.~Liendo, W.~Peelaers, L.~Rastelli and B.~C.~van Rees, \emph{Infinite chiral symmetry in four dimensions},
Comm.\ Math.\ Phys.\ 336 (2015) 1359--1433, arXiv:1312.5344. The identification of the chiral algebras of the
Argyres--Douglas theories $(A_1,A_{2n})$ with $M(2,2n+3)$ and of $(A_1,D_{2n+1})$ with $L_{-4n/(2n+1)}(\ltwentyninesl)$ is
due to C.~Cordova and S.-H.~Shao, \emph{Schur indices, BPS particles, and Argyres--Douglas theories}, JHEP 01 (2016) 040,
arXiv:1506.00265; see also T.~Creutzig, \emph{W-algebras for
Argyres--Douglas theories}, Eur.\ J.\ Math.\ 3 (2017) 659--690, arXiv:1701.05926. The BRST coset construction and the
Hodge package: C.~Beem and S.~Garner, arXiv:2509.10364 (2025), as cited in the source files. The conjecture that the
boundary minimal models are graded unitary goes back to R.~M.~Kaufmann, hep-th/9405041 (1994) (see Chapter~\ref{ch:L19}).

\paragraph{Cosets and admissible levels.} P.~Goddard, A.~Kent and D.~Olive, \emph{Unitary representations of the
Virasoro and super-Virasoro algebras}, Comm.\ Math.\ Phys.\ 103 (1986) 105--119 (the coset construction); V.~G.~Kac and
M.~Wakimoto, \emph{Modular invariant representations of infinite-dimensional Lie algebras and superalgebras}, Proc.\ Natl.\
Acad.\ Sci.\ USA 85 (1988) 4956--4960 (admissible levels). T.~Arakawa and C.~Dai, on classical freeness and characters at
boundary admissible levels, as in Chapter~\ref{ch:L20}.

\paragraph{Drinfeld--Sokolov reduction.} B.~Feigin and E.~Frenkel, \emph{Quantization of the Drinfeld--Sokolov
reduction}, Phys.\ Lett.\ B 246 (1990) 75--81; T.~Arakawa, \emph{Representation theory of $W$-algebras}, Invent.\ Math.\
169 (2007) 219--320, where (in the ``$-$'' reduction) the reduction of an irreducible highest-weight module is shown to be
irreducible or zero; for admissible $k$ this gives $H_{DS}(L_k(\ltwentyninesl))\cong L(c_{\mathrm{DS}}(k),0)$. No result of
Arakawa's is assumed by the Lean development; only the ``descends'' half is proved, for the explicit slice map. The Whittaker coinvariant description is in the spirit of
B.~Kostant, \emph{On Whittaker vectors and representation theory}, Invent.\ Math.\ 48 (1978) 101--184. The slice map of this
chapter is the project's own construction.

\paragraph{Characters.} G.~E.~Andrews, \emph{An analytic generalization of the Rogers--Ramanujan identities for odd
moduli}, Proc.\ Natl.\ Acad.\ Sci.\ USA 71 (1974) 4082--4085, and \emph{The Theory of Partitions}, Addison-Wesley 1976,
Chapter~7 (Andrews' recursion); B.~Gordon, \emph{A combinatorial generalization of the Rogers--Ramanujan identities},
Amer.\ J.\ Math.\ 83 (1961) 393--399. A.~Rocha-Caridi, \emph{Vacuum vector representations of the Virasoro algebra}, in
\emph{Vertex Operators in Mathematics and Physics}, MSRI Publ.\ 3 (1985) 451--473. A.~Benoit and Y.~Saint-Aubin, \emph{Degenerate
conformal field theories and explicit expressions for some null vectors}, Phys.\ Lett.\ B 215 (1988) 517--522. B.~Feigin and
A.~Stoyanovsky, \emph{Quasi-particles models for the representations of Lie algebras and geometry of flag manifold},
arXiv:hep-th/9308079 (functional realisations). The fermionic formula for the boundary levels: S.~Chern, D.~Tran and A.~Wakhare,
as cited in the source files (we do not reproduce the title). The Bailey chain: G.~E.~Andrews, \emph{Multiple series
Rogers--Ramanujan type identities}, Pacific J.\ Math.\ 114 (1984) 267--283. P.~A.~MacMahon, \emph{Combinatory Analysis}
(1915--16), for the inversion generating function of words.

\paragraph{The $n$-equal arrangement.} S.-Y.~R.~Li and W.~C.~W.~Li, \emph{Independence numbers of graphs and generators of
ideals}, Combinatorica 1 (1981) 55--61 (generators of the ideal of the $n$-equal arrangement). The point that the jet-ideal
Poincar\'e lemma reduces to this arrangement, via a fermionic realisation, is suggested by Ait El Manssour--Pogudin's work on
arcs of fat points.

\endgroup
